\chapter{Filtrations, the \texorpdfstring{$C_2$}{C2} algebra and presentations}\label{ch:L11}
\begingroup\setlength{\emergencystretch}{3em}

\providecommand{\lelevenvac}{\mathbf{1}}
\providecommand{\lelevenR}{R_V}
\providecommand{\lelevengr}{\operatorname{gr}}
\providecommand{\lelevenTR}{T(R)}
\providecommand{\lelevenF}{F}
\providecommand{\lelevenJ}{J}

This chapter documents the release of the LeanVOA library named in the front matter. Every Lean quotation below was copied by a script that reads the
source lines (so that no declaration was retyped by hand) and is machine-checked by
\texttt{scripts/textbook/sync.py} to occur verbatim, up to whitespace and comments, in the source
files of that release. Checks that were run at an earlier state of the library and not re-run for this edition are marked as such. The chapter covers the 38 files assigned to it in
\texttt{Notes/Textbook/assignment.txt}: the filtration, symbol and presentation files of
\leanpath{VOA/\allowbreak Core/} (\leanpath{GoodFiltration*}, \leanpath{LinearFiltration},
\leanpath{FilteredTensorSymbols}, \leanpath{FilteredWordSymbols}, \leanpath{FiniteAssociatedGraded},
\leanpath{FiniteHilbertSqueeze}, \leanpath{Polynomial*}, \leanpath{OrderedExponentWord},
\leanpath{IntegerJetConvolution}, \leanpath{PowerSeriesCoefficientDerivation},
\leanpath{SparseRational}, \leanpath{NonlinearConformal}, \leanpath{NonlinearVirasoro},
\leanpath{NonlinearZeta}, \leanpath{NPExtension}, \leanpath{NonlinearMR}, \leanpath{VAPresentation},
\leanpath{PresentationDoctrine}), the two directories \leanpath{VOA/\allowbreak Commutative/} and
\leanpath{VOA/\allowbreak Filtered/}, and the example files \leanpath{VOA/\allowbreak Examples/\allowbreak LiDelta.lean},
\leanpath{LiDelta2.lean}, \leanpath{NonlinearVirasoroTransport.lean} and
\leanpath{NonlinearVirasoroGenuine.lean}. All together these are about fourteen thousand lines of
Lean. Section~\ref{l11:sec:genwords} documents three further files of
\leanpath{VOA/\allowbreak Filtered/} (\leanpath{GeneratorWord}\allowbreak\leanpath{LiSurjectivity.lean},
\leanpath{GeneratorCurrent}\allowbreak\leanpath{LiSpanning.lean}, \leanpath{AffineSl2Generator}\allowbreak\leanpath{CurrentInstance.lean}), first quoted at an
earlier commit on the branch \texttt{deligne-lean}.

\paragraph{What the chapter is about.} A vertex operator algebra $V$ is a very non-commutative
object: its products $u_{(n)}v$, $n\in\Z$, satisfy the Borcherds identity, a tower of
relations with no classical analogue. There are two standard ways to tame it, and this chapter
documents the library's version of both.
\begin{enumerate}
\item \emph{Pass to a commutative shadow.} Filter $V$ in a way compatible with the products and
  take the associated graded object. For Li's filtration the associated graded
  $\lelevengr^F V$ is a commutative algebra with a derivation (indeed a vertex Poisson
  algebra), and its degree-zero part is Zhu's $C_2$-algebra $\lelevenR = V/C_2(V)$. Arakawa's
  \emph{associated variety} of $V$ is the spectrum of $\lelevenR$, and a basic theorem of Li and
  Arakawa says that the arc (jet) algebra of $\lelevenR$ surjects onto $\lelevengr^F V$.
  ``Classical freeness'' is the statement that this surjection is an isomorphism.
\item \emph{Present the algebra.} Write $V$ as generated by finitely many fields with a given
  singular OPE, modulo explicit relations. This is how physicists specify a chiral algebra, and De
  Sole and Kac turned it into mathematics: the universal enveloping vertex algebra of a
  non-linear Lie conformal algebra. (The library's datum \lean{OPEData} is the weaker non-linear \emph{conformal} algebra of
  De Sole--Kac's Definition~3.1: skew-symmetry and the Jacobi identity are not imposed; see Definition~\ref{l11:def:opedata}.)
\end{enumerate}
The library has substantial, carefully staged material on both. Sections~\ref{l11:sec:flags}
to~\ref{l11:sec:symbols} treat filtrations, the Li filtration, jets and the polynomial ``symbol''
toolkit. Within them, Section~\ref{l11:sec:genwords} proves that the canonical jet map onto the Li associated graded is
surjective whenever generator words span Li's filtration, for any VOA. It also shows how to establish that spanning for
currents and for quotients; Chapters~\ref{ch:L20} and~\ref{ch:L24} use this; Sections~\ref{l11:sec:dsk} to~\ref{l11:sec:doctrine} treat presentations; a short
Section~\ref{l11:sec:lidelta} documents two files on Li's $\Delta$-operator that were filed here.

\paragraph{The findings of this chapter, in brief.} Most of what is documented here is proved
without hypotheses, \texttt{\#print axioms}-clean, and without \lean{VOA.pbw} (the PBW theorem of Chapter~\ref{ch:L02}; formerly a project axiom, now a theorem). The fidelity checks, however, turned up the following points, each explained in a
red warning box where it occurs.
\begin{itemize}
\item \textbf{Every VOA has a \lean{GoodFiltration}.} The library's \lean{GoodFiltration} (a
  reconstruction, flagged as such in its own docstring, of the ``good filtration'' clause of
  Arabi~Ardehali--Beem--Lemos--Rastelli) is satisfied by the coarse two-step flag
  ``vacuum in degree $0$, everything else in degree $1$''. We prove this in Lean (checked against the library; the check is not part of this
  release). So the existence of a good filtration carries no information; only
  a \emph{specific} filtration does. The mode condition also omits the second clause of Li's and
  Arakawa's notion of a good filtration (non-negative modes lower the degree by one).
\item \textbf{Li's $F^1V = C_2(V)$ is proved in the library} (\lean{VOA.}\allowbreak\lean{liFiltration_}\allowbreak\lean{one_}\allowbreak\lean{eq_}\allowbreak\lean{O₂},
  \leanpath{VOA/\allowbreak Filtered/\allowbreak LiC2Bridge.lean}); it was first checked in Lean for this chapter, outside the library, at a
  commit where the library stated only $C_2(V)\subseteq F^1V$. Hence
  the degree-zero layer of the library's Li associated graded \emph{is} Zhu's $C_2$-algebra. The
  same check also shows that an explicit \lean{LWGenerated} hypothesis carried by two Li-filtration
  theorems is redundant.
\item \textbf{The De Sole--Kac interface \lean{NPExtension} has no instance anywhere in the
  library.} Every theorem quantified over an \lean{NP : NPExtension D} (the full ideal $M(R)$, the
  presentation ideal of \lean{VAPresentation}, the vacuum guards, the ``$\omega$-counterexample'')
  is therefore, as a Lean proposition, vacuously true today. The library says so itself; we repeat
  it prominently. Likewise \lean{simple_}\allowbreak\lean{realization_}\allowbreak\lean{unique} is a \lean{def : Prop}, stated and not
  proved, and \lean{TheVOA} takes an \emph{arbitrary} VOA structure on the presented space as input,
  unrelated to the presentation.
\item \textbf{The ``generators $+$ OPE $=$ Borcherds'' theorem for Virasoro is genuine but
  narrower than its name.} The ideal \lean{MRtwo} is defined from the OPE table alone and the
  library proves \lean{MRtwo_eq_ker_PhiVir} honestly. But the relations do not depend on the
  central charge (we check this in Lean, outside the library), and the
  annihilation half of the fields is obtained by transporting the independently constructed
  Virasoro VOA, not from the OPE. So the theorem identifies the \emph{vector space and creation
  operators} of the presented object; the OPE data enter only through the realization.
\item \textbf{The \lean{hdec} surjectivity criterion of Theorem~\ref{l11:thm:hits} cannot be used outside the
  $C_2$-cofinite world.} Its hypothesis implies \lean{C2Cofinite}; we check this in Lean, outside the library. Hence
  it fails for every non-lisse VOA, e.g.\ $L_{-4/3}(\mathfrak{sl}_2)$, and also for the Virasoro minimal models
  $M(2,2n+1)$, $n\ge3$ (the $M(2,2m+3)$ of Chapter~\ref{ch:L20}, with $n=m+1$), with the single generator $\omega$. The
  replacement, spanning of Li's filtration by generator words,
  is Section~\ref{l11:sec:genwords}. This is a recorded dead end (warning box on page~\pageref{l11:warn:hdec}).
\item The left-ideal route \lean{MRbot} of the original De Sole--Kac staging is refuted in Lean
  (\lean{not_}\allowbreak\lean{nonempty_}\allowbreak\lean{genFields_}\allowbreak\lean{MRbot_}\allowbreak\lean{virasoro}): this is a recorded dead end.
\end{itemize}

%======================================================================
\section{Conventions recalled}\label{l11:sec:conv}
%======================================================================

Throughout, $V$ is a vertex operator algebra in the sense of the library: a Lean term
\lean{W : VertexOperatorAlgebra 𝒱} (Chapter~\ref{ch:L03}, Definition~\ref{l03:def:VOA}) with
underlying complex vector space \lean{𝒱}. We use:
\begin{itemize}
\item $V = \bigoplus_{n\in\Z}V_n$ with $V_n = $ \lean{W.grading n}, each finite-dimensional,
  $V_n = 0$ for $n<0$ (\lean{positivity}) and $V_0 = \C\lelevenvac$ (\lean{vacuum_grade}). Every
  library VOA is of CFT type.
\item The $n$-th product $u_{(n)}v$ is \lean{W.smul u n v}, and \lean{mode f n} is the coefficient
  $f_{-n-1}$ of a formal series $f=\sum_kf_kz^k$, so that \lean{mode (Y u) n} $=u_{(n)}$
  (Chapter~\ref{ch:01}, Example~\ref{ch01:ex:mode}). Conversely the \emph{$z^j$-coefficient}
  \lean{(Y u) j} of $Y(u,z)$ is $u_{(-j-1)}$. Both indexings occur in this chapter and we always say
  which one is meant.
\item $T = L_{-1}$ is the translation operator (\lean{W.L (-1)}, or
  \lean{W.toPreVertexAlgebra.T}); $(Tu)_{(n)} = -n\,u_{(n-1)}$ and $u_{(-n-1)}\lelevenvac =
  T^nu/n!$.
\item $C_2(V) = \operatorname{span}\{u_{(-2)}v\}$ is \lean{O₂ W}, and $\lelevenR = V/C_2(V)$ is
  \lean{A₂ W} (Chapter~\ref{ch:L12}, Definition~\ref{l12:def:c2}; recall that \lean{A₂} is
  \emph{not} the second higher Zhu algebra). $\lelevenR$ is a commutative Poisson algebra with
  product $[u][v] = [u_{(-1)}v]$ and bracket $\{[u],[v]\} = [u_{(0)}v]$
  (Theorem~\ref{l12:thm:poisson}).
\item \lean{Submodule ℂ X} is a linear subspace; \lean{⊥}, \lean{⊤} are $0$ and $X$; \lean{⊔},
  \lean{⊓} are sum and intersection; \lean{⨆ p, S p} is the sum of a family; \lean{S.map f},
  \lean{S.comap f} are image and preimage under a linear map; \lean{X ⧸ S} is the quotient;
  \lean{Module.finrank} is dimension (and is $0$ on infinite-dimensional spaces). \lean{MvPolynomial
  σ ℂ} is the polynomial ring $\C[x_s : s\in\sigma]$ and \lean{PowerSeries R} is $R[[t]]$.
  Generic Lean idioms (structures, \lean{∃}, coercions, \lean{↥}, junk values) are explained in
  Chapter~\ref{ch:01}, Section~\ref{ch01:sec:lean}.
\item \emph{Normal ordering.} For fields $a(z) = \sum_ja_jz^j$ (coefficient convention), ${:}a(z)b(z){:} = a(z)_+b(z)+b(z)a(z)_-$,
  where $a(z)_+$ collects the creation part (the modes $a_{(n)}$, $n<0$, which need not kill the vacuum) and $a(z)_-$ the
  annihilation part ($a_{(n)}$, $n\ge0$): annihilation operators are moved to the right. Its value on states is the $(-1)$-product,
  $Y(a_{(-1)}b,z) = {:}Y(a,z)Y(b,z){:}$; the library's \lean{nprod} is defined in Chapter~\ref{ch:L03}.
\item \emph{Strong generators.} A set $S\subseteq V$ \emph{strongly generates} $V$ if $V$ is spanned by the states
  $s^1_{(-n_1-1)}\cdots s^r_{(-n_r-1)}\lelevenvac$ with $s^i\in S$, $n_i\ge0$; equivalently every field is a normal-ordered polynomial
  in the fields of $S$ and their derivatives. The library's version is \lean{iterProdSpan} (Chapter~\ref{ch:L12},
  Theorem~\ref{l12:thm:fingen}).
\item Fields, locality and $T$-covariance: \lean{IsField a} says $a(z)x$ has only finitely many
  negative powers of $z$ for every $x$; \lean{Local a b} says $(z-w)^N[a(z),b(w)] = 0$ for some
  $N$; \lean{DCovariant T a} says $[T,a(z)] = \partial_za(z)$. All three are defined in
  Chapter~\ref{ch:L03} (Section~\ref{l03:sec:conventions} and
  Definition~\ref{l03:def:dcov}), as is the operator-first reconstruction theorem
  \lean{Reconstruction.}\allowbreak\lean{reconstruct} (Theorem~\ref{l03:thm:reconstruct}).
\end{itemize}

\begin{physics}
The whole chapter is about the relation between a \emph{chiral algebra} and its \emph{classical
limit}. Think of $V$ as the operator algebra of the holomorphic fields of a two-dimensional
conformal field theory. A filtration by ``how many derivatives'' (Li's filtration) or by ``how many
fields'' (a length or $R$-charge filtration) organises the operators the way a physicist organises
a semiclassical expansion; the associated graded is the algebra of \emph{classical} fields, in which
normal ordering becomes ordinary commutative multiplication and the singular part of the OPE
survives only as a Poisson bracket. Its degree-zero part $\lelevenR = V/C_2(V)$ is the ring of
functions on the \emph{associated variety} $X_V = \operatorname{Specm}\lelevenR$. For the chiral
algebras of four-dimensional $\mathcal N=2$ superconformal theories, Beem and Rastelli conjectured
(and checked in many examples) that $X_V$ is the \emph{Higgs branch} of the four-dimensional theory:
the classical moduli space of vacua. Strong generators of $V$ are a minimal set of fields from which
all others are obtained by normal ordering and derivatives (they generate $\lelevenR$ as an algebra),
and the \emph{null relations} among them --- the singular vectors that vanish in the simple quotient
--- are the chiral-algebra version of equations of motion, or of the defining equations of the Higgs
branch.
\end{physics}

%======================================================================
\section{Graded flags, good filtrations and filtered maps}\label{l11:sec:flags}
%======================================================================

We start with the most general notion: a filtration of each weight space. The library introduced it
for the graded-unitarity programme of Arabi~Ardehali, Beem, Lemos and Rastelli (ABLR), which is
treated in Chapter~\ref{ch:L19}; here we only document the definitions and the facts about them.

\begin{definition}[Graded flag]\label{l11:def:flag}
A \emph{graded flag} on a VOA $V$ is a family of subspaces $G_pV_n\subseteq V_n$ ($p\in\N$,
$n\in\Z$) such that
\begin{enumerate}
\item $G_pV_n\subseteq G_qV_n$ whenever $p\le q$ (increasing);
\item $G_0V_n = 0$ for $n>0$, and $G_0V_0 = \C\lelevenvac$ (normalisation);
\item $\sum_pG_pV_n = V_n$ for every $n$ (exhaustive).
\end{enumerate}
It is \emph{mode-compatible} if for all $p,q\in\N$, $m,n,k\in\Z$, $u\in G_pV_m$ and $v\in G_qV_n$,
\begin{equation}\label{l11:eq:modecompat}
u_{(k)}v\in G_{p+q}V_{m+n-k-1}.
\end{equation}
A \emph{good filtration} in the library's sense is a mode-compatible graded flag.
\end{definition}
\noindent File \leanpath{VOA/\allowbreak Core/\allowbreak GoodFiltration.lean}.
\begin{leancode}
structure GradedFlag (W : VertexOperatorAlgebra 𝒱) where
  piece : ℕ → ℤ → Submodule ℂ 𝒱
  le_grading : ∀ (p : ℕ) (n : ℤ), piece p n ≤ W.grading n
  mono : ∀ {p q : ℕ} (n : ℤ), p ≤ q → piece p n ≤ piece q n
  bot : ∀ n : ℤ, 0 < n → piece 0 n = ⊥
  vacuum_zero : piece 0 0 = ℂ ∙ W.vacuum
  exhaustive : ∀ n : ℤ, (⨆ p : ℕ, piece p n) = W.grading n

def ModeCompatible : Prop :=
  ∀ {p q : ℕ} {m n k : ℤ} {u v : 𝒱}, u ∈ F.piece p m → v ∈ F.piece q n →
    W.smul u k v ∈ F.piece (p + q) (m + n - k - 1)

structure GoodFiltration (W : VertexOperatorAlgebra 𝒱) extends GradedFlag W where
  mode_mem : toGradedFlag.ModeCompatible
\end{leancode}

\begin{dictionary}
\lean{piece p n} is $G_pV_n$, a subspace of the whole space \lean{𝒱} together with the field
\lean{le_grading} saying it lies in $V_n$. \lean{mono} is (1); \lean{bot} and \lean{vacuum_zero} are
(2), with \lean{ℂ ∙ W.vacuum} the line spanned by the vacuum; \lean{exhaustive} is (3), \lean{⨆} being
the sum of subspaces. Note that in negative weight nothing has to be imposed, because $V_n = 0$ there.
\lean{ModeCompatible} is (\ref{l11:eq:modecompat}) with all six variables universally quantified
(the curly braces make them implicit, which changes nothing logically); \lean{W.smul u k v} is
$u_{(k)}v$ and the target weight \lean{m + n - k - 1} is the standard weight of $u_{(k)}v$.
\lean{GoodFiltration W} \lean{extends} \lean{GradedFlag W}: it is a graded flag plus the single field
\lean{mode_mem}.

\textbf{Comparison with the literature.} Li (2005) and Arakawa (2012) call an increasing filtration
$\{G_pV\}$ \emph{good} if $\lelevenvac\in G_0V$ and, for $u\in G_pV$, $v\in G_qV$,
\begin{equation*}
u_{(k)}v\in G_{p+q}V\ (k\in\Z),\qquad u_{(k)}v\in G_{p+q-1}V\ (k\ge0).
\end{equation*}
The second clause is what makes the associated graded a \emph{vertex Poisson} algebra (the non-negative
modes become a bracket of degree $-1$). The library's \lean{ModeCompatible} is only the first clause.
It is therefore \emph{weaker} than the Li--Arakawa notion. ABLR's own definition is not in the
repository, and the library's docstring says explicitly that \lean{ModeCompatible} ``IS A
RECONSTRUCTION, NOT A TRANSCRIPTION'' of ABLR's clause (iv). The graded-unitarity theorems of
Chapter~\ref{ch:L19} use only the underlying \lean{GradedFlag}, so this weakness does not affect them.
\end{dictionary}

\begin{warning}
\textbf{Every VOA has a \lean{GoodFiltration} (moderate).} Put $G_0V_n = 0$ for $n>0$ and $G_pV_n = V_n$
in all other cases. This is a graded flag (with $G_0V_0 = V_0 = \C\lelevenvac$ by CFT type), and it is
mode-compatible: if $p+q\ge1$ the target $G_{p+q}V_{m+n-k-1}$ is the whole weight space, which
contains $u_{(k)}v$; if $p=q=0$ then $u,v\in\C\lelevenvac$ (or are $0$), and
$\lelevenvac_{(k)}\lelevenvac = \delta_{k,-1}\lelevenvac$ lies in $G_0V_0$. A Lean check made for this chapter (not part of this release) constructs this
\lean{GoodFiltration W} for an arbitrary \lean{W}; \texttt{\#print axioms}: standard three. (The coarse
flag even satisfies the stronger Li--Arakawa condition.) Consequently, a hypothesis
``\lean{F : GoodFiltration W}'' never restricts \lean{W}; whatever content a theorem has must come from
further conditions on the specific \lean{F} (in Chapter~\ref{ch:L19}: nondegeneracy of a form on each
piece). The structure is not vacuous --- it has instances --- but it is not informative by itself: only specific instances (the
length filtration of the Virasoro minimal models, the $R$-filtration of chiral algebras of 4d $\mathcal N=2$ theories, the affine
filtration of \leanpath{AffineSl2GoodFiltration.lean}) carry content, and the theorems of Chapter~\ref{ch:L19} add a
nondegeneracy condition on each piece.
\end{warning}

\textbf{Why the definition is still reasonable.} In physics, the filtration that matters is a specific
one: for the boundary Virasoro minimal models ABLR take $G_pV$ to be the span of
$L_{-n_1}\cdots L_{-n_k}\lelevenvac$ with $k\le p$ (the \emph{length} of the partition), and for chiral
algebras of $\mathcal N=2$ theories the $R$-filtration by $SU(2)_R$ charge. The library constructs such
filtrations (\leanpath{VOA/\allowbreak Examples/\allowbreak VirasoroGoodFiltration.lean},
\leanpath{AffineSl2GoodFiltration.lean}, Chapters~\ref{ch:L07} and~\ref{ch:L19}); the structure here is
only the container.

\begin{physics}
A graded flag is a way of counting ``how composite'' a state is. For the Virasoro vacuum module the
length filtration counts the number of stress-tensor modes needed to create a state; for a chiral
algebra of a four-dimensional $\mathcal N=2$ theory, the Schur operators carry an $SU(2)_R$ charge $R$,
and the $R$-filtration is the filtration by $R$. Mode compatibility says that the OPE of two operators
of ``lengths'' $p$ and $q$ contains only operators of length at most $p+q$: the OPE cannot create more
constituents than it was given. Graded unitarity (Chapter~\ref{ch:L19}) is the statement that a natural
Hermitian form is positive on the associated graded --- the chiral-algebra shadow of four-dimensional
unitarity.
\end{physics}

\begin{theorem}[Flags and good filtrations descend to quotients]\label{l11:thm:flagquot}
Let $I\subseteq V$ be a graded ideal and $\pi:V\to V/I$ the quotient VOA map. If $\{G_pV_n\}$ is a graded
flag on $V$, then $\{\pi(G_pV_n)\}$ is a graded flag on $V/I$; if the former is mode-compatible, so is
the latter. In particular every good filtration of $V$ induces one of $V/I$.
\end{theorem}
\noindent File \leanpath{VOA/\allowbreak Core/\allowbreak GoodFiltrationQuotient.lean}.
\begin{leancode}
noncomputable def quotient {W : VertexOperatorAlgebra V} (F : GradedFlag W)
    (I : Submodule ℂ V) (hI : W.IsIdeal I) (hgr : QuotVOA.IsGradedIdeal W I) :
    GradedFlag (QuotVOA.quotientVOA W I hI hgr) where
  piece p n := (F.piece p n).map I.mkQ

theorem quotient_modeCompatible {W : VertexOperatorAlgebra V} (F : GradedFlag W)
    (I : Submodule ℂ V) (hI : W.IsIdeal I) (hgr : QuotVOA.IsGradedIdeal W I)
    (hmode : F.ModeCompatible) : (F.quotient I hI hgr).ModeCompatible := by

noncomputable def quotient {W : VertexOperatorAlgebra V} (F : GoodFiltration W)
    (I : Submodule ℂ V) (hI : W.IsIdeal I) (hgr : QuotVOA.IsGradedIdeal W I) :
    GoodFiltration (QuotVOA.quotientVOA W I hI hgr) where
  toGradedFlag := F.toGradedFlag.quotient I hI hgr
\end{leancode}

\begin{dictionary}
\lean{W.IsIdeal I} and \lean{QuotVOA.}\allowbreak\lean{IsGradedIdeal W I} say that $I$ is an ideal of the VOA and is
graded; \lean{QuotVOA.}\allowbreak\lean{quotientVOA W I hI hgr} is the quotient VOA $V/I$ (Chapter~\ref{ch:L04},
Definition~\ref{l04:def:quot}). \lean{I.mkQ} is the quotient map $\pi$, and \lean{(F.}\allowbreak\lean{piece p n).}\allowbreak\lean{map
I.}\allowbreak\lean{mkQ} is $\pi(G_pV_n)$. The two \lean{quotient} constructions are definitions (they produce the
induced flag) and \lean{quotient_}\allowbreak\lean{modeCompatible} is the theorem; the induced filtration is \emph{the}
image filtration, as in the literature. Nothing is claimed about whether the induced filtration agrees
with some independently defined filtration of $V/I$ (the docstring says so).
\end{dictionary}

\textbf{Proof on paper.} Images of an increasing exhaustive family are increasing and exhaustive, and
$\pi(\C\lelevenvac) = \C\pi(\lelevenvac)$. For mode compatibility, write $\bar u = \pi(u)$,
$\bar v=\pi(v)$ with $u\in G_pV_m$, $v\in G_qV_n$; since $\pi$ is a VOA homomorphism,
$\bar u_{(k)}\bar v = \pi(u_{(k)}v)\in\pi(G_{p+q}V_{m+n-k-1})$. \qed

\begin{status}
Proved, unconditional; \texttt{\#print axioms} of \lean{VOA.}\allowbreak\lean{GoodFiltration.}\allowbreak\lean{quotient} and
\lean{VOA.}\allowbreak\lean{GradedFlag.}\allowbreak\lean{quotient_}\allowbreak\lean{modeCompatible}: \lean{[propext, Classical.}\allowbreak\lean{choice, Quot.}\allowbreak\lean{sound]}.
\end{status}

\begin{physics}
Passing to a quotient VOA is setting null fields to zero (going from the universal to the simple
chiral algebra, say). The statement says that a ``number of constituents'' count on the universal
algebra still makes sense after imposing the null relations, which is how one defines, for example,
the $R$-filtration on a simple quotient once it is known on a free-field or universal model.
\end{physics}

\subsection{Filtrations indexed by a preorder, and strict maps}

For the inverse Drinfeld--Sokolov programme (Chapter~\ref{ch:L19}) the library needs filtrations
indexed by $\Z$ rather than $\N$, and maps between filtered spaces.

\begin{definition}[Linear filtration; filtered and strictly filtered maps]\label{l11:def:linfilt}
Let $\iota$ be a preordered set. A \emph{linear filtration} of an $R$-module $M$ indexed by $\iota$ is
a monotone family $(F_i)_{i\in\iota}$ of submodules. Given filtrations $F$ on $M$ and $G$ on $N$
(indexed by $\iota$ and $\kappa$), a map $\varphi:\iota\to\kappa$ and a linear map $f:M\to N$, $f$ is
\emph{filtered along $\varphi$} if $f(F_i)\subseteq G_{\varphi(i)}$ for all $i$, and \emph{strictly
filtered along $\varphi$} if $f(F_i) = f(M)\cap G_{\varphi(i)}$ for all $i$.
\end{definition}
\noindent File \leanpath{VOA/\allowbreak Core/\allowbreak LinearFiltration.lean}.
\begin{leancode}
structure LinearFiltration (ι : Type uι) (R : Type uR) (M : Type uM)
    [Preorder ι] [Semiring R] [AddCommMonoid M] [Module R M] where
  piece : ι → Submodule R M
  mono' : Monotone piece

def IsFilteredAlong (F : LinearFiltration ι R M) (G : LinearFiltration κ R N)
    (phi : ι → κ) (f : M →ₗ[R] N) : Prop :=
  ∀ i, F.piece i ≤ (G.piece (phi i)).comap f

def IsStrictFilteredAlong (F : LinearFiltration ι R M) (G : LinearFiltration κ R N)
    (phi : ι → κ) (f : M →ₗ[R] N) : Prop :=
  ∀ i, (F.piece i).map f = LinearMap.range f ⊓ G.piece (phi i)
\end{leancode}

\begin{dictionary}
\lean{Monotone piece} is $i\le j\Rightarrow F_i\subseteq F_j$. \lean{(G.}\allowbreak\lean{piece (phi i)).}\allowbreak\lean{comap f} is the
preimage $f^{-1}(G_{\varphi(i)})$, so \lean{IsFilteredAlong} is $F_i\subseteq f^{-1}(G_{\varphi(i)})$,
i.e.\ $f(F_i)\subseteq G_{\varphi(i)}$. \lean{LinearMap.}\allowbreak\lean{range f ⊓ G.}\allowbreak\lean{piece (phi i)} is $f(M)\cap
G_{\varphi(i)}$. These are the textbook notions of a filtered and a strict morphism of filtered modules
(Bourbaki, \emph{Alg\`ebre commutative} III; the strictness notion of Deligne's mixed Hodge theory). No
monotonicity of $\varphi$ is required to \emph{state} them. Strictness is genuinely stronger: the
inclusion of $\C$ into $\C$ with $F_0 = 0$, $F_1=\C$ on the source and $G_0 = G_1 = \C$ on the target is
filtered but not strict.
\end{dictionary}

\begin{physics}
A filtration indexed by $\Z$ is natural for a free-field (Wakimoto-type) algebra, where a ``charge'' or ``number of $\beta\gamma$
pairs minus number of bosons'' can be negative; a filtered map is a realisation that respects the leading-order counting, and a
strict one is a realisation in which leading terms never cancel accidentally.
\end{physics}


\begin{theorem}[Strictness criteria]\label{l11:thm:strict}
\begin{enumerate}
\item If $f$ is injective, $f$ is strictly filtered along $\varphi$ iff $F_i = f^{-1}(G_{\varphi(i)})$
  for all $i$.
\item (Dimension criterion.) Over a division ring $K$: if $f$ is injective and filtered along
  $\varphi$, and for every $i$ the space $f(M)\cap G_{\varphi(i)}$ is finite-dimensional of the same
  dimension as $F_i$, then $f$ is strictly filtered along $\varphi$.
\end{enumerate}
\end{theorem}
\noindent File \leanpath{VOA/\allowbreak Core/\allowbreak LinearFiltration.lean}.
\begin{leancode}
theorem strictFilteredAlong_iff_comap_eq
    {F : LinearFiltration ι R M} {G : LinearFiltration κ R N}
    {phi : ι → κ} {f : M →ₗ[R] N} (hf : Function.Injective f) :
    F.IsStrictFilteredAlong G phi f ↔
      ∀ i, F.piece i = (G.piece (phi i)).comap f := by

theorem IsFilteredAlong.isStrictFilteredAlong_of_finrank_eq_range_inf
    {F : LinearFiltration ι K V} {G : LinearFiltration κ K W}
    {phi : ι → κ} {f : V →ₗ[K] W}
    (hfiltered : F.IsFilteredAlong G phi f) (hf : Function.Injective f)
    (hfinite : ∀ i, FiniteDimensional K ↥(LinearMap.range f ⊓ G.piece (phi i)))
    (hfinrank : ∀ i, Module.finrank K (F.piece i) =
      Module.finrank K ↥(LinearMap.range f ⊓ G.piece (phi i))) :
    F.IsStrictFilteredAlong G phi f := by
\end{leancode}

\begin{dictionary}
\lean{Function.Injective f} is injectivity; \lean{↥(LinearMap.}\allowbreak\lean{range f ⊓ G.}\allowbreak\lean{piece (phi i))} is that
intersection regarded as a vector space in its own right (the up-arrow \lean{↥} turns a subspace into a
type). Both statements are exactly as written above; the second only needs finite-dimensionality of the
\emph{individual} pieces, which is the situation of a graded VOA (infinite-dimensional in total, finite
in each weight).
\end{dictionary}

\textbf{Proof on paper.} (1) If $f$ is injective, $f(F_i) = f(M)\cap G_{\varphi(i)}$ pulls back to
$F_i = f^{-1}(G_{\varphi(i)})$ and conversely. (2) $f(F_i)\subseteq f(M)\cap G_{\varphi(i)}$ by
filteredness, and $\dim f(F_i) = \dim F_i$ by injectivity; a subspace of a finite-dimensional space of
the same dimension is the whole space. \qed

\begin{status}
Proved, unconditional; \texttt{\#print axioms} of both: standard three.
\end{status}

\begin{physics}
Strictness is the precise form of ``no hidden cancellations'': when a VOA is realised inside a
free-field algebra (Wakimoto, Adamovi\'c's inverse Drinfeld--Sokolov), each operator has a leading
free-field term, and strictness says that a combination whose leading terms cancel really is of lower
order in the source as well. The dimension criterion turns this into a \emph{character} computation:
compare filtered characters on both sides.
\end{physics}

\begin{definition}[Exact layer]\label{l11:def:layer}
For subspaces $Q\subseteq P$ of a vector space, the \emph{exact layer} is $P/Q$, realised as the quotient
of $P$ by the preimage of $Q$ under the inclusion $P\hookrightarrow V$.
\end{definition}
\noindent File \leanpath{VOA/\allowbreak Core/\allowbreak FiniteAssociatedGraded.lean}.
\begin{leancode}
abbrev Layer (P Q : Submodule K V) (_h : Q ≤ P) : Type _ :=
  P ⧸ Q.comap P.subtype

theorem finrank_layer_add_finrank
    (P Q : Submodule K V) (h : Q ≤ P) [FiniteDimensional K P] :
  finrank K (Layer P Q h) + finrank K Q = finrank K P := by
\end{leancode}

\begin{dictionary}
\lean{Q.comap P.subtype} is $Q$ seen as a subspace of $P$ (the preimage under the inclusion
\lean{P.subtype}); the hypothesis \lean{_h : Q ≤ P} is recorded only to document the intended use. The
theorem is $\dim(P/Q)+\dim Q = \dim P$ for finite-dimensional $P$. This is the layer $\lelevengr_p =
F_p/F_{p+1}$ of a filtration, one weight space at a time.
\end{dictionary}

\textbf{Proof on paper.} Rank--nullity for the quotient map $P\to P/Q$. \qed

\begin{status}
Proved, unconditional; standard three.
\end{status}

\begin{physics}
The layers of a filtration are the ``classical'' pieces of the operator algebra: $\lelevengr_p$ is the
space of operators of exactly $p$ units of the filtration degree, modulo higher corrections. The
dimension identity is the bookkeeping that turns a filtered character into an ordinary one.
\end{physics}


%======================================================================
\section{Li's filtration and its associated graded algebra}\label{l11:sec:li}
%======================================================================

\subsection{Motivation}

Zhu's $C_2$-algebra $\lelevenR = V/C_2(V)$ throws away every state containing a mode $u_{(-n)}$ with
$n\ge2$, that is, every state containing a derivative of a field. Haisheng Li observed that $\lelevenR$
is only the bottom layer of a \emph{decreasing} filtration of $V$ that keeps track of \emph{how many}
derivatives a state contains:
\begin{equation}\label{l11:eq:lidef}
F^pV = \operatorname{span}\bigl\{u^1_{(-n_1-1)}u^2_{(-n_2-1)}\cdots u^r_{(-n_r-1)}\lelevenvac\ :\
r\ge0,\ u^i\text{ homogeneous},\ n_i\ge0,\ n_1+\cdots+n_r\ge p\bigr\}.
\end{equation}
Since $u_{(-n-1)}\lelevenvac = T^nu/n!$, the number $n_1+\cdots+n_r$ is the number of derivatives. Li
proved: $F^0V = V$, $F^pV\supseteq F^{p+1}V$, $F^1V = C_2(V)$, $T F^pV\subseteq F^{p+1}V$,
\begin{equation}\label{l11:eq:limult}
F^p_{(n)}F^q\subseteq F^{p+q-n-1}\ (n\in\Z),\qquad F^p_{(n)}F^q\subseteq F^{p+q-n}\ (n\ge0),
\end{equation}
(with $F^{k} = V$ for $k\le0$). Consequently $\lelevengr^FV = \bigoplus_pF^pV/F^{p+1}V$ is a commutative
algebra under the $(-1)$-product, $T$ induces a derivation of degree $1$, and the $(n)$-products with
$n\ge0$ induce a vertex Poisson structure. Its degree-zero piece is $F^0/F^1 = V/C_2(V) = \lelevenR$, and
it is generated by $\lelevenR$ as a differential algebra. That last statement is the source of the
surjection from the arc algebra $J_\infty(\lelevenR)$ onto $\lelevengr^FV$
(Section~\ref{l11:sec:jets}).

\subsection{The definition in the library}

The library's Li filtration is defined for an arbitrary weak module (it is the engine of Buhl's spanning
theorem, Chapter~\ref{ch:L15}) and specialised to the adjoint module. We quote the general definition,
which lives in a file of Chapter~\ref{ch:L15}, because every statement in this section depends on it.

In words: a \emph{letter} is a formal token $(c,\beta,n)$ standing for the mode $c_{(-n-1)}$ of a state $c$ of weight $\beta$; a
\emph{word} is an ordered list of letters, acting on a state with the first letter as the outermost operator; its \emph{level} is
$\sum n_i$, the total number of derivatives in the composite operator. (``Admissible'' below means only that the weight label is
correct; it has nothing to do with the ``admissible relation ideals'' of Section~\ref{l11:sec:dsk}.)

\begin{definition}[Letters, words, level, Li's filtration]\label{l11:def:gfil}
A \emph{letter} is a triple $(c,\beta,n)\in V\times\N\times\N$ standing for the mode $c_{(-n-1)}$; it is
\emph{admissible} if $c\in V_\beta$. A \emph{word} is a list of letters; it acts on a module $M$ by
applying the modes with the first letter outermost, and its \emph{level} is $\sum n_i$. For a set
$B\subseteq M$ and $p\in\Z$, $G^p(M,B)$ is the span of all admissible words of level $\ge\max(p,0)$
applied to elements of $B$. The module is \emph{generated by $B$} (\lean{LWGenerated}) if $G^0(M,B) = M$.
The \emph{Li filtration} of $V$ is $F^pV = G^p(V,\{\lelevenvac\})$ for the adjoint module.
\end{definition}
\noindent File \leanpath{VOA/\allowbreak Modularity/\allowbreak LiFiltration.lean} (Chapter~\ref{ch:L15}).
\begin{leancode}
abbrev Letter (𝒱 : Type u) [AddCommGroup 𝒱] [Module ℂ 𝒱] := 𝒱 × ℕ × ℕ

def letterOK (W : VertexOperatorAlgebra 𝒱) (p : Letter 𝒱) : Prop :=
  p.1 ∈ W.grading (p.2.1 : ℤ)

def WordOK (W : VertexOperatorAlgebra 𝒱) (l : List (Letter 𝒱)) : Prop := ∀ p ∈ l, letterOK W p

def wordAct (A : WeakModule W M) : List (Letter 𝒱) → M → M
  | [], x => x
  | p :: t, x => mode (A.YM p.1) (-(p.2.2 : ℤ) - 1) (wordAct A t x)

def wordLevel (l : List (Letter 𝒱)) : ℕ := (l.map (fun p => p.2.2)).sum

def Gset (A : WeakModule W M) (B : Set M) (p : ℕ) : Set M :=
  {x | ∃ (l : List (Letter 𝒱)) (w : M), WordOK W l ∧ w ∈ B ∧ p ≤ wordLevel l ∧
    x = wordAct A l w}

def Gfil (A : WeakModule W M) (B : Set M) (p : ℤ) : Submodule ℂ M :=
  Submodule.span ℂ (Gset A B p.toNat)

def LWGenerated (A : WeakModule W M) (B : Set M) : Prop := Gfil A B 0 = ⊤
\end{leancode}
\noindent File \leanpath{VOA/\allowbreak Filtered/\allowbreak LiAssociatedGraded.lean}.
\begin{leancode}
def liFiltration (W : VertexOperatorAlgebra 𝒱) (p : ℕ) : Submodule ℂ 𝒱 :=
  Gfil (adjointModule W).toWeakModule ({W.vacuum} : Set 𝒱) (p : ℤ)

def liWeightPiece (W : VertexOperatorAlgebra 𝒱) (p N : ℕ) : Submodule ℂ 𝒱 :=
  liFiltration W p ⊓ W.grading N

abbrev liWeightLayer (W : VertexOperatorAlgebra 𝒱) (p N : ℕ) : Type _ :=
  Submodule.Layer (P := liWeightPiece W p N) (Q := liWeightPiece W (p + 1) N)
    (liWeightPiece_antitone (W := W) (N := N) (Nat.le_succ p))
\end{leancode}

\begin{dictionary}
\lean{Letter 𝒱} is the product type $V\times\N\times\N$; \lean{p.1}, \lean{p.2.1}, \lean{p.2.2} are its
three components $c,\beta,n$. \lean{wordAct A (p :: t) x} applies \lean{mode (A.}\allowbreak\lean{YM p.}\allowbreak\lean{1) (-(p.}\allowbreak\lean{2.}\allowbreak\lean{2) - 1)},
which is the module mode $c^M_{(-n-1)}$, to \lean{wordAct A t x}: the head of the list is the outermost
operator. \lean{wordLevel} is $\sum n_i$. In \lean{Gset A B p}, the existential packages a word \lean{l}
and a bottom vector \lean{w}, with \lean{WordOK W l} (all letters homogeneous of the recorded weight), $w\in
B$ and level $\ge p$. \lean{Gfil} takes the span and converts the integer index with \lean{p.toNat} (so
every $p\le0$ gives the level-$0$ piece). With \lean{A = adjointWeakModule W} (the module $V$ itself, its
modes being \lean{W.smul}) and $B = \{\lelevenvac\}$ this is exactly Li's (\ref{l11:eq:lidef}); the
weight label $\beta$ is redundant data that makes ``homogeneous'' checkable. \lean{liWeightPiece W p N} is
$F^pV\cap V_N$ and \lean{liWeightLayer W p N} is $(F^pV\cap V_N)/(F^{p+1}V\cap V_N)$, the
$(p,N)$-component of $\lelevengr^FV$, built with the exact layer of Definition~\ref{l11:def:layer}.

\textbf{Same as Li?} Yes: Li (2005, §2) defines $F_pV$ as the span of the same
monomials, with $u^i$ arbitrary (not necessarily homogeneous); since every $u$ is a finite sum of
homogeneous components and modes are linear in $u$, the spans agree. The library's convention ``weights
in $\N$'' is harmless because a library VOA has no negative weights.
\end{dictionary}

\begin{warning}
\textbf{A redundant hypothesis (minor) --- since removed.} Two theorems of
\leanpath{VOA/\allowbreak Filtered/\allowbreak LiAssociatedGraded.lean},
\lean{liWeightPiece_}\allowbreak\lean{zero_}\allowbreak\lean{eq_}\allowbreak\lean{grading} and \lean{sum_}\allowbreak\lean{liWeightLayer_}\allowbreak\lean{finrank_}\allowbreak\lean{eq_}\allowbreak\lean{grading}, carry the
hypothesis \lean{hLW : LWGenerated (adjointWeakModule W) ({W.vacuum} : Set 𝒱)} when this chapter was written; the library
has since dropped it from both (the statements quoted in this chapter are the current ones). This hypothesis is always
true: the library proves it as \lean{lwGenerated_adjoint} (in
\leanpath{VOA/\allowbreak Modularity/\allowbreak BuhlSpanningData.lean}), and in the same chapter
\lean{liFiltration_}\allowbreak\lean{zero_}\allowbreak\lean{eq_}\allowbreak\lean{top} is proved without it. A Lean check made for this chapter (not part of this release)
states the first theorem without \lean{hLW}. Nothing
becomes false; a reader should simply know that the hypothesis costs nothing.
\end{warning}

\begin{physics}
Li's filtration counts derivatives. A state $u^1_{(-n_1-1)}\cdots u^r_{(-n_r-1)}\lelevenvac$ is the
normal-ordered product ${:}\partial^{n_1}u^1\cdots\partial^{n_r}u^r{:}$ (up to factorials), so $F^pV$ is
spanned by composite operators carrying at least $p$ derivatives in total. In a semiclassical (large
central charge, or $\hbar\to0$) picture the fields become classical functions on a phase space and
derivatives become jet coordinates; $\lelevengr^FV$ is the algebra of classical local functionals, and
$F^0/F^1 = \lelevenR$ is the algebra of functions on the space of constant field configurations --- the
classical moduli space, the associated variety.
\end{physics}

\begin{theorem}[Finiteness and the weight-wise telescoping identity]\label{l11:thm:litele}
For all $p,N\in\N$: $F^pV\cap V_N$ is finite-dimensional; $F^pV\cap V_N = 0$ if $N<p$; and
\begin{equation*}
\sum_{p=0}^{N}\dim\bigl(F^pV\cap V_N\big/F^{p+1}V\cap V_N\bigr) = \dim V_N .
\end{equation*}
\end{theorem}
\noindent File \leanpath{VOA/\allowbreak Filtered/\allowbreak LiAssociatedGraded.lean}.
\begin{leancode}
theorem liWeightPiece_finiteDimensional (W : VertexOperatorAlgebra 𝒱) (p N : ℕ) :
    FiniteDimensional ℂ (liWeightPiece W p N) := by

theorem liWeightPiece_eq_bot_of_lt (W : VertexOperatorAlgebra 𝒱) {p N : ℕ} (h : N < p) :
    liWeightPiece W p N = ⊥ := by

theorem sum_liWeightLayer_finrank_eq_grading (W : VertexOperatorAlgebra 𝒱) (N : ℕ) :
    (∑ p ∈ Finset.range (N + 1), Module.finrank ℂ (liWeightLayer W p N)) =
      Module.finrank ℂ (W.grading N) := by
\end{leancode}

\begin{dictionary}
\lean{Finset.range (N + 1)} is $\{0,\dots,N\}$; the conclusion is exactly the displayed identity, with
\lean{W.grading N} the weight space $V_N$ (the cast $\N\to\Z$ is implicit). There is no further hypothesis (the
redundant \lean{hLW} of an earlier version has been removed; warning above). The identity says that the $N$-th coefficient of the character of $V$ equals
the total dimension of the weight-$N$ part of $\lelevengr^FV$: passing to the associated graded preserves
characters.
\end{dictionary}

\textbf{Proof on paper.} $F^pV\cap V_N\subseteq V_N$, which is finite-dimensional. A word of level $\ge p$
raises weight by $\sum(\beta_i+n_i)\ge p$, so it lands in weights $\ge p$; hence $F^pV\cap V_N = 0$ for
$N<p$. Telescoping the exact-layer identity of Definition~\ref{l11:def:layer} from $p=0$ (where
$F^0V\cap V_N = V_N$) to $p=N+1$ (where the piece is $0$) gives the sum. \qed

\begin{status}
Proved, unconditional (no hypotheses beyond the VOA $W$); \texttt{\#print axioms
VOA.sum\_liWeightLayer\_finrank\_eq\_grading}: standard three.
\end{status}

\begin{physics}
Taking the classical limit does not change the counting of states: the partition function of the chiral
algebra equals that of its classical shadow, graded by weight. Refining it by the Li degree gives a
two-variable character, which is what the ``Hilbert squeeze'' of Section~\ref{l11:sec:jets} compares.
\end{physics}

\subsection{Multiplicativity}

\begin{theorem}[Li's multiplicativity estimates]\label{l11:thm:limult}
Let $x\in F^pV$, $y\in F^qV$. Then
\begin{enumerate}
\item for every $N\in\Z$, $x_{(N)}y\in F^{p+q-N-1}V$ (indices $\le0$ meaning $V$);
\item $x_{(-1)}y\in F^{p+q}V$ and, more generally, $x_{(-d-1)}y\in F^{p+q+d}V$ for $d\in\N$;
\item for $n\in\N$ with $n\le p+q$, $x_{(n)}y\in F^{p+q-n}V$;
\item $Tx\in F^{p+1}V$; and $F^0V = V$.
\end{enumerate}
\end{theorem}
\noindent File \leanpath{VOA/\allowbreak Filtered/\allowbreak LiMultiplicativity.lean}.
\begin{leancode}
theorem adjointGfil_smul_mem (W : VertexOperatorAlgebra V)
    (hLW : LWGenerated (adjointWeakModule W) ({W.vacuum} : Set V))
    {p q N : ℤ} {x y : V}
    (hx : x ∈ Gfil (adjointWeakModule W) ({W.vacuum} : Set V) p)
    (hy : y ∈ Gfil (adjointWeakModule W) ({W.vacuum} : Set V) q) :
    W.smul x N y ∈
      Gfil (adjointWeakModule W) ({W.vacuum} : Set V) (p + q - N - 1) := by

theorem liFiltration_smul_negSucc_mem (W : VertexOperatorAlgebra V)
    {p q : ℕ} (d : ℕ) {x y : V}
    (hx : x ∈ liFiltration W p) (hy : y ∈ liFiltration W q) :
    W.smul x (-(d : ℤ) - 1) y ∈ liFiltration W (p + q + d) := by

theorem liFiltration_smul_nonneg_mem (W : VertexOperatorAlgebra V)
    {p q n : ℕ} (hn : n ≤ p + q) {x y : V}
    (hx : x ∈ liFiltration W p) (hy : y ∈ liFiltration W q) :
    W.smul x (n : ℤ) y ∈ liFiltration W (p + q - n) := by

theorem liFiltration_L_neg_one_mem (W : VertexOperatorAlgebra V) {p : ℕ} {x : V}
    (hx : x ∈ liFiltration W p) : W.L (-1) x ∈ liFiltration W (p + 1) := by

theorem liFiltration_zero_eq_top (W : VertexOperatorAlgebra V) : liFiltration W 0 = ⊤ := by
\end{leancode}

\begin{dictionary}
(1) is stated with integer indices for the general \lean{Gfil}; \lean{(p + q - N - 1)} is an integer and
\lean{Gfil} of a non-positive index is everything. It carries the always-true \lean{hLW}. (3) uses
natural-number subtraction \lean{p + q - n}, which is why the hypothesis \lean{hn : n ≤ p + q} is there:
without it \lean{p + q - n} would truncate to $0$ and the statement would be weaker but still true. (2)
with \lean{d = 0} is \lean{liFiltration_}\allowbreak\lean{smul_}\allowbreak\lean{neg_}\allowbreak\lean{one_}\allowbreak\lean{mem}. All four are exactly Li's (\ref{l11:eq:limult})
(Li 2005, §2; Arakawa 2012, §2). (1) is the ``all modes'' estimate; (3) is the sharper
non-negative-mode estimate, one step better than (1) at $N = n\ge0$, which is what makes the bracket on
$\lelevengr^FV$ have degree $-n$ rather than $-n-1$.
\end{dictionary}

\textbf{Proof on paper.} By linearity take $x$ a word $a^1_{(-m_1-1)}\cdots a^r_{(-m_r-1)}\lelevenvac$ of
level $\ell\ge p$. The key one-letter estimate (\lean{act_mem_Gfil}, Chapter~\ref{ch:L15}) is that for
homogeneous $u$, $u_{(k)}F^r\subseteq F^{r-k-1}$: induction on the word that spans the vector acted on,
using the commutator formula $[u_{(k)},a_{(-m-1)}] = \sum_j\binom kj(u_{(j)}a)_{(k-m-1-j)}$, whose
correction terms land one level better. For a composite $x$, write $x = a_{(-m-1)}x'$ and use the iterate
(associativity) formula
\begin{equation*}
(a_{(-m-1)}x')_{(N)}y = \sum_{i\ge0}(-1)^i\tbinom{-m-1}{i}\Bigl(a_{(-m-1-i)}x'_{(N+i)}y-(-1)^{-m-1}x'_{(-m-1+N-i)}a_{(i)}y\Bigr),
\end{equation*}
induct on the length of $x$ and apply the one-letter estimate to each term; the levels add up as claimed.
For $Tx$: $Tx = x_{(-2)}\lelevenvac$ and $\lelevenvac\in F^0$, so (2) with $d=1$ gives $Tx\in F^{p+1}$. For
(3), Li's argument uses skew-symmetry $x_{(n)}y = \sum_j(-1)^{n+j+1}T^{(j)}(y_{(n+j)}x)$ together with the
fact that $T$ raises the filtration by one, which recovers the lost step. \qed

\begin{status}
Proved, unconditional; \texttt{\#print axioms} of \lean{VOA.}\allowbreak\lean{adjointGfil_}\allowbreak\lean{smul_}\allowbreak\lean{mem},
\lean{VOA.}\allowbreak\lean{liFiltration_}\allowbreak\lean{smul_}\allowbreak\lean{neg_}\allowbreak\lean{one_}\allowbreak\lean{mem}, \lean{VOA.}\allowbreak\lean{liFiltration_}\allowbreak\lean{smul_}\allowbreak\lean{nonneg_}\allowbreak\lean{mem}: standard three.
\end{status}

\begin{physics}
Each derivative carried by a field survives in the product, and a singular OPE term (a mode $u_{(n)}$ with
$n\ge0$) can only \emph{remove} derivatives (it is a contraction, lowering the classical degree). That is
exactly why the leading, classical part of an OPE is a commutative product (from $u_{(-1)}$) and a Poisson
bracket (from $u_{(0)}$), and why derivatives act as a derivation.
\end{physics}

\subsection{The associated graded algebra}

\begin{theorem}[The Li associated graded is a commutative differential algebra]\label{l11:thm:ligr}
Let $\lelevengr^FV = \bigoplus_{(p,N)\in\N^2}\lelevengr^F_{p,N}V$, where $\lelevengr^F_{p,N}V$ is the
layer of Theorem~\ref{l11:thm:litele}. The $(-1)$-product descends to bilinear maps
$\lelevengr^F_{p,N}\times\lelevengr^F_{q,M}\to\lelevengr^F_{p+q,N+M}$, which make $\lelevengr^FV$ a
commutative, associative, unital, bigraded $\C$-algebra with unit the class of $\lelevenvac$; $T$ descends to
a linear map of bidegree $(1,1)$ which is a derivation; each component is finite-dimensional and
$\lelevengr^F_{p,N}V = 0$ when $N<p$.
\end{theorem}
\noindent Files \leanpath{VOA/\allowbreak Filtered/\allowbreak LiMultiplicativity.lean},
\leanpath{VOA/\allowbreak Filtered/\allowbreak LiGradedAlgebra.lean}.
\begin{leancode}
noncomputable def liWeightLayerMul (W : VertexOperatorAlgebra V) (p q N M : ℕ) :
    liWeightLayer W p N →ₗ[ℂ]
      liWeightLayer W q M →ₗ[ℂ] liWeightLayer W (p + q) (N + M) :=

def LiNegOneCommutatorEstimate (W : VertexOperatorAlgebra V) : Prop :=
  ∀ {p q : ℕ} {x y : V}, x ∈ liFiltration W p → y ∈ liFiltration W q →
    W.smul x (-1) y - W.smul y (-1) x ∈ liFiltration W (p + q + 1)

theorem liNegOneCommutatorEstimate_proven (W : VertexOperatorAlgebra V) :
    LiNegOneCommutatorEstimate W := by

theorem liWeightLayerMul_comm (W : VertexOperatorAlgebra V) (p q N M : ℕ)
    (xbar : liWeightLayer W p N) (ybar : liWeightLayer W q M) :
    HEq (liWeightLayerMul W p q N M xbar ybar)
      (liWeightLayerMul W q p M N ybar xbar) := by

abbrev liAssociatedGraded (W : VertexOperatorAlgebra V) : Type u :=
  ⨁ d : ℕ × ℕ, liWeightLayer W d.1 d.2

noncomputable instance liWeightLayerGCommRing (W : VertexOperatorAlgebra V) :
    DirectSum.GCommRing (fun d : ℕ × ℕ => liWeightLayer W d.1 d.2) where

@[simp] theorem liAssociatedGraded_of_mul_of (W : VertexOperatorAlgebra V)
    (p q N M : ℕ) (x : liWeightLayer W p N) (y : liWeightLayer W q M) :
    (DirectSum.of (fun d : ℕ × ℕ => liWeightLayer W d.1 d.2) (p, N) x :
        liAssociatedGraded W) *
      DirectSum.of (fun d : ℕ × ℕ => liWeightLayer W d.1 d.2) (q, M) y =
    DirectSum.of (fun d : ℕ × ℕ => liWeightLayer W d.1 d.2) (p + q, N + M)
      (liWeightLayerMul W p q N M x y) := by

noncomputable def liAssociatedGradedTranslation (W : VertexOperatorAlgebra V) :
    liAssociatedGraded W →ₗ[ℂ] liAssociatedGraded W :=

theorem liAssociatedGradedTranslation_leibniz (W : VertexOperatorAlgebra V)
    (a b : liAssociatedGraded W) :
    liAssociatedGradedTranslation W (a * b) =
      liAssociatedGradedTranslation W a * b + a * liAssociatedGradedTranslation W b := by

theorem liAssociatedGraded_component_eq_zero_of_weight_lt_degree
    (W : VertexOperatorAlgebra V) {p N : ℕ} (h : N < p)
    (x : liWeightLayer W p N) : x = 0 := by
\end{leancode}

\begin{dictionary}
\lean{liWeightLayerMul W p q N M} is the descended product, a bilinear map (two nested \lean{→ₗ[ℂ]}). The
commutativity statement \lean{liWeightLayerMul_comm} uses \lean{HEq} (``heterogeneous equality'') because
the two sides live in the layers $(p+q,N+M)$ and $(q+p,M+N)$, which are equal but not \emph{definitionally}
the same type in Lean; \lean{HEq} is the honest way to say ``equal after identifying the indices''. The
direct sum \lean{⨁ d : ℕ × ℕ, …} is the external direct sum of the layers (finitely supported families).
\lean{DirectSum.GCommRing} is Mathlib's notion of a graded commutative ring structure on a family indexed by
an additive monoid; with it, Mathlib makes the direct sum a commutative ring, and
\lean{liAssociatedGraded_}\allowbreak\lean{of_}\allowbreak\lean{mul_}\allowbreak\lean{of} says the product of two homogeneous elements is the homogeneous element
given by \lean{liWeightLayerMul}. A further instance (\lean{liWeightLayerGAlgebra}) makes it a
$\C$-algebra. \lean{liAssociatedGradedTranslation} is assembled from the descended $L_{-1}$ on each layer,
shifting bidegree by $(1,1)$; \lean{_leibniz} is the derivation property. \lean{LiNegOneCommutatorEstimate}
(the one-step commutator estimate $x_{(-1)}y - y_{(-1)}x\in F^{p+q+1}$, which is what commutativity of the
graded product means) is introduced as a \lean{Prop} and then proved
(\lean{liNegOneCommutatorEstimate_}\allowbreak\lean{proven}), so it is not an open hypothesis.

\textbf{What is not packaged.} Li's vertex Poisson structure --- the brackets induced by the $(n)$-products,
$n\ge0$, of degree $-n$ --- is not defined on the library's $\lelevengr^FV$. The estimate that would make it
well defined is Theorem~\ref{l11:thm:limult}(3); the descent is not carried out. So the library's
associated graded is a commutative differential algebra, not yet a vertex Poisson algebra. In particular nothing in this chapter
gives the associated variety its Poisson structure; the Poisson structure on the degree-zero part $\lelevenR$ alone is
Chapter~\ref{ch:L12}, Theorem~\ref{l12:thm:poisson}.
\end{dictionary}

\textbf{Proof on paper.} The descended product is well defined by Theorem~\ref{l11:thm:limult}(2): if $x$
or $y$ moves by an element of the next filtration step, the product moves by an element of $F^{p+q+1}$.
Commutativity: by skew-symmetry $x_{(-1)}y - y_{(-1)}x = \sum_{j\ge1}(-1)^{j+1}T^{(j)}(y_{(j-1)}x)$, and
each term lies in $F^{p+q+1}$ because $y_{(j-1)}x\in F^{p+q-j+1}$ (part (3)) and $T^{(j)}$ raises by $j$.
Associativity: the iterate formula at $(-1,-1)$ writes $(x_{(-1)}y)_{(-1)}z - x_{(-1)}(y_{(-1)}z)$ as a
sum of terms $x_{(-1-i)}y_{(-1+i)}z$, $i\ge1$, and their mirror images, each one step higher. The unit is
$[\lelevenvac]$ since $\lelevenvac_{(-1)} = \mathrm{id}$. $T$ is a derivation of the $(-1)$-product
($[T,x_{(-1)}] = (Tx)_{(-1)}$) and raises the filtration by one. \qed

\begin{status}
Proved, unconditional; \texttt{\#print axioms} of \lean{VOA.}\allowbreak\lean{liWeightLayerMul_}\allowbreak\lean{comm},
\lean{VOA.}\allowbreak\lean{liWeightLayerGCommRing}, \lean{VOA.}\allowbreak\lean{liAssociatedGradedTranslation_}\allowbreak\lean{leibniz}: standard three.
\end{status}

\begin{physics}
This is the statement that the classical limit of a chiral algebra is a commutative algebra of classical
fields with a derivative: normal ordering becomes multiplication, and $\partial$ acts by the Leibniz rule. The
missing Poisson bracket is the classical remnant of the singular OPE, $\{a_\lambda b\}$ in the language of
Poisson vertex algebras. (Update: the Poisson vertex algebra structure on $\operatorname{gr}^FV$ has since been
formalized, \lean{liPoissonVertexAlgebra}, as part of the Arakawa--Dai campaign; see Chapter~\ref{ch:L20}. The statement
quoted here is unchanged.)
\end{physics}

\subsection{The bottom layer is Zhu's \texorpdfstring{$C_2$}{C2}-algebra}

\begin{theorem}[Li: $F^1V = C_2(V)$; $\lelevengr^F_0V = \lelevenR$]\label{l11:thm:f1}
$C_2(V) = F^1V$. Consequently the degree-zero layers $\bigoplus_N\lelevengr^F_{0,N}V$ form the algebra
$\lelevenR = V/C_2(V)$, and the class map $V_N\to\lelevengr^F_{0,N}V$ kills $C_2(V)\cap V_N$.
\end{theorem}
\noindent File \leanpath{VOA/\allowbreak Filtered/\allowbreak ClassicalFreeness.lean} (the inclusion $C_2\subseteq F^1$).
\begin{leancode}
theorem O₂_le_liFiltration_one (W : VertexOperatorAlgebra V) :
    O₂ W ≤ liFiltration W 1 := by

noncomputable def liDegreeZeroClass (W : VertexOperatorAlgebra V) (N : ℕ)
    (x : W.grading (N : ℤ)) : liAssociatedGraded W :=
  DirectSum.of (fun d : ℕ × ℕ => liWeightLayer W d.1 d.2) (0, N)

theorem liDegreeZeroClass_eq_zero_of_mem_O₂ (W : VertexOperatorAlgebra V) (N : ℕ)
    (x : W.grading (N : ℤ)) (hx : x.1 ∈ O₂ W) :
    liDegreeZeroClass W N x = 0 := by
\end{leancode}
\noindent File \leanpath{VOA/\allowbreak Filtered/\allowbreak LiC2Bridge.lean} (the reverse inclusion and the equality, now
library theorems):
\begin{leancode}
theorem liFiltration_one_le_O₂ (W : VertexOperatorAlgebra V) :
    liFiltration W 1 ≤ O₂ W := by

theorem liFiltration_one_eq_O₂ (W : VertexOperatorAlgebra V) :
    liFiltration W 1 = O₂ W :=
\end{leancode}

\begin{dictionary}
\lean{O₂ W} is $C_2(V)$. \lean{liDegreeZeroClass W N x} is the class of a weight-$N$ vector in the
$(0,N)$-component of \lean{liAssociatedGraded W}. The inclusion $C_2(V)\subseteq F^1V$ gives
a surjection $\lelevenR\twoheadrightarrow\bigoplus_N\lelevengr^F_{0,N}V$. Li's theorem is the equality, which
upgrades the surjection to an isomorphism; the library proves it (\lean{liFiltration_one_eq_O₂}, quoted above; the
same file also gives \lean{liDegreeZeroClass_eq_zero_iff_mem_O₂} and the linear equivalence
\lean{liZerothQuotientEquivC2}), and a Lean check made for this chapter outside the library, when the library had
only one inclusion, checked it first. With the equality, the
library's $\lelevengr^F_0V$ is literally Zhu's \lean{A₂ W}, and the product on it is $[u][v] = [u_{(-1)}v]$ ---
the same product as the \lean{CommRing (A₂ W)} instance of Chapter~\ref{ch:L12}.
\end{dictionary}

\textbf{Proof on paper.} $C_2\subseteq F^1$: $u_{(-2)}v$ is $u_{(-2)}$ applied to $v\in F^0$, and a letter
with $n=1$ raises the level by one (Theorem~\ref{l11:thm:limult}(2) with $d=1$). $F^1\subseteq C_2$ (\lean{liFiltration_one_le_O₂}; first in our
check outside the library): a word $a^1_{(-n_1-1)}\cdots a^r_{(-n_r-1)}\lelevenvac$ of level $\ge1$ has some $n_i\ge1$. If
the outermost letter has $n_1\ge1$, the word is $a^1_{(-n_1-1)}(\dots)$ with $-n_1-1\le-2$, which lies in
$C_2(V)$ (Zhu: $u_{(-k)}v\in C_2(V)$ for all $k\ge2$). Otherwise $n_1 = 0$, the remaining word has level
$\ge1$ and lies in $C_2(V)$ by induction, and $C_2(V)$ is stable under $a_{(-1)}$ (it is an ideal for the
$(-1)$-product, Theorem~\ref{l12:thm:poisson}). \qed

\begin{status}
\textbf{Proved in the library: both inclusions and the equality.} $C_2(V)\subseteq F^1V$ is
\lean{O₂_le_liFiltration_one} (\leanpath{Filtered/\allowbreak ClassicalFreeness.lean}); $F^1V\subseteq C_2(V)$ and the
equality $F^1V = C_2(V)$ are \lean{VOA.}\allowbreak\lean{liFiltration_}\allowbreak\lean{one_}\allowbreak\lean{le_}\allowbreak\lean{O₂} and \lean{VOA.}\allowbreak\lean{liFiltration_}\allowbreak\lean{one_}\allowbreak\lean{eq_}\allowbreak\lean{O₂}
(\leanpath{Filtered/\allowbreak LiC2Bridge.lean}), for every VOA, with no hypothesis. When this chapter was first written, at an earlier state of the library, the equality
existed only in a Lean check made for this chapter outside the library; that gap has
since been closed in the library. \texttt{\#print axioms} for the library declarations of \leanpath{LiC2Bridge.lean} was not
re-run for this chapter. Inclusion $C_2\subseteq F^1$: \texttt{\#print axioms} \lean{VOA.}\allowbreak\lean{O₂_}\allowbreak\lean{le_}\allowbreak\lean{liFiltration_}\allowbreak\lean{one} and
\lean{VOA.}\allowbreak\lean{liDegreeZeroClass_}\allowbreak\lean{eq_}\allowbreak\lean{zero_}\allowbreak\lean{of_}\allowbreak\lean{mem_}\allowbreak\lean{O₂}: standard three. Reverse inclusion and equality, version
checked outside the library: compiled with \texttt{lake env lean -DmaxSynthPendingDepth=3} for this chapter; standard three.
\end{status}

\begin{physics}
$F^1 = C_2$ says: ``a composite operator has at least one derivative'' and ``it is a normal-ordered product
with a derivative field in it'' are the same condition. So the bottom of the classical tower, the functions
on constant field configurations, is exactly Zhu's $C_2$-algebra, and $X_V = \operatorname{Specm}\lelevenR$
is the space of classical vacua. For the free boson, $\lelevenR = \C[J]$ and $X_V = \C$ (a line of vacua
$\langle J\rangle$); for a lisse ($C_2$-cofinite) theory, $X_V$ is a point. For the chiral algebra of a 4d
$\mathcal N=2$ theory, Beem--Rastelli's conjecture is $X_V\cong$ the Higgs branch.
\end{physics}


%======================================================================
\section{Jets, presentations of commutative algebras, and classical freeness}\label{l11:sec:jets}
%======================================================================

\subsection{The commutative-algebra side}

If $\lelevenR = \C[x_1,\dots,x_g]/I$, the \emph{arc algebra} (infinite jet algebra) is
$J_\infty\lelevenR = \C[x_i^{(r)} : 1\le i\le g,\ r\ge0]/J_\infty I$, where $J_\infty I$ is generated by the
coefficients of $t^n$ in $f\bigl(\sum_rx_1^{(r)}t^r,\dots,\sum_rx_g^{(r)}t^r\bigr)$ for $f\in I$, $n\ge0$. It is
the universal commutative differential algebra generated by $\lelevenR$: the derivation is
$\partial x_i^{(r)} = (r+1)x_i^{(r+1)}$. Li and Arakawa show that there is a canonical surjection of differential
algebras $J_\infty\lelevenR\twoheadrightarrow\lelevengr^FV$ extending $\lelevenR = \lelevengr^F_0V$; $V$ is called
\emph{classically free} when it is an isomorphism: the classical limit $\lelevengr^FV$ is then exactly the arc space of the
associated variety, i.e.\ $V$ has no relations in the classical limit beyond the derivatives of those of $\lelevenR$. The library works with the finite truncations of order $m$
(variables $x_i^{(r)}$ with $r\le m$), which is enough for any fixed Li degree $p\le m$.

\begin{definition}[Presentations of a commutative algebra]\label{l11:def:compres}
For an algebra map $\varphi:A\to B$ and an ideal $I\subseteq A$: the relations \emph{vanish} if
$I\subseteq\ker\varphi$; the presentation is \emph{exact} if $\ker\varphi = I$.
\end{definition}
\noindent File \leanpath{VOA/\allowbreak Commutative/\allowbreak Presentation.lean}.
\begin{leancode}
def RelationsVanish (I : Ideal A) (φ : A →ₐ[R] B) : Prop :=
  I ≤ RingHom.ker φ

def IsExactPresentation (I : Ideal A) (φ : A →ₐ[R] B) : Prop :=
  RingHom.ker φ = I

noncomputable def quotientEquiv (I : Ideal A) (φ : A →ₐ[R] B)
    (hsurj : Function.Surjective φ) (hexact : IsExactPresentation I φ) :
    (A ⧸ I) ≃ₐ[R] B := by
\end{leancode}

\begin{dictionary}
\lean{RingHom.ker φ} is the kernel. \lean{(A ⧸ I) ≃ₐ[R] B} is an isomorphism of $R$-algebras. The theorem
\lean{quotientEquiv} is the first isomorphism theorem: a surjective $\varphi$ with exact presentation $I$ gives
$A/I\cong B$. These are the textbook notions, kept separate so that ``the relations hold'' is never confused
with ``these are all the relations''.
\end{dictionary}

\textbf{Proof on paper.} First isomorphism theorem. \qed

\begin{status}
Proved, unconditional; \texttt{\#print axioms} \lean{VOA.}\allowbreak\lean{Presentation.}\allowbreak\lean{quotientEquiv}: standard three.
\end{status}

\begin{physics}
``Relations vanish'' is the statement that certain composite operators are null; ``exact presentation'' is the
much stronger claim that the listed null relations generate \emph{all} of them --- that one knows the complete
set of equations of the associated variety (for a Higgs branch: the complete chiral ring relations).
\end{physics}

\begin{definition}[Finite jets; the jet ideal; presented jets]\label{l11:def:jets}
Fix generators $\sigma$ and an order $m\in\N$. The \emph{jet polynomial ring} is
$\C[x_i^{(r)} : i\in\sigma,\ 0\le r\le m]$; the truncated series of $x_i$ is $x_i(t) = \sum_{r=0}^mx_i^{(r)}t^r$;
for $f\in\C[x_i]$ the \emph{relation series} is $f(x(t))$ and its $n$-th coefficient is the Hasse--Schmidt
coefficient $f^{(n)}$. For an ideal $I$, the \emph{order-$m$ jet ideal} $J_mI$ is generated by $f^{(j)}$ for $f\in I$,
$0\le j\le m$, and the \emph{presented jet algebra} is $J_m(\C[x]/I) := \C[x_i^{(r)}]/J_mI$.
\end{definition}
\noindent File \leanpath{VOA/\allowbreak Commutative/\allowbreak FiniteJet.lean}.
\begin{leancode}
structure JetVar (σ : Type uσ) (m : ℕ) where
  generator : σ
  order : Fin (m + 1)

abbrev JetPoly (R : Type uR) [CommSemiring R] (σ : Type uσ) (m : ℕ) :=
  MvPolynomial (JetVar σ m) R

noncomputable def jetVariable (m : ℕ) (i : σ) (r : ℕ) : JetPoly R σ m :=
  if h : r ≤ m then
    MvPolynomial.X ⟨i, ⟨r, Nat.lt_succ_iff.mpr h⟩⟩
  else
    0

noncomputable def jetSeries (m : ℕ) (i : σ) : Polynomial (JetPoly R σ m) :=
  ∑ r ∈ Finset.range (m + 1), Polynomial.monomial r (jetVariable (R := R) m i r)

noncomputable def jetSubst (m : ℕ) :
    MvPolynomial σ R →ₐ[R] Polynomial (JetPoly R σ m) :=
  MvPolynomial.aeval (jetSeries (R := R) m)

noncomputable def relationCoeff (m : ℕ) (n : ℕ) (f : MvPolynomial σ R) :
    JetPoly R σ m :=
  (relationSeries (R := R) m f).coeff n

noncomputable def jetRelationSet (m : ℕ) (I : Ideal (MvPolynomial σ R)) :
    Set (JetPoly R σ m) :=
  Set.range fun fj : I × Fin (m + 1) => jetCoeff (R := R) m fj.2 fj.1.1

noncomputable def jetIdeal (m : ℕ) (I : Ideal (MvPolynomial σ R)) :
    Ideal (JetPoly R σ m) :=
  Ideal.span (jetRelationSet (R := R) m I)

abbrev PresentedJet (m : ℕ) (I : Ideal (MvPolynomial σ R)) :=
  JetPoly R σ m ⧸ jetIdeal (R := R) m I
\end{leancode}

\begin{dictionary}
\lean{JetVar σ m} is the set of pairs $(i,r)$ with $r\in$ \lean{Fin (m + 1)} $=\{0,\dots,m\}$, and \lean{JetPoly ℂ σ m}
the polynomial ring on them. \lean{jetVariable m i r} is $x_i^{(r)}$ for $r\le m$ and $0$ for $r>m$ (a convenient
junk value). \lean{jetSeries} is the polynomial $\sum_{r\le m}x_i^{(r)}t^r\in(\text{JetPoly})[t]$
(\lean{Polynomial.}\allowbreak\lean{monomial r a} is $at^r$). \lean{jetSubst} substitutes these series (\lean{MvPolynomial.aeval}
is evaluation). \lean{relationCoeff m n f} is the coefficient of $t^n$ in $f(x(t))$. \lean{jetRelationSet} ranges over
pairs (element of $I$, order $j\le m$), and \lean{jetCoeff m j f} is \lean{relationCoeff m j f}
(\lean{jetCoeff_apply}). \lean{Ideal.span} is the ideal generated.

\textbf{Same as the literature?} Yes, for the order-$m$ jet scheme $J_m(\operatorname{Spec}\C[x]/I)$: the ideal of
$J_m$ is generated by the coefficients of $t^0,\dots,t^m$ of $f(x(t))$ (Ishii; Arakawa--Moreau). The normalisation
$x(t) = \sum x^{(r)}t^r$ means $x^{(r)}$ corresponds to $\partial^rx/r!$ (divided powers), which matches the library's
divided translates $T^r a/r!$ below. The file \leanpath{FiniteJet.lean} has no transition map $J_{m+1}\to J_m$ between finite orders and no
infinite arc algebra (its docstring says so); fixed-weight stabilisation (next theorem) replaces the transition maps. The
infinite arc algebra \lean{InfiniteJet.PresentedJet} is introduced elsewhere in the library and documented in
Section~\ref{l11:sec:genwords-surj} (Definition~\ref{l11:def:infjetmap}).
\end{dictionary}

\begin{physics}
A jet of order $m$ of a field configuration $x(z)$ is its Taylor polynomial $\sum_{r\le m}x^{(r)}z^r$; the jet
scheme of the associated variety is the space of such Taylor polynomials that satisfy the classical equations
$f(x(z)) = 0$ to order $m$. In a chiral algebra, $x_i^{(r)}$ is the classical shadow of the field
$\partial^rx_i/r!$, and the jet relations are the derivatives of the classical null relations --- the equations of
motion and all their $z$-derivatives.
\end{physics}

\begin{theorem}[Weighted jets: homogeneity, independence of generators, stabilisation]\label{l11:thm:jetweights}
Give $x_i^{(r)}$ the weight $w_i+r$, with all $w_i>0$. Then
\begin{enumerate}
\item if $f$ is weighted-homogeneous of degree $D$, then $f^{(n)}$ is homogeneous of weight $D+n$; hence for a set
  $S$ of homogeneous relations, $J_m(\operatorname{span}S)$ is a homogeneous ideal;
\item $J_m(\text{ideal generated by }S)$ equals the ideal generated by the $f^{(j)}$, $f\in S$, $j\le m$;
\item for weight $N\le m$, the order-raising map identifies the weight-$N$ parts of the jet polynomial rings of
  orders $m$ and $m+1$, and (for homogeneous relations) the weight-$N$ images in the presented jet algebras;
\item with finitely many generators, every bidegree (Li degree, weight) piece of the presented jet algebra is
  finite-dimensional.
\end{enumerate}
\end{theorem}
\noindent Files \leanpath{VOA/\allowbreak Commutative/\allowbreak FiniteWeightedJetRelations.lean},
\leanpath{VOA/\allowbreak Commutative/\allowbreak FiniteWeightedJet.lean},
\leanpath{VOA/\allowbreak Commutative/\allowbreak MvPolynomialFiniteHomogeneous.lean}.
\begin{leancode}
theorem relationCoeff_isWeightedHomogeneous (baseWeight : σ → ℕ) (m D n : ℕ)
    {f : MvPolynomial σ R} (hf : f.IsWeightedHomogeneous baseWeight D) :
    (relationCoeff (R := R) m n f).IsWeightedHomogeneous

theorem jetIdeal_span_isHomogeneous (baseWeight : σ → ℕ) (m : ℕ)
    (S : Set (MvPolynomial σ R)) (hS : IsWeightedHomogeneousRelationSet baseWeight S) :
    (jetIdeal (R := R) m (Ideal.span S)).IsHomogeneous

theorem jetIdeal_span_eq_jetIdealOfSet (m : ℕ) (S : Set (MvPolynomial σ R)) :
    jetIdeal (R := R) m (Ideal.span S) = jetIdealOfSet (R := R) m S := by

noncomputable def jetWeightPieceOrderEquiv (baseWeight : σ → ℕ)
    (hpositive : ∀ i, 0 < baseWeight i) (m N : ℕ) (hlarge : N ≤ m) :
    JetWeightPiece R baseWeight m N ≃ₗ[R] JetWeightPiece R baseWeight (m + 1) N :=

theorem presentedJetWeightImageOrderHom_injective_of_homogeneousRelations
    (baseWeight : σ → ℕ) (m N : ℕ) (hlarge : N ≤ m)
    (S : Set (MvPolynomial σ R)) (hS : IsWeightedHomogeneousRelationSet baseWeight S) :
    Function.Injective
      (presentedJetWeightImageOrderHom (R := R) baseWeight m N (Ideal.span S)) := by

theorem finiteDimensional_presentedJetLiBidegreeImage
    (baseWeight : σ → ℕ) (hpositive : ∀ i, 0 < baseWeight i)
    (m p N : ℕ) (I : Ideal (MvPolynomial σ K)) :
    FiniteDimensional K (PresentedJetLiBidegreeImage
      (R := K) baseWeight m p N I) := by

def jetWeight (baseWeight : σ → ℕ) (m : ℕ) : JetVar σ m → ℕ :=
  fun v => baseWeight v.generator + v.order

def jetLiBidegree (baseWeight : σ → ℕ) (m : ℕ) : JetVar σ m → ℕ × ℕ :=
  fun v => (v.order, baseWeight v.generator + v.order)
\end{leancode}

\begin{dictionary}
\lean{jetWeight} gives $x_i^{(r)}$ weight $w_i+r$; \lean{jetLiBidegree} gives it bidegree $(r,w_i+r)$ --- the Li
degree is the number of derivatives. \lean{IsWeightedHomogeneous w D} is Mathlib's ``every monomial has
$w$-weight $D$''. \lean{IsWeightedHomogeneousRelationSet baseWeight S} says every $f\in S$ is homogeneous of
some degree, and \lean{(…).IsHomogeneous (…)} says the ideal is graded. \lean{jetIdealOfSet m S} is the ideal
generated by the $f^{(j)}$, $f\in S$, $j\le m$. \lean{jetWeightPieceOrderEquiv} is a linear isomorphism
(\lean{≃ₗ[R]}) between the weight-$N$ pieces at orders $m$ and $m+1$ when $N\le m$ (a monomial of weight $N$
cannot contain $x^{(r)}$ with $r\ge N$ because all $w_i\ge1$). The presented version is stated as injectivity
of the order-raising map on weight images; surjectivity is \lean{presentedJetWeightImageOrderHom_}\allowbreak\lean{surjective}.
\lean{PresentedJetLiBidegreeImage} is the image in the presented algebra of the bihomogeneous polynomials of
bidegree $(p,N)$.
\end{dictionary}

\textbf{Proof on paper.} (1) If $f$ has weight $D$ then $f(x(t))$, with $t$ given weight $-1$, is homogeneous of
weight $D$, so the $t^n$ coefficient has weight $D+n$. (2) By the Leibniz (Cauchy) formula
$(af)^{(n)} = \sum_{i+j=n}a^{(i)}f^{(j)}$, the coefficients of multiples of $f$ lie in the ideal generated by those
of $f$. (3) At weight $N\le m$ no variable of order $>N$ can occur, and the new top relations $f^{(m+1)}$ of
order $m+1$ have weight $>m\ge N$, so they cannot create a relation in weight $N$ of a homogeneous ideal.
(4) Finitely many monomials have a given positive weight. \qed

\begin{status}
Proved, unconditional; \texttt{\#print axioms} of
\lean{VOA.}\allowbreak\lean{FiniteJet.}\allowbreak\lean{jetIdeal_}\allowbreak\lean{span_}\allowbreak\lean{eq_}\allowbreak\lean{jetIdealOfSet}, \lean{…jetWeightPieceOrderEquiv},
\lean{…presentedJetWeightImageOrderEquiv_}\allowbreak\lean{of_}\allowbreak\lean{homogeneousRelations},
\lean{…finiteDimensional_}\allowbreak\lean{presentedJetLiBidegreeImage}: standard three.
\end{status}

\begin{physics}
Weight-$N$ states of a chiral algebra can carry at most $N$ derivatives, so to describe the classical shadow up
to weight $N$ one needs jets only up to order $N$. This is the reason the infinite arc space can be replaced by
finite jet schemes in every concrete character computation.
\end{physics}

\subsection{The canonical jet map to the Li associated graded}

\begin{definition}[Divided iterates; the jet map]\label{l11:def:jetmap}
For a linear map $D$ of a commutative algebra $A$ put $D^{(r)}(a) = D^r(a)/r!$. For $x:\sigma\to A$, the
\emph{derivation jet evaluation} is the algebra map $\C[x_i^{(r)}]\to A$, $x_i^{(r)}\mapsto D^{(r)}(x_i)$. For
$A = \lelevengr^FV$ with $D$ the descended translation, and an ideal $I$ of relations vanishing on the $x_i$, it
descends to the \emph{Li jet map} $J_m(\C[x]/I)\to\lelevengr^FV$.
\end{definition}
\noindent File \leanpath{VOA/\allowbreak Filtered/\allowbreak ClassicalFreeness.lean}.
\begin{leancode}
noncomputable def dividedIterate (D : A →ₗ[ℂ] A) (r : ℕ) (a : A) : A :=
  ((r.factorial : ℂ)⁻¹) • ((D ^ r) a)

noncomputable def derivationJetEval (D : A →ₗ[ℂ] A) (m : ℕ) (x : σ → A) :
    JetPoly ℂ σ m →ₐ[ℂ] A :=
  MvPolynomial.aeval fun v => dividedIterate D v.order (x v.generator)

noncomputable def liFiniteJetPolynomialMap (W : VertexOperatorAlgebra V) (m : ℕ)
    (x : σ → liAssociatedGraded W) :
    FiniteJet.JetPoly ℂ σ m →ₐ[ℂ] liAssociatedGraded W :=
  FiniteJet.derivationJetEval (liAssociatedGradedTranslation W) m x

noncomputable def liPresentedJetMap (W : VertexOperatorAlgebra V) (m : ℕ)
    (x : σ → liAssociatedGraded W) (I : Ideal (MvPolynomial σ ℂ))
    (hI : I ≤ RingHom.ker (MvPolynomial.aeval x).toRingHom) :
    FiniteJet.PresentedJet (R := ℂ) m I →ₐ[ℂ] liAssociatedGraded W :=
  FiniteJet.presentedJetToDifferentialAlgebra (liAssociatedGradedTranslation W)
    (liAssociatedGradedTranslation_leibniz W) (liAssociatedGradedTranslation_one W) m x I hI
\end{leancode}

\begin{dictionary}
\lean{v.order} and \lean{v.generator} are the two components $r$ and $i$ of a jet variable. The hypothesis
\lean{hI : I ≤ RingHom.}\allowbreak\lean{ker (MvPolynomial.}\allowbreak\lean{aeval x).}\allowbreak\lean{toRingHom} says that the relations in $I$ vanish when
$x_i\mapsto x(i)\in\lelevengr^FV$ (at order zero). Note that \lean{x} is an \emph{arbitrary} family of elements of
the associated graded; the intended choice is the degree-zero classes of strong generators (below). The map is
\emph{canonical}: no basis or normal form is chosen. Neither surjectivity nor injectivity is part of the
definition (the docstring says so).
\end{dictionary}

\begin{physics}
The map sends the classical field $x_i^{(r)}$ to the classical shadow of $\partial^ra^i/r!$. It is the precise form of the slogan that
the classical limit of a chiral algebra is generated by the classical limits of its strong generators and their derivatives.
\end{physics}


\begin{theorem}[The jet map is well defined; its kernel criterion]\label{l11:thm:jetmap}
Let $D$ be a derivation of a commutative $\C$-algebra $A$ with $D(1) = 0$ and $x:\sigma\to A$.
\begin{enumerate}
\item For $f\in\C[x]$ and $n\le m$, the evaluation of $f^{(n)}$ is $D^{(n)}(f(x))$. Hence if $I\subseteq\ker(f\mapsto
  f(x))$, then $J_mI\subseteq\ker$ of the jet evaluation, and the jet evaluation descends to $J_m(\C[x]/I)$.
\item The Li jet map is injective iff the kernel of the polynomial jet evaluation is exactly $J_mI$.
\end{enumerate}
\end{theorem}
\noindent File \leanpath{VOA/\allowbreak Filtered/\allowbreak ClassicalFreeness.lean}.
\begin{leancode}
theorem derivationJetEval_relationCoeff (D : A →ₗ[ℂ] A)
    (hleib : ∀ a b : A, D (a * b) = D a * b + a * D b) (hOne : D 1 = 0)
    (m n : ℕ) (hn : n ≤ m)
    (x : σ → A) (f : MvPolynomial σ ℂ) :
    derivationJetEval D m x (relationCoeff (R := ℂ) m n f) =

theorem derivationJetEval_jetIdeal_le_ker (D : A →ₗ[ℂ] A)
    (hleib : ∀ a b : A, D (a * b) = D a * b + a * D b) (hOne : D 1 = 0)
    (m : ℕ) (x : σ → A)
    (I : Ideal (MvPolynomial σ ℂ))
    (hI : I ≤ RingHom.ker (MvPolynomial.aeval x).toRingHom) :
    jetIdeal (R := ℂ) m I ≤ RingHom.ker (derivationJetEval D m x).toRingHom := by

theorem liPresentedJetMap_injective_iff_polynomial_kernel
    (W : VertexOperatorAlgebra V) (m : ℕ)
    (x : σ → liAssociatedGraded W) (I : Ideal (MvPolynomial σ ℂ))
    (hI : I ≤ RingHom.ker (MvPolynomial.aeval x).toRingHom) :
    Function.Injective (liPresentedJetMap W m x I hI) ↔
      RingHom.ker (liFiniteJetPolynomialMap W m x).toRingHom =
        FiniteJet.jetIdeal (R := ℂ) m I := by
\end{leancode}

\begin{dictionary}
\lean{hleib} and \lean{hOne} say $D$ is a derivation killing $1$ (for the Li graded algebra they are
\lean{liAssociatedGradedTranslation_}\allowbreak\lean{leibniz} and \lean{liAssociatedGradedTranslation_}\allowbreak\lean{one}). The conclusion of (1)
is the higher Leibniz rule $D^{(n)}(ab) = \sum_{i+j=n}D^{(i)}(a)D^{(j)}(b)$ applied to polynomials. (2) is the
first-isomorphism theorem for the composite $\C[x^{(r)}]\to J_m\to\lelevengr^FV$.
\end{dictionary}

\textbf{Proof on paper.} (1) Induction on $f$: both sides are algebra maps in $f$ (by the higher Leibniz rule), and
agree on variables, $x_i(t)\mapsto\sum_rD^{(r)}(x_i)t^r$, and constants ($D^{(r)}c = 0$ for $r\ge1$). (2) The
quotient map is surjective. \qed

\begin{status}
Proved, unconditional; \texttt{\#print axioms} \lean{VOA.liPresentedJetMap} and
\lean{VOA.}\allowbreak\lean{liPresentedJetMap_}\allowbreak\lean{injective_}\allowbreak\lean{iff_}\allowbreak\lean{polynomial_}\allowbreak\lean{kernel}: standard three.
\end{status}

\begin{physics}
The jet map sends the classical field $x^{(r)}$ to the class of $\partial^rx/r!$: it says that the classical
shadow of the chiral algebra is \emph{generated} by the classical shadows of the strong generators and their
derivatives, subject at least to the derivatives of the classical relations. Injectivity (classical freeness) says
there are no further ``quantum'' relations in the classical limit. It holds in some important cases (van Ekeren and Heluani
proved it for the $(2,2k+1)$ Virasoro minimal models) and fails in others (Andrews, van Ekeren and Heluani showed that the Ising
model is not classically free): the classical shadow can have relations that do not come from $\lelevenR$.
\end{physics}

\begin{theorem}[Generator words hit every Li layer]\label{l11:thm:hits}
Let $a^1,\dots,a^g\in V$ be homogeneous of weights $w_i$, and suppose
\begin{equation*}
V = \operatorname{span}\{a^1,\dots,a^g,\lelevenvac\} + C_2(V).
\end{equation*}
Let $I$ be any ideal of relations vanishing on the degree-zero classes $[a^i]\in\lelevengr^F_0V$. Then for every
$p\le m$ and every $N$, every element of the layer $\lelevengr^F_{p,N}V$ is in the image of the Li jet map with
$x_i = [a^i]$. Moreover the jet monomial $\prod x_{i_k}^{(r_k)}$ maps to the class of the state
$a^{i_1}_{(-r_1-1)}\cdots a^{i_s}_{(-r_s-1)}\lelevenvac$, exactly.
\end{theorem}
\noindent File \leanpath{VOA/\allowbreak Filtered/\allowbreak LiGeneratorWords.lean}; generator words
\leanpath{VOA/\allowbreak Modularity/\allowbreak OrderedSpan.lean} (Chapter~\ref{ch:L15}).
\begin{leancode}
def genWord {g : ℕ} (gen : Fin g → 𝒱) (wt : Fin g → ℕ) (l : List (Fin g × ℕ)) :
    List (Letter 𝒱) :=
  l.map (fun q => (gen q.1, wt q.1, q.2))

def pairsLevel {g : ℕ} (l : List (Fin g × ℕ)) : ℕ := (l.map Prod.snd).sum

noncomputable def dividedTranslateState (W : VertexOperatorAlgebra V) (r : ℕ) (a : V) : V :=
  ((r.factorial : ℂ)⁻¹) • ((W.toPreVertexAlgebra.T ^ r) a)

theorem prod_dividedIterate_degreeZeroClass_eq_generatorWord
    (W : VertexOperatorAlgebra V)
    {g : ℕ} (gen : Fin g → V) (wt : Fin g → ℕ)
    (hgen : ∀ i, gen i ∈ W.grading (wt i : ℤ)) (l : List (Fin g × ℕ)) :
    (l.map fun q =>
      FiniteJet.dividedIterate (liAssociatedGradedTranslation W) q.2
        (liDegreeZeroClass W (wt q.1) ⟨gen q.1, hgen q.1⟩)).prod =
      liHomogeneousClass W (pairsLevel l) (wordDeg (genWord gen wt l))
        ⟨wordAct (adjointWeakModule W) (genWord gen wt l) W.vacuum,
          generatorWord_mem_liWeightPiece W gen wt hgen l⟩ := by

theorem liPresentedJetMap_hits_homogeneous_class_of_generators
    (W : VertexOperatorAlgebra V) (m : ℕ)
    (I : Ideal (MvPolynomial (Fin g) ℂ))
    (gen : Fin g → V) (wt : Fin g → ℕ)
    (hgen : ∀ i, gen i ∈ W.grading (wt i : ℤ))
    (hI : I ≤ RingHom.ker
      (MvPolynomial.aeval fun i => liDegreeZeroClass W (wt i) ⟨gen i, hgen i⟩).toRingHom)
    (hdec : (⊤ : Submodule ℂ V) =
      Submodule.span ℂ (Set.range gen ∪ {W.vacuum}) ⊔ O₂ W)
    (p N : ℕ) (hp : p ≤ m) :
    ∀ z : liWeightLayer W p N,
      ∃ y : FiniteJet.PresentedJet (R := ℂ) m I,
\end{leancode}

\begin{dictionary}
A list \lean{l : List (Fin g × ℕ)} of pairs $(i,r)$ encodes the word $a^{i_1}_{(-r_1-1)}\cdots$; \lean{genWord}
turns it into a word of letters, \lean{pairsLevel} is $\sum r_k$ (the Li level), \lean{wordDeg} its weight.
\lean{dividedTranslateState W r a} is $T^ra/r!$. \lean{liHomogeneousClass W p N ⟨z, h⟩} is the class of $z$ in
$\lelevengr^F_{p,N}V$. \lean{hdec} is the displayed spanning hypothesis, with \lean{⊤} the whole space and
\lean{⊔} the sum of subspaces. The conclusion quantifies over \lean{z : liWeightLayer W p N} and asserts a
preimage \lean{y} of its inclusion \lean{DirectSum.}\allowbreak\lean{of … (p, N) z} into $\lelevengr^FV$.

\textbf{Compared with Li's theorem.} Li's theorem (Li 2005; Arakawa 2012) assumes only that
the $[a^i]$ \emph{generate} $\lelevenR$ as an algebra, which can be achieved with finitely many $a^i$ whenever $V$
is finitely strongly generated. The library's hypothesis \lean{hdec} instead asks that the $a^i$ \emph{linearly
span} $V$ modulo $C_2(V)$ and $\lelevenvac$. Since the $a^i$ are finitely many (\lean{Fin g}), \lean{hdec} forces
$\dim\lelevenR\le g+1$: \textbf{the theorem applies only to $C_2$-cofinite $V$}. This is narrower than Li's
statement (moderate, scope). The exact word identity \lean{prod_}\allowbreak\lean{dividedIterate_}\allowbreak\lean{degreeZeroClass_}\allowbreak\lean{eq_}\allowbreak\lean{generatorWord}
has no such hypothesis and holds for every $V$.
\end{dictionary}

\textbf{Proof on paper.} The identity $(T^ra/r!)_{(-1)}b = a_{(-r-1)}b$ shows that the jet monomial evaluates to
the class of the generator word. The hypothesis \lean{hdec} feeds the ordered spanning theorem of
Gaberdiel--Neitzke type (Chapter~\ref{ch:L15}): weight by weight, $F^pV\cap V_N$ is spanned modulo $F^{p+1}V$ by
ordered generator words of level exactly $p$; each of these is hit. \qed

\begin{status}
Proved, unconditional; \texttt{\#print axioms} \lean{VOA.}\allowbreak\lean{liPresentedJetMap_}\allowbreak\lean{hits_}\allowbreak\lean{homogeneous_}\allowbreak\lean{class_}\allowbreak\lean{of_}\allowbreak\lean{generators}:
standard three.
\end{status}

\begin{physics}
``Every classical operator is a polynomial in the generators and their derivatives'': the classical shadow of a
lisse chiral algebra is generated as a differential algebra by the finitely many fields that span $\lelevenR$. This
is the classical counterpart of strong generation (Chapter~\ref{ch:L12}, Theorem~\ref{l12:thm:fingen}).
\end{physics}

\subsection{The Hilbert squeeze}

To prove classical freeness one needs injectivity of the jet map; the library's method is a dimension squeeze,
bidegree by bidegree. The theorem below is used in Theorem~\ref{l11:thm:freeness}: on each bidegree, surjectivity of the jet map
(Theorem~\ref{l11:thm:hits}) plus a dimension bound (from a character identity) forces bijectivity.

\begin{theorem}[Finite-dimensional Hilbert squeeze]\label{l11:thm:squeeze}
\begin{enumerate}
\item A surjective linear map between finite-dimensional spaces with $\dim(\text{source})\le\dim(\text{target})$ is
  bijective.
\item If $c_i\le a_i\le b_i$ for $i$ in a finite set $s$ and $\sum_{i\in s}c_i = \sum_{i\in s}b_i$, then $c_i = a_i = b_i$
  for all $i\in s$.
\end{enumerate}
\end{theorem}
\noindent File \leanpath{VOA/\allowbreak Core/\allowbreak FiniteHilbertSqueeze.lean}.
\begin{leancode}
theorem bijective_of_surjective_of_finrank_le
    (f : A →ₗ[K] B) (hsurj : Function.Surjective f)
    (hle : Module.finrank K A ≤ Module.finrank K B) :
    Function.Bijective f :=

theorem pointwise_eq_of_sandwich_of_sum_eq {ι : Type*} (s : Finset ι)
    (c a b : ι → ℕ)
    (hca : ∀ i ∈ s, c i ≤ a i) (hab : ∀ i ∈ s, a i ≤ b i)
    (hsum : ∑ i ∈ s, c i = ∑ i ∈ s, b i) :
    ∀ i ∈ s, c i = a i ∧ a i = b i := by
\end{leancode}

\begin{dictionary}
Exactly as stated. In application, $s$ is the finite set of Li degrees at a fixed weight $N$, $c_i$ a normal-form
(lower) count, $a_i$ the actual dimension of the Li layer, and $b_i$ the dimension of the jet layer; one character
identity (equality of the \emph{total} sums at weight $N$) then forces equality in every Li degree. No two-variable
character identity is needed.
\end{dictionary}

\textbf{Proof on paper.} (1) Rank--nullity. (2) If some $c_i<b_i$ then $\sum c<\sum b$. \qed

\begin{status}
Proved, unconditional; \texttt{\#print axioms}: standard three (the second even avoids \lean{Classical.choice}).
\end{status}

\begin{physics}
Counting states twice: once in the quantum theory (the character of $V$) and once in the classical theory (the
Hilbert series of the jet scheme). If the classical count is never smaller and the totals agree, there is no room
for extra quantum relations.
\end{physics}

\begin{theorem}[Classical freeness from a bidegree squeeze]\label{l11:thm:freeness}
Let $\sigma$ be finite with positive weights $w_i$, let $x:\sigma\to\lelevengr^FV$ and $I$ be as in
Definition~\ref{l11:def:jetmap}, and suppose
\begin{enumerate}
\item the kernel of the polynomial jet evaluation is bihomogeneous (stable under taking (Li degree, weight)
  components);
\item for every bidegree $(p,N)$, the induced map from the $(p,N)$-image in $J_m(\C[x]/I)$ to
  $\lelevengr^F_{p,N}V$ is surjective;
\item for every bidegree, $\dim(\text{jet }(p,N)\text{-image})\le\dim\lelevengr^F_{p,N}V$.
\end{enumerate}
Then the kernel of the polynomial jet evaluation equals $J_mI$, i.e.\ the Li jet map is injective.
\end{theorem}
\noindent File \leanpath{VOA/\allowbreak Filtered/\allowbreak ClassicalFreeness.lean}.
\begin{leancode}
theorem liFiniteJetPolynomialMap_kernel_eq_of_bidegree_hilbert_squeeze
    [Finite σ]
    (W : VertexOperatorAlgebra V) (m : ℕ)
    (x : σ → liAssociatedGraded W) (I : Ideal (MvPolynomial σ ℂ))
    (hI : I ≤ RingHom.ker (MvPolynomial.aeval x).toRingHom)
    (baseWeight : σ → ℕ) (hpositive : ∀ i, 0 < baseWeight i)
    (hkernelComponent :
      ∀ f ∈ RingHom.ker (liFiniteJetPolynomialMap W m x).toRingHom,
        ∀ d : ℕ × ℕ,
          MvPolynomial.weightedHomogeneousComponent
              (FiniteJet.jetLiBidegree baseWeight m) d f ∈
            RingHom.ker (liFiniteJetPolynomialMap W m x).toRingHom)
    (hsurj : ∀ p N, Function.Surjective
      (liPresentedJetBidegreeMap W m x I hI baseWeight p N))
    (hle : ∀ p N,
      Module.finrank ℂ
          (FiniteJet.PresentedJetLiBidegreeImage baseWeight m p N I) ≤
        Module.finrank ℂ (liWeightLayer W p N)) :
    RingHom.ker (liFiniteJetPolynomialMap W m x).toRingHom =
      FiniteJet.jetIdeal (R := ℂ) m I := by
\end{leancode}

\begin{dictionary}
\lean{hkernelComponent} is (1), with \lean{weightedHomogeneousComponent (jetLiBidegree baseWeight m) d f} the
bidegree-$d$ component of $f$; \lean{hsurj} is (2) and \lean{hle} is (3), with \lean{liPresentedJetBidegreeMap} the
$(p,N)$-component of the jet map restricted to the $(p,N)$-image. All three are \emph{hypotheses}: the theorem is
a reduction of classical freeness to a surjectivity statement (supplied by Theorem~\ref{l11:thm:hits} in the
$C_2$-cofinite case) and a dimension count (a character identity). No VOA in this chapter's files is proved
classically free; the consumers are the boundary affine $\mathfrak{sl}_2$ files of Chapter~\ref{ch:L09}, and classical
freeness itself (boundary affine $\mathfrak{sl}_2$, Virasoro $M(2,2m+3)$) is proved in Chapter~\ref{ch:L20} (Arakawa--Dai).
\end{dictionary}

\textbf{Proof on paper.} By (1) it suffices to show that each bihomogeneous kernel element lies in $J_mI$. On a
fixed bidegree the jet map is surjective (2) with source of dimension $\le$ target (3), hence injective
(Theorem~\ref{l11:thm:squeeze}); so a bihomogeneous polynomial mapping to $0$ is $0$ in $J_m(\C[x]/I)$, i.e.\ lies in
$J_mI$. \qed

\begin{status}
Proved, conditional on the three displayed hypotheses (which are the mathematical content of classical freeness);
\texttt{\#print axioms}: standard three.
\end{status}

\begin{physics}
Classical freeness of a chiral algebra means its classical limit is exactly the arc space of its associated
variety, with no extra constraints. Physically it means that the full operator content of the chiral algebra is
captured by the classical moduli space and its jets --- the Higgs branch ``knows'' the whole Schur sector (a heuristic reading; whether a given chiral algebra of a
four-dimensional theory is classically free is a case-by-case question).
\end{physics}

\subsection{An example: the affine currents span the Li filtration}

\begin{theorem}[Current words and Li's filtration for $V^k(\mathfrak{sl}_2)$]\label{l11:thm:affineli}
Let $k\ne-2$ and let $C_p\subseteq V^k(\mathfrak{sl}_2)$ be the span of the current words
$J^{a_1}_{(-r_1-1)}\cdots J^{a_s}_{(-r_s-1)}\lelevenvac$ with $r_1+\cdots+r_s\ge p$. Then $F^pV^k(\mathfrak{sl}_2) =
C_p$ for all $p$. For any graded ideal $I$, the Li filtration of the quotient $V^k(\mathfrak{sl}_2)/I$ is the image
of that of $V^k(\mathfrak{sl}_2)$, and equals the image of $C_p$.
\end{theorem}
\noindent File \leanpath{VOA/\allowbreak Filtered/\allowbreak AffineCurrentLi.lean}.
\begin{leancode}
def currentWordPiece (k : ℂ) (p : ℕ) : Submodule ℂ (WakModule k) :=
  Submodule.span ℂ
    {x | ∃ w : List (Fin 3 × ℕ), p ≤ currentWordLevel w ∧ wordStateA k w = x}

theorem liFiltration_eq_currentWordPiece (k : ℂ) (hk : k + 2 ≠ 0) (p : ℕ) :
    liFiltration (affineSl2VOA_closed k hk) p = currentWordPiece k p :=

theorem quotientLiFiltration_eq_map (p : ℕ) :
    liFiltration (QuotVOA.quotientVOA W I hI hgr) p =
      (liFiltration W p).map I.mkQ :=

theorem quotientLiFiltration_eq_currentWordPiece (p : ℕ) :
    liFiltration (QuotVOA.quotientVOA (affineSl2VOA_closed k hk) I hI hgr) p =
      quotientCurrentWordPiece k I p :=
\end{leancode}

\begin{dictionary}
\lean{WakModule k} is (despite its name) the universal affine vacuum module $V^k(\mathfrak{sl}_2)$, and
\lean{affineSl2VOA_}\allowbreak\lean{closed k hk} the universal affine VOA (Chapter~\ref{ch:L08}, Definition~\ref{l08:def:Vk});
\lean{wordStateA k w} is the current word $J^{a_1}_{-r_1-1}\cdots\lelevenvac$ attached to the list
$w = [(a_1,r_1),\dots]$ (Definition~\ref{l08:def:words}); \lean{currentWordLevel} is $\sum r_i$. The quotient
statements are in a \lean{section} with variables \lean{W}, \lean{I}, \lean{hI}, \lean{hgr} (a graded ideal of a
general VOA) and, in the last, \lean{k}, \lean{hk} with \lean{W = affineSl2VOA_}\allowbreak\lean{closed k hk}; they are stated
for \emph{every} graded ideal, in particular for the maximal one (the simple affine VOA). The hypothesis $k\ne-2$
(non-critical level) is the standing hypothesis of the construction.
\end{dictionary}

\textbf{Proof on paper.} $C_p\subseteq F^p$ is immediate: a current word is a generator word with generators $J^a$ of
weight one. For $F^p\subseteq C_p$, one shows that the action of any mode of any current word on $C_q$ lands in
$C_{q-N-1}$ (respectively the sharper bound for negative modes), by normal-ordering with the affine commutation
relations; applying this to the letters of an arbitrary Li word gives the reverse inclusion. For quotients, a word
on $V/I$ lifts to a word on $V$ of the same level. \qed

\begin{status}
Proved, unconditional for $k\neq-2$; \texttt{\#print axioms} of \lean{VOA.}\allowbreak\lean{AffineVOA.}\allowbreak\lean{CurrentLi.}\allowbreak\lean{liFiltration_}\allowbreak\lean{eq_}\allowbreak\lean{currentWordPiece}
and \lean{…quotientLiFiltration_}\allowbreak\lean{eq_}\allowbreak\lean{currentWordPiece}: standard three. No PBW basis and no \lean{VOA.pbw} is used.
\end{status}

\begin{physics}
For a Wess--Zumino--Witten current algebra, ``number of derivatives'' can be counted on the currents alone: every
operator is a normal-ordered polynomial in $\partial^rJ^a$, and the Li degree is the total number of
derivatives. The associated variety of $V^k(\mathfrak{sl}_2)$ is $\mathfrak{sl}_2^*$ (all current
configurations), and for the simple quotient at admissible level it is the nilpotent cone (Arakawa) --- the Higgs
branch of Argyres--Douglas theories in the Beem--Rastelli dictionary.
\end{physics}


%======================================================================
\section{Generator-word spanning of Li's filtration}\label{l11:sec:genwords}
%======================================================================

This section documents three files:
\leanpath{VOA/\allowbreak Filtered/\allowbreak GeneratorWordLiSurjectivity.lean},
\leanpath{VOA/\allowbreak Filtered/\allowbreak GeneratorCurrentLiSpanning.lean} and
\leanpath{VOA/\allowbreak Filtered/\allowbreak AffineSl2GeneratorCurrentInstance.lean}, quoted at the commit named in
the chapter's first paragraph. Declarations from other files that the section needs are re-quoted from the same
commit, with their paths.

\paragraph{Notation recalled.} The section is meant to be readable on its own, so we recall the four objects it
works with. Zhu's subspace $C_2(V)$ is the span of all $u_{(-2)}v$, $u,v\in V$, and $\lelevenR = V/C_2(V)$ is
Zhu's $C_2$-algebra, a commutative algebra with product $[u][v] = [u_{(-1)}v]$
(Section~\ref{l11:sec:li}). Li's filtration $V = F^0V\supseteq F^1V\supseteq\cdots$ is defined by letting $F^pV$ be
the span of all words $u^1_{(-n_1-1)}\cdots u^s_{(-n_s-1)}\lelevenvac$ with homogeneous $u^j\in V$ and $n_1+\dots+n_s\ge p$
(Definition~\ref{l11:def:gfil}); the integer $n_j$ counts the derivatives carried by the letter $u^j$, because
$u_{(-n-1)} = (T^nu/n!)_{(-1)}$. Its associated graded is $\lelevengr^FV = \bigoplus_{p,N}\lelevengr^F_{p,N}V$ with
$\lelevengr^F_{p,N}V = (F^pV\cap V_N)/(F^{p+1}V\cap V_N)$, a commutative algebra with a derivation $D$ induced by $T$
(Theorem~\ref{l11:thm:ligr}); its degree-zero part is $\lelevenR$, because $F^1V = C_2(V)$ (Theorem~\ref{l11:thm:f1}).
Finally, \lean{hdec} is the hypothesis of Theorem~\ref{l11:thm:hits}, recalled in the next paragraph; in Lean it is
stated with \lean{O₂ W}, the library's name for $C_2(V)$.

\paragraph{Why a second surjectivity criterion is needed.} Theorem~\ref{l11:thm:hits} proves that the Li jet
map hits every layer of $\lelevengr^FV$, but only under the hypothesis \lean{hdec}: the finitely many states
$a^1,\dots,a^g$ together with $\lelevenvac$ must \emph{linearly span} $V$ modulo $C_2(V)$. As remarked there,
this forces $\dim\lelevenR\le g+1$, so it can only hold for $C_2$-cofinite (lisse) $V$, and even then only
if the generators are chosen to span $\lelevenR$ linearly rather than merely to generate it as an algebra.
For the cases the library actually needs, it fails: for the Virasoro minimal models $M(2,2n+1)$ with the single
generator $\omega$ once $n\ge3$ (Chapter~\ref{ch:L20} writes these models as $M(2,2m+3)$, so $n=m+1$; there
$\lelevenR\cong\C[T]/(T^{m+1})$, here $\C[T]/(T^n)$), and for every affine VOA that is not lisse, whatever finite family is chosen. The
warning box at the end of this section (page~\pageref{l11:warn:hdec}) records the dead end precisely.

The criterion of this section asks something that holds in all these cases: that Li's filtration is
\emph{equal} to the filtration by generator words, i.e.\ that every state of Li degree $\ge p$ is a linear
combination of generator words with at least $p$ derivatives in total. This is a property of \emph{strong
generation} (all of $V$ is spanned by the words), not of $\lelevenR$ alone. The section proceeds in four steps.
\begin{enumerate}
\item The word filtration $\mathcal W_p$ and the property ``$F^pV\subseteq \mathcal W_p$ for all $p$''
  (Definition~\ref{l11:def:genwordpiece}, Proposition~\ref{l11:prop:wordpieceli}).
\item If the property holds, the finite-order bidegree maps and the infinite jet map are surjective, for any
  VOA and any finite family of homogeneous generators (Theorem~\ref{l11:thm:genwordsurj}).
\item How to prove the property: word spanning plus a ``tracked'' estimate on generator modes
  (Theorem~\ref{l11:thm:tracked}); the tracked estimate holds automatically for currents with Lie-type
  commutators (Theorem~\ref{l11:thm:liecurrents}); the property descends to every graded quotient
  (Theorem~\ref{l11:thm:wordquot}).
\item A consistency check: the generic theorem recovers Theorem~\ref{l11:thm:affineli} for
  $V^k(\mathfrak{sl}_2)$ (Theorem~\ref{l11:thm:sl2genwords}), with a worked example in weight two.
\end{enumerate}
These results are used later in the book. The Virasoro minimal models $M(2,2m+3)$ are shown classically free in
Chapter~\ref{ch:L20}, and the jet-surjectivity half of that proof is
\lean{liInfinitePresentedJetMap_}\allowbreak\lean{surjective_}\allowbreak\lean{of_}\allowbreak\lean{genWordSpan}
applied to the Virasoro words. For every Deligne type, the package \lean{CurrentJetSurjectivity} of
Chapter~\ref{ch:L24} (Definition~\ref{l24:def:packages}) is discharged in Theorem~\ref{l24:thm:jetsurj} through
the three generic theorems of Section~\ref{l11:sec:genwords-spanning} below (see also
Definition~\ref{l24:def:genword}).

\subsection{The word filtration}\label{l11:sec:genwords-filtration}

\begin{definition}[Generator words and the word filtration]\label{l11:def:genwordpiece}
Let $V$ be a VOA and $a^1,\dots,a^g\in V$ a finite family of states, with weight labels $w_1,\dots,w_g\in\N$. Encode
a word by a finite list $l = [(i_1,r_1),\dots,(i_s,r_s)]$ of pairs $(i,r)\in\{1,\dots,g\}\times\N$. Its
\emph{state} is
\begin{equation*}
  a_l := a^{i_1}_{(-r_1-1)}a^{i_2}_{(-r_2-1)}\cdots a^{i_s}_{(-r_s-1)}\lelevenvac ,
\end{equation*}
and its \emph{level} is $|l| := r_1+\cdots+r_s$ (the total number of derivatives, since
$a_{(-r-1)}b = (T^ra/r!)_{(-1)}b$). For $p\in\N$ the \emph{word piece} is
\begin{equation*}
  \mathcal W_p := \operatorname{span}_\C\{\,a_l : |l|\ge p\,\}\subseteq V .
\end{equation*}
The family \emph{spans Li's filtration} if $F^pV\subseteq \mathcal W_p$ for every $p\in\N$.
\end{definition}
\noindent\emph{Notation.} The calligraphic $\mathcal W_p$ (``word piece'') is used to keep it apart from Zhu's
$C_2(V)$, which appears in the same arguments: $\mathcal W_2$ is the span of generator words with at least two
derivatives, $C_2(V)$ the span of all $u_{(-2)}v$. For the three $\mathfrak{sl}_2$ currents, $\mathcal W_p$ is the
current-word piece written $C_p$ in Theorem~\ref{l11:thm:affineli} (Lean \lean{currentWordPiece}); in this section the
letter $C$ is reserved for $C_2(V)$.
\noindent File \leanpath{VOA/\allowbreak Filtered/\allowbreak GeneratorWordLiSurjectivity.lean}.
\begin{leancode}
abbrev genWordState (l : List (Fin g × ℕ)) : V :=
  wordAct (adjointWeakModule W) (genWord gen wt l) W.vacuum

def genWordPiece (p : ℕ) : Submodule ℂ V :=
  Submodule.span ℂ {u | ∃ l : List (Fin g × ℕ), p ≤ pairsLevel l ∧ genWordState W gen wt l = u}

def GeneratorWordsSpanLiFiltration : Prop :=
  ∀ p : ℕ, liFiltration W p ≤ genWordPiece W gen wt p
\end{leancode}

\begin{dictionary}
\textbf{Section variables.} The file opens with \lean{variable (W : VertexOperatorAlgebra V) {g : ℕ} (gen : Fin g → V)
(wt : Fin g → ℕ)}, so all three declarations take \lean{W}, \lean{gen}, \lean{wt} as explicit arguments: \lean{gen i} is
$a^{i+1}$ (Lean's \lean{Fin g} is $\{0,\dots,g-1\}$) and \lean{wt i} its weight label.

\textbf{Decoding.} \lean{genWord gen wt l} turns the list of pairs into a list of letters
$(a^{i},w_i,r)$ (Theorem~\ref{l11:thm:hits}); \lean{wordAct (adjointWeakModule W)} applies the letters as modes of the
adjoint module, head of the list outermost, letter $(c,\beta,n)$ acting by $c_{(-n-1)}$ (Definition~\ref{l11:def:gfil}).
Concretely, \lean{genWord} (\leanpath{VOA/\allowbreak Modularity/\allowbreak OrderedSpan.lean}) is \lean{l.map}, sending the
pair $(i,r)$ to the letter $(a^i,w_i,r)$, and \lean{wordAct} is defined by recursion on the list:
$[\,]$ acts as the identity, and $p::t$ acts as the mode of $p$ applied \emph{after} $t$. So the word
$[(i_1,r_1),\dots,(i_s,r_s)]$ is applied to the vacuum from right to left: first $a^{i_s}_{(-r_s-1)}$, last
$a^{i_1}_{(-r_1-1)}$. Note the index shift: the pair records the number $r\ge0$ of derivatives, and the mode applied
is $a^i_{(-r-1)}$ (\lean{mode (A.YM p.1) (-(p.2.2 : ℤ) - 1)} in the source, where \lean{mode f n} is $f_{(n)}$).
So \lean{genWordState W gen wt l} is exactly $a_l$. \lean{pairsLevel l} is $|l| = \sum r_k$. \lean{genWordPiece W gen wt p}
is $\mathcal W_p$, the span (\lean{Submodule.span ℂ}) of the set of word states of level $\ge p$. \lean{liFiltration W p} is Li's
$F^pV$ (Definition~\ref{l11:def:gfil}), and \lean{≤} between submodules is inclusion. So
\lean{GeneratorWordsSpanLiFiltration W gen wt} is literally ``$F^pV\subseteq \mathcal W_p$ for all $p$''.

\textbf{The weight labels play no role here.} \lean{wordAct} reads only the state and the mode index of a letter,
not its weight label. Hence $a_l$, $\mathcal W_p$ and the property do not depend on \lean{wt}, and nothing so far requires the
$a^i$ to be homogeneous. The labels become meaningful only when the hypothesis \lean{hgen} below ties them to the
actual weights.

\textbf{Could it be vacuous?} No. The property is a genuine condition, and it is satisfiable: the empty family
($g=0$) satisfies it exactly when $V=\C\lelevenvac$. It is satisfied by the currents of
$V^k(\mathfrak{sl}_2)$ (Theorem~\ref{l11:thm:sl2genwords}), and the later chapters prove it for the Virasoro minimal
models and the Deligne-series affine VOAs.
\end{dictionary}

\begin{physics}
$\mathcal W_p$ is the space of normal-ordered products of the generating fields carrying at least $p$ derivatives in total,
such as ${:}\partial^{r_1}a^{i_1}\cdots\partial^{r_s}a^{i_s}{:}$ up to factorials. Li's filtration counts derivatives
in an intrinsic way, using \emph{all} states of $V$ as letters. The property says that counting derivatives on
the generators alone gives the same answer: no operator acquires ``hidden'' derivatives by being written through a
composite field. This is what a physicist assumes without comment when writing the classical limit of a chiral
algebra as a differential polynomial algebra in the generating fields.
\end{physics}

\begin{proposition}[Words always lie in Li's filtration; the property is an equality]\label{l11:prop:wordpieceli}
If each $a^i$ is homogeneous of weight $w_i$, then $\mathcal W_p\subseteq F^pV$ for every $p$. Hence, if the family spans Li's
filtration, $F^pV = \mathcal W_p$ for every $p$.
\end{proposition}
\noindent File \leanpath{VOA/\allowbreak Filtered/\allowbreak GeneratorWordLiSurjectivity.lean}.
\begin{leancode}
theorem genWordPiece_le_liFiltration (hgen : ∀ i, gen i ∈ W.grading (wt i : ℤ)) (p : ℕ) :
    genWordPiece W gen wt p ≤ liFiltration W p := by

theorem liFiltration_eq_genWordPiece (hgen : ∀ i, gen i ∈ W.grading (wt i : ℤ))
    (hspan : GeneratorWordsSpanLiFiltration W gen wt) (p : ℕ) :
    liFiltration W p = genWordPiece W gen wt p :=
  le_antisymm (hspan p) (genWordPiece_le_liFiltration W gen wt hgen p)
\end{leancode}

\begin{dictionary}
\lean{hgen} says $a^i\in V_{w_i}$ (\lean{W.grading n} is the weight-$n$ space $V_n$). It is the only hypothesis of
the first statement, and it is not even needed mathematically: $a_l$ is by definition one of the spanning words of
$F^{|l|}V$ (letters from all of $V$, total level $\ge|l|$). The library proves it through the word lemma
\lean{genWord_mem_Gfil} (\leanpath{VOA/\allowbreak Modularity/\allowbreak OrderedSpan.lean}, Chapter~\ref{ch:L15}), which carries the weights. The second statement is
\lean{le_antisymm} of the two inclusions, as the displayed one-line proof shows.
\end{dictionary}

\textbf{Proof on paper.} $F^pV$ is spanned by the words $u^1_{(-n_1-1)}\cdots u^s_{(-n_s-1)}\lelevenvac$ with arbitrary
states $u^j\in V$ and $\sum n_j\ge p$ (Definition~\ref{l11:def:gfil}); taking $u^j = a^{i_j}$ shows $a_l\in F^{|l|}V\subseteq
F^pV$ when $|l|\ge p$. \qed

\begin{status}
Proved, unconditional. \texttt{\#print axioms} of \lean{VOA.}\allowbreak\lean{GeneratorWordsSpanLiFiltration} and
\lean{VOA.}\allowbreak\lean{liFiltration_}\allowbreak\lean{eq_}\allowbreak\lean{genWordPiece}: \lean{[propext, Classical.}\allowbreak\lean{choice,
Quot.}\allowbreak\lean{sound]} (scratch file importing
\texttt{VOA.\allowbreak Filtered.\allowbreak AffineSl2GeneratorCurrentInstance}, run with \texttt{lake env lean -DmaxSynthPendingDepth=3};
commit and procedure in the verification record, Section~\ref{l11:sec:checklist}).
\end{status}

\begin{physics}
Every normal-ordered product of the generators with $p$ derivatives has Li degree at least $p$: derivatives can only
add to the Li degree. The reverse inclusion is the real content (``derivatives cannot hide''), and that is the
hypothesis.
\end{physics}

\subsection{Spanning gives a surjective jet map}\label{l11:sec:genwords-surj}

We recall the two maps whose surjectivity is at stake. The first is the bidegree restriction of the order-$m$ Li jet
map of Definition~\ref{l11:def:jetmap}, already used as hypothesis (2) of Theorem~\ref{l11:thm:freeness}; the second
is the jet map from the \emph{infinite} jet algebra.

\begin{definition}[The bidegree jet map; the infinite jet algebra and the infinite Li jet map]\label{l11:def:infjetmap}
Let $\sigma$ be a set of generators with base weights $w:\sigma\to\N$, let $x:\sigma\to\lelevengr^FV$ and let $I\subseteq\C[x_i]$
be an ideal with $f(x)=0$ for $f\in I$.
\begin{enumerate}
\item For $p,N\in\N$, let $J_m(\C[x]/I)_{(p,N)}$ be the image in $J_m(\C[x]/I)$ of the polynomials of jet bidegree
  $(p,N)$, where $x_i^{(r)}$ has bidegree $(r,w_i+r)$ (Theorem~\ref{l11:thm:jetweights}). The \emph{bidegree jet map}
  is the composite
  \begin{equation*}
    J_m(\C[x]/I)_{(p,N)}\hookrightarrow J_m(\C[x]/I)\xrightarrow{\ \text{Li jet map}\ }\lelevengr^FV
    \xrightarrow{\ \text{component}\ }\lelevengr^F_{p,N}V .
  \end{equation*}
\item The \emph{infinite jet polynomial ring} is $\C[x_i^{(r)} : i\in\sigma,\ r\in\N]$ with no bound on $r$. Its
  \emph{jet ideal} $J_\infty I$ is the ideal generated by all finite jet ideals $J_mI$, $m\in\N$, viewed inside it.
  The \emph{infinite presented jet algebra} is $J_\infty(\C[x]/I) := \C[x_i^{(r)}]/J_\infty I$, and the \emph{infinite
  Li jet map} $J_\infty(\C[x]/I)\to\lelevengr^FV$ is the algebra map $x_i^{(r)}\mapsto D^{(r)}x(i)$, with $D$ the
  translation of $\lelevengr^FV$. Its restriction along $J_m(\C[x]/I)\to J_\infty(\C[x]/I)$ is the order-$m$ Li jet map.
\end{enumerate}
\end{definition}
\noindent File \leanpath{VOA/\allowbreak Filtered/\allowbreak ClassicalFreeness.lean}.
\begin{leancode}
noncomputable def liPresentedJetBidegreeMap
    (W : VertexOperatorAlgebra V) (m : ℕ)
    (x : σ → liAssociatedGraded W) (I : Ideal (MvPolynomial σ ℂ))
    (hI : I ≤ RingHom.ker (MvPolynomial.aeval x).toRingHom)
    (baseWeight : σ → ℕ) (p N : ℕ) :
    FiniteJet.PresentedJetLiBidegreeImage baseWeight m p N I →ₗ[ℂ]
      liWeightLayer W p N :=
  (DirectSum.component ℂ (ℕ × ℕ)
      (fun d : ℕ × ℕ => liWeightLayer W d.1 d.2) (p, N)).comp
    ((liPresentedJetMap W m x I hI).toLinearMap.comp
      (Submodule.subtype
        (FiniteJet.PresentedJetLiBidegreeImage baseWeight m p N I)))
\end{leancode}
\noindent File \leanpath{VOA/\allowbreak Commutative/\allowbreak InfiniteJet.lean}.
\begin{leancode}
abbrev JetPoly (R : Type uR) [CommSemiring R] (σ : Type uσ) :=
  MvPolynomial (σ × ℕ) R

noncomputable def relationSet (I : Ideal (MvPolynomial σ R)) : Set (JetPoly R σ) :=
  {f | ∃ (m : ℕ) (g : FiniteJet.JetPoly R σ m),
    g ∈ FiniteJet.jetIdeal (R := R) m I ∧ ofFinite (R := R) m g = f}

noncomputable def jetIdeal (I : Ideal (MvPolynomial σ R)) : Ideal (JetPoly R σ) :=
  Ideal.span (relationSet (R := R) I)

abbrev PresentedJet (I : Ideal (MvPolynomial σ R)) :=
  JetPoly R σ ⧸ jetIdeal (R := R) I
\end{leancode}
\noindent File \leanpath{VOA/\allowbreak Filtered/\allowbreak InfiniteClassicalFreeness.lean}.
\begin{leancode}
noncomputable def liInfinitePresentedJetMap (W : VertexOperatorAlgebra V)
    (x : σ → liAssociatedGraded W) (I : Ideal (MvPolynomial σ ℂ))
    (hI : I ≤ RingHom.ker (MvPolynomial.aeval x).toRingHom) :
    InfiniteJet.PresentedJet (R := ℂ) I →ₐ[ℂ] liAssociatedGraded W :=
  InfiniteJet.presentedJetToDifferentialAlgebra (liAssociatedGradedTranslation W)
    (liAssociatedGradedTranslation_leibniz W) (liAssociatedGradedTranslation_one W) x I hI
\end{leancode}

\begin{dictionary}
\textbf{The bidegree map.} \lean{FiniteJet.PresentedJetLiBidegreeImage baseWeight m p N I} is the submodule
$J_m(\C[x]/I)_{(p,N)}$ (Theorem~\ref{l11:thm:jetweights}); \lean{Submodule.subtype} is its inclusion;
\lean{DirectSum.component … (p, N)} is the projection of $\lelevengr^FV = \bigoplus_{(p,N)}\lelevengr^F_{p,N}V$
(\lean{liAssociatedGraded}, Theorem~\ref{l11:thm:ligr}) onto one summand. The map is linear (\lean{→ₗ[ℂ]}), not an algebra map,
since its source is only a subspace. The source docstring notes that no homogeneity of the full jet map is assumed in
the definition.

\textbf{The infinite jet algebra.} These declarations live in namespace \lean{VOA.InfiniteJet}, distinct from
\lean{VOA.FiniteJet} of Definition~\ref{l11:def:jets}. Here \lean{JetPoly R σ} is the polynomial ring on the variables
$(i,r)\in\sigma\times\N$, the pair $(i,r)$ standing for $x_i^{(r)}$. \lean{ofFinite m} is the renaming map
$\C[x_i^{(r)}:r\le m]\to\C[x_i^{(r)}]$, so \lean{relationSet I} is the union over $m$ of the images of the finite jet ideals
$J_mI$, and \lean{jetIdeal I} is the ideal it spans. This is the ideal generated by all Hasse--Schmidt coefficients
$f^{(n)}$, $f\in I$, $n\in\N$, which is the usual arc ideal. \lean{PresentedJet I} is the quotient ring (\lean{⧸}).

\textbf{The infinite Li jet map.} \lean{presentedJetToDifferentialAlgebra D hleib hOne x I hI} is the algebra map
$x_i^{(r)}\mapsto D^{(r)}(x(i))$ descended to the quotient (the infinite analogue of Theorem~\ref{l11:thm:jetmap}). Here
$D$ is the translation of $\lelevengr^FV$ together with its Leibniz rule and $D(1)=0$. As in the finite case, \lean{x} is
arbitrary and neither surjectivity nor injectivity is part of the definition.
\end{dictionary}

\begin{physics}
The infinite jet algebra is the algebra of functions on the arc space of $\operatorname{Spec}(\C[x]/I)$: classical
fields $x_i(t) = \sum_rx_i^{(r)}t^r$ constrained only by the Taylor coefficients of the classical relations. The infinite
Li jet map compares this classical phase space of field configurations with the classical limit of the chiral
algebra. The bidegree map looks at one sector at a time: fixed number of derivatives $p$ and fixed conformal weight $N$.
\end{physics}

\begin{theorem}[Spanning of Li's filtration gives surjective jet maps]\label{l11:thm:genwordsurj}
Let $V$ be a VOA, let $a^1,\dots,a^g$ be homogeneous of weights $w_1,\dots,w_g$, and assume that they span Li's filtration
(Definition~\ref{l11:def:genwordpiece}). Let $x_i := [a^i]\in\lelevengr^F_{0,w_i}V$ be their degree-zero classes, and let
$I\subseteq\C[x_1,\dots,x_g]$ be any ideal with $f([a^1],\dots,[a^g]) = 0$ for all $f\in I$. Then:
\begin{enumerate}
\item for all $p,N$, the layer $\lelevengr^F_{p,N}V$ is spanned by the classes of the words $a_l$ with $|l| = p$ and
  $\sum_k(w_{i_k}+r_k) = N$;
\item for all $p,N$, every element of $\lelevengr^F_{p,N}V$ is the image of an element of $J_p(\C[x]/I)$ under the
  order-$p$ Li jet map;
\item for all $p,N$, the bidegree jet map $J_p(\C[x]/I)_{(p,N)}\to\lelevengr^F_{p,N}V$ (order $m=p$, base weights $w$) is
  surjective;
\item the infinite Li jet map $J_\infty(\C[x]/I)\to\lelevengr^FV$ is surjective.
\end{enumerate}
\end{theorem}
\noindent File \leanpath{VOA/\allowbreak Filtered/\allowbreak GeneratorWordLiSurjectivity.lean}.
\begin{leancode}
def genWordClassSet (p N : ℕ) : Set (liAssociatedGraded W) :=
  {x | ∃ l : List (Fin g × ℕ),
      pairsLevel l = p ∧ wordDeg (genWord gen wt l) = N ∧
      x = liHomogeneousClass W (pairsLevel l) (wordDeg (genWord gen wt l))
        ⟨genWordState W gen wt l, generatorWord_mem_liWeightPiece W gen wt hgen l⟩}

theorem genWordClassSet_spans_layer (hspan : GeneratorWordsSpanLiFiltration W gen wt)
    (p N : ℕ) (z : liWeightLayer W p N) :
    DirectSum.of (fun d : ℕ × ℕ => liWeightLayer W d.1 d.2) (p, N) z ∈
      Submodule.span ℂ (genWordClassSet W gen wt hgen p N) := by

theorem liPresentedJetMap_hits_homogeneous_class_of_genWordSpan
    (hspan : GeneratorWordsSpanLiFiltration W gen wt)
    (I : Ideal (MvPolynomial (Fin g) ℂ))
    (hI : I ≤ RingHom.ker
      (MvPolynomial.aeval fun i => liDegreeZeroClass W (wt i) ⟨gen i, hgen i⟩).toRingHom) :
    ∀ p N (z : liWeightLayer W p N),
      ∃ y : FiniteJet.PresentedJet (R := ℂ) p I,
        liPresentedJetMap W p (fun i => liDegreeZeroClass W (wt i) ⟨gen i, hgen i⟩) I hI y =
          DirectSum.of (fun d : ℕ × ℕ => liWeightLayer W d.1 d.2) (p, N) z := by

theorem liPresentedJetBidegreeMap_surjective_of_genWordSpan
    (hspan : GeneratorWordsSpanLiFiltration W gen wt)
    (I : Ideal (MvPolynomial (Fin g) ℂ))
    (hI : I ≤ RingHom.ker
      (MvPolynomial.aeval fun i => liDegreeZeroClass W (wt i) ⟨gen i, hgen i⟩).toRingHom)
    (p N : ℕ) :
    Function.Surjective
      (liPresentedJetBidegreeMap W p (fun i => liDegreeZeroClass W (wt i) ⟨gen i, hgen i⟩)
        I hI wt p N) :=

theorem liInfinitePresentedJetMap_surjective_of_genWordSpan
    (hspan : GeneratorWordsSpanLiFiltration W gen wt)
    (I : Ideal (MvPolynomial (Fin g) ℂ))
    (hI : I ≤ RingHom.ker
      (MvPolynomial.aeval fun i => liDegreeZeroClass W (wt i) ⟨gen i, hgen i⟩).toRingHom) :
    Function.Surjective
      (liInfinitePresentedJetMap W (fun i => liDegreeZeroClass W (wt i) ⟨gen i, hgen i⟩) I hI) := by
\end{leancode}

\begin{dictionary}
\textbf{Hidden arguments.} Besides \lean{W}, \lean{gen}, \lean{wt}, every declaration here also takes the section variable
\lean{(hgen : ∀ i, gen i ∈ W.grading (wt i : ℤ))}, declared just before \lean{genWordClassSet}: the generators are
homogeneous with the stated weights. Lean includes it because it occurs in the statements, inside
\lean{⟨gen i, hgen i⟩}. So the full hypothesis list of each theorem is: a VOA \lean{W}, homogeneous generators (\lean{hgen}),
spanning of Li's filtration (\lean{hspan}), and a relation ideal \lean{I} vanishing on the generator classes (\lean{hI}). There
is no $C_2$-cofiniteness, no PBW basis, no normal form and no classical-freeness hypothesis.

\textbf{Two hypotheses named \lean{hspan}.} In this theorem \lean{hspan : GeneratorWordsSpanLiFiltration W gen wt} is
\emph{Li spanning}, $F^pV\subseteq\mathcal W_p$ for all $p$. Theorem~\ref{l11:thm:tracked} below has a hypothesis with
the same Lean name but a different, weaker content, \emph{word spanning}: the words $a_l$ span $V$, i.e.\
$\mathcal W_0 = V$. The Lean names are local to each theorem and cannot be changed in a verbatim quotation; in the prose
we always say ``Li spanning'' or ``word spanning''.

\textbf{Decoding.} \lean{liDegreeZeroClass W (wt i) ⟨gen i, hgen i⟩} is $[a^i]\in\lelevengr^F_{0,w_i}V$
(Theorem~\ref{l11:thm:f1}). \lean{hI} says the evaluation $f\mapsto f([a^1],\dots,[a^g])$ kills $I$
(\lean{MvPolynomial.aeval}). \lean{wordDeg (genWord gen wt l)} is the weight of the word, $\sum_k(w_{i_k}+r_k)$, computed
from the labels; \lean{generatorWord_mem_liWeightPiece} certifies that $a_l\in F^{|l|}V\cap V_{\mathrm{wordDeg}}$, so that
\lean{liHomogeneousClass W p N ⟨·,·⟩}, the class in $\lelevengr^F_{p,N}V$ placed in the direct sum, makes sense.
\lean{genWordClassSet W gen wt hgen p N} is therefore the set of classes of words of level exactly $p$ and weight exactly
$N$, and (1) says their span contains every \lean{DirectSum.of … (p, N) z}, i.e.\ every element of the $(p,N)$ summand.

\textbf{The orders.} In (2) and (3) the jet order is the Li degree itself (\lean{PresentedJet … p I}, \lean{liPresentedJetMap W p},
\lean{liPresentedJetBidegreeMap W p … wt p N}). This is the smallest order that can work, since a level-$p$ word needs
the variables $x^{(r)}$ with $r\le p$. Larger orders follow by the order-stabilisation of
Theorem~\ref{l11:thm:jetweights}. In (3) the base weights of the jet bigrading are the labels \lean{wt}, which by
\lean{hgen} are the true weights.

\textbf{Compared with Theorem~\ref{l11:thm:hits}.} The conclusion of (2) is the conclusion of
\lean{liPresentedJetMap_}\allowbreak\lean{hits_}\allowbreak\lean{homogeneous_}\allowbreak\lean{class_}\allowbreak\lean{of_}\allowbreak\lean{generators}
at $m=p$, and (3) is the conclusion of the \lean{hdec}-based
\lean{liPresentedJetBidegreeMap_}\allowbreak\lean{surjective_}\allowbreak\lean{of_}\allowbreak\lean{generators}
(\leanpath{VOA/\allowbreak Filtered/\allowbreak LiJetBidegree.lean}). Only the hypothesis changes: \lean{hdec} (span of
$V$ modulo $C_2(V)$, a statement about $\lelevenR$) is replaced by Li spanning (span of each $F^pV$ by generator words, a
statement about all of $V$). How the two hypotheses compare is the subject of Remark~\ref{l11:rem:hdechspan} below: on
paper, Li spanning is equivalent to the $[a^i]$ generating $\lelevenR$ as an \emph{algebra}, while \lean{hdec} asks
them (with $[\lelevenvac]$) to span it as a \emph{vector space}; so \lean{hdec} implies Li spanning. That comparison is
not formalized. In each application below Li spanning is proved, not assumed.

\textbf{Same as the textbook?} (4) is the surjectivity half of the Li--Arakawa comparison
$J_\infty\lelevenR\twoheadrightarrow\lelevengr^FV$, for a presentation $\C[x]/I$ whose relations hold in $\lelevenR$. When
$I$ is exactly the kernel of $\C[x]\to\lelevenR$, the source is $J_\infty\lelevenR$. The theorem is stated for every
smaller $I$ as well; that is harmless, because a smaller $I$ only makes the source bigger.
\end{dictionary}

\textbf{Proof on paper.} (1) Take $z\in F^pV\cap V_N$ representing a class in $\lelevengr^F_{p,N}V$. By Li spanning (\lean{hspan}),
$z = \sum_lc_la_l$ with $|l|\ge p$. Apply the projection $P_N$ onto $V_N$. Each $a_l$ is homogeneous, so
$P_Na_l$ is $a_l$ if the word has weight $N$ and $0$ otherwise. Since $P_Nz = z$, we get
$z = \sum_{\mathrm{wt}(l)=N}c_la_l$. The words with $|l|>p$ lie in $F^{p+1}V$ (Proposition~\ref{l11:prop:wordpieceli}) and
vanish in the quotient $\lelevengr^F_{p,N}V$. What remains is a combination of classes of words with $|l| = p$ and
weight $N$. (2) By Theorem~\ref{l11:thm:hits} (the exact word identity, which needs no \lean{hdec}), the jet monomial
$x_{i_1}^{(r_1)}\cdots x_{i_s}^{(r_s)}$ maps to the class of $a_l$, and it lives in order $\max r_k\le|l| = p$. So every
spanning class of (1) is hit, and the image of a linear map is a subspace. (3) A preimage from (2) can be replaced by
its $(p,N)$ bihomogeneous component, which still maps to the same element of the $(p,N)$ summand
(\lean{liPresentedJetBidegreeMap_}\allowbreak\lean{surjective_}\allowbreak\lean{of_}\allowbreak\lean{homogeneous_}\allowbreak\lean{hits},
\leanpath{LiJetBidegree.lean}). (4) $\lelevengr^FV$ is the direct sum of its $(p,N)$ summands. Each summand is hit from a
finite order by (2), and the infinite map restricts to the finite ones along $J_p\to J_\infty$. \qed

\begin{status}
Proved, unconditional (the hypotheses \lean{hgen}, \lean{hspan}, \lean{hI} are the content of the statement, not open
inputs). \texttt{\#print axioms} of \lean{VOA.}\allowbreak\lean{genWordClassSet_}\allowbreak\lean{spans_}\allowbreak\lean{layer},
\lean{VOA.}\allowbreak\lean{liPresentedJetMap_}\allowbreak\lean{hits_}\allowbreak\lean{homogeneous_}\allowbreak\lean{class_}\allowbreak\lean{of_}\allowbreak\lean{genWordSpan},
\lean{VOA.}\allowbreak\lean{liPresentedJetBidegreeMap_}\allowbreak\lean{surjective_}\allowbreak\lean{of_}\allowbreak\lean{genWordSpan},
\lean{VOA.}\allowbreak\lean{liInfinitePresentedJetMap_}\allowbreak\lean{surjective_}\allowbreak\lean{of_}\allowbreak\lean{genWordSpan}:
\lean{[propext, Classical.}\allowbreak\lean{choice, Quot.}\allowbreak\lean{sound]}. No \lean{VOA.pbw}.
\end{status}

\begin{physics}
Every operator in the classical limit of the chiral algebra is a differential polynomial in the classical generating
fields. Whatever relations the classical fields obey, there are no ``extra'' classical operators that cannot be built from
the generators and their derivatives. This is the surjective half of the comparison between the chiral algebra and the
arc space of its associated variety (the Higgs branch, in the four-dimensional dictionary). (A remark on proof difficulty, not on
physics: this surjectivity half is the easy one; the other half, injectivity, i.e.\ classical freeness, is a genuinely
different statement, which the library approaches through a dimension count, Theorem~\ref{l11:thm:freeness}.)
\end{physics}

\begin{remark}[Li spanning versus \texttt{hdec}: an argument on paper, not formalized]\label{l11:rem:hdechspan}
Let $V$ be a VOA and $a^1,\dots,a^g$ homogeneous of weights $w_i$. Write $x_i = [a^i]\in\lelevenR = \lelevengr^F_0V$.
Then:
\begin{enumerate}
\item the family spans Li's filtration ($F^pV\subseteq\mathcal W_p$ for all $p$) if and only if $x_1,\dots,x_g$ generate
  $\lelevenR$ as a commutative $\C$-algebra;
\item consequently \lean{hdec} implies Li spanning, since \lean{hdec} says that $1,x_1,\dots,x_g$ span $\lelevenR$
  linearly, which is more than generating it as an algebra.
\end{enumerate}
Neither direction uses $C_2$-cofiniteness, a PBW basis, or finite-dimensionality of $\lelevenR$. \textbf{None of this is
proved in the library}: the theorems of this section take Li spanning as a hypothesis, and the applications prove it
directly (Theorems~\ref{l11:thm:tracked}--\ref{l11:thm:wordquot} and Chapter~\ref{ch:L20}), never through (1). The
equivalence is due to Li (\emph{Abelianizing vertex algebras}, 2005) and is used in this form by Arakawa (\emph{A remark
on the $C_2$-cofiniteness condition on vertex algebras}, 2012); the argument below is a self-contained version of it.
\end{remark}

\textbf{The argument on paper.} Three facts about $\lelevengr^FV$ are used (Section~\ref{l11:sec:li}).
(a) The class of $u\in F^pV$ in $\lelevengr^F_pV$ is written $[u]_p$; the product is
$[u]_p[v]_q = [u_{(-1)}v]_{p+q}$ and the derivation is $D[u]_p = [Tu]_{p+1}$, so $D^{(n)} := D^n/n!$ sends $[u]_0$ to
$[T^{(n)}u]_n$ with $T^{(n)} = T^n/n!$. (b) For all $u,v$ and $n\ge0$ one has the \emph{exact} identity
$u_{(-n-1)}v = (T^{(n)}u)_{(-1)}v$. Hence the class of a Li word of level exactly $p$ is a product of divided derivatives:
\begin{equation*}
  \bigl[u^1_{(-n_1-1)}\cdots u^s_{(-n_s-1)}\lelevenvac\bigr]_p = D^{(n_1)}[u^1]_0\cdots D^{(n_s)}[u^s]_0
  \qquad(n_1+\dots+n_s = p),
\end{equation*}
and for generator words this is the exact word identity of Theorem~\ref{l11:thm:hits}:
$[a_l]_{|l|} = \prod_kD^{(r_k)}x_{i_k}$. (c) A Li word of weight $N$ has level at most $N$, because its letters have
weight $\ge0$ and the letter $u_{(-n-1)}$ raises the weight by $\mathrm{wt}\,u+n\ge n$; so $F^pV\cap V_N = 0$ for $p>N$.

\emph{Li spanning $\Rightarrow$ algebra generation.} Take $p=0$: every $u\in V$ is a combination of words $a_l$. Words with
$|l|\ge1$ lie in $\mathcal W_1\subseteq F^1V = C_2(V)$ (Proposition~\ref{l11:prop:wordpieceli} and
Theorem~\ref{l11:thm:f1}), so they vanish in $\lelevenR$. A word of level $0$ is
$a^{i_1}_{(-1)}\cdots a^{i_s}_{(-1)}\lelevenvac$, whose class is the monomial $x_{i_1}\cdots x_{i_s}$ (empty word: $1$). So
$\lelevenR$ is spanned by monomials in the $x_i$. (Only the case $p=0$ of Li spanning, i.e.\ word spanning, was used.)

\emph{Algebra generation $\Rightarrow$ Li spanning.} First, for each $(p,N)$, $\lelevengr^F_{p,N}V$ is spanned by the
classes of generator words of level exactly $p$ and weight $N$. Indeed, it is spanned by classes of Li words of level
exactly $p$ (those of level $>p$ lie in $F^{p+1}V$), and by (b) each is a product $\prod_jD^{(n_j)}[u^j]_0$. By hypothesis,
and since $\lelevengr^F_0V = \lelevenR$ (Theorem~\ref{l11:thm:f1}), each $[u^j]_0$ is a polynomial in the $x_i$. By the Leibniz rule $D^{(n)}(fg) = \sum_{a+b=n}D^{(a)}f\,D^{(b)}g$, the element
$D^{(n_j)}$ of a monomial in the $x_i$ is a combination of products $\prod_kD^{(r_k)}x_{i_k}$ with $\sum_kr_k = n_j$. So
$\lelevengr^F_pV$ is spanned by products $\prod_kD^{(r_k)}x_{i_k}$ with $\sum r_k = p$, which by (b) are the classes
$[a_l]_p$ with $|l| = p$; taking the weight-$N$ component keeps the words of weight $N$, since each $a_l$ is
homogeneous. Now let $z\in F^pV\cap V_N$. By the first step there is a combination $y_p$ of words of level $p$ with
$z-y_p\in F^{p+1}V\cap V_N$. Repeat with $z-y_p$ in level $p+1$, and so on. By (c) the process stops after at most $N-p+1$
steps, with remainder $0$, so $z$ is a combination of words of level $\ge p$: $z\in\mathcal W_p$. \qed

\emph{What this explains.} Li spanning is the natural hypothesis: it asks the generators to generate the classical
algebra $\lelevenR$ in the usual sense, whereas \lean{hdec} asks them, together with the unit, to span it as a vector space, which
already fails for the single Virasoro generator $\omega$ once $\lelevenR\cong\C[T]/(T^n)$ with $n\ge3$ (the monomial
$T^2$ is generated by $T$ but not in the span of $1,T$). The argument also shows why the step from $\lelevenR$
to all of $V$ is possible at all: every class in $\lelevengr^FV$ is a differential polynomial in classes of degree zero.

\subsection{How to prove spanning: the tracked estimate}\label{l11:sec:genwords-spanning}

Spanning of Li's filtration is usually proved by moving modes through words. Write
\begin{equation*}
  t(q,m) := \max(q-m-1,\,0)\qquad(q\in\N,\ m\in\Z).
\end{equation*}
For a homogeneous state $a$, Li's estimates (Theorem~\ref{l11:thm:limult}) give $a_{(m)}F^q\subseteq F^{t(q,m)}$ when $a\in F^0$.
For a negative mode $m=-n-1$ this reads $t(q,-n-1) = q+n$, i.e.\ the mode adds $n$ derivatives. For $m\ge0$ it reads
$q-m-1$: a non-negative mode can remove up to $m+1$ derivatives (it contracts with one of them).

\begin{definition}[The tracked estimate]\label{l11:def:tracked}
A family $a^1,\dots,a^g$ satisfies the \emph{tracked estimate} if $a^i_{(m)}\mathcal W_q\subseteq \mathcal W_{t(q,m)}$ for all
$i$, all $m\in\Z$ and all $q\in\N$: the word filtration obeys the generators' half of Li's estimate.
\end{definition}
\noindent File \leanpath{VOA/\allowbreak Filtered/\allowbreak AffineCurrentLi.lean}.
\begin{leancode}
def currentModeTarget (q : ℕ) (m : ℤ) : ℕ := ((q : ℤ) - m - 1).toNat
\end{leancode}
\noindent File \leanpath{VOA/\allowbreak Filtered/\allowbreak GeneratorCurrentLiSpanning.lean}.
\begin{leancode}
def GeneratorModeTracked : Prop :=
  ∀ (i : Fin g) (m : ℤ) (q : ℕ) (v : V), v ∈ genWordPiece W gen wt q →
    W.smul (gen i) m v ∈ genWordPiece W gen wt (currentModeTarget q m)
\end{leancode}

\begin{dictionary}
\lean{Int.toNat} sends negative integers to $0$, so \lean{currentModeTarget q m} is $t(q,m)$. It was introduced for the
$\mathfrak{sl}_2$ currents (hence the name) and is reused here. \lean{W.smul u m v} is the product $u_{(m)}v$, i.e.\ the
$m$-th mode of $Y(u,z)$ applied to $v$ (Chapter~\ref{ch:L01}). The quantifier order is $\forall i,m,q,v$, with no hypothesis
on $m$: the estimate is required for negative and non-negative modes alike. For negative modes it is automatic for
words (prepend a letter), but it must also hold on linear combinations, which is the same thing since $\mathcal W_q$ is a span.
\end{dictionary}

\begin{physics}
Acting with a generating field's mode on a normal-ordered product: a creation mode $a_{(-n-1)}$ adds $n$ derivatives,
and an annihilation mode $a_{(m)}$, $m\ge0$, removes at most $m+1$, because it contracts with one constituent through the
singular part of the OPE. The tracked estimate says the OPE of the generators with themselves never produces
operators with fewer derivatives than this naive count.
\end{physics}

\begin{theorem}[Word spanning plus the tracked estimate give spanning of Li's filtration]\label{l11:thm:tracked}
If the words $a_l$ span $V$ and the family satisfies the tracked estimate, then it spans Li's filtration, i.e.\
$F^pV\subseteq \mathcal W_p$ for every $p$.
\end{theorem}
\noindent File \leanpath{VOA/\allowbreak Filtered/\allowbreak GeneratorCurrentLiSpanning.lean}.
\begin{leancode}
theorem generatorWordsSpanLiFiltration_of_tracked (htrack : GeneratorModeTracked W gen wt)
    (hspan : Submodule.span ℂ (Set.range (genWordState W gen wt)) = ⊤) :
    GeneratorWordsSpanLiFiltration W gen wt := by
\end{leancode}

\begin{dictionary}
\textbf{Do not confuse the two \lean{hspan}.} Here \lean{hspan} is \emph{word spanning}, not the Li spanning that
carries the same Lean name in Theorem~\ref{l11:thm:genwordsurj} (Li spanning is this theorem's \emph{conclusion}).
\lean{Set.range (genWordState W gen wt)} is the set of all word states $a_l$ (all lists $l$), and \lean{hspan} says
they span $V$ (\lean{⊤} is the whole space), i.e.\ $\mathcal W_0 = V$: the family \emph{strongly generates} $V$ in the
linear sense. Word spanning is the case $p=0$ of Li spanning; the theorem says that, together with the tracked
estimate, it gives all $p$. Note that word spanning is about all words, not ordered (PBW-type) words. \lean{htrack} is Definition~\ref{l11:def:tracked}. Here
\lean{W}, \lean{gen}, \lean{wt} are implicit (\lean{variable {W gen wt}} precedes the theorem in the file). No homogeneity
hypothesis on the $a^i$ is needed.
\end{dictionary}

\textbf{Proof on paper.} \emph{Step 1: all modes of all words.} We show by induction on the word $l$ that
$(a_l)_{(N)}\mathcal W_q\subseteq \mathcal W_{t(q,N)}$ for all $N\in\Z$. For $l=[\,]$, $a_l = \lelevenvac$ and $\lelevenvac_{(N)}v = \delta_{N,-1}v$,
with $t(q,-1) = q$. For $l = (i,r)::l'$, $a_l = a^i_{(-r-1)}a_{l'}$, and the Borcherds iterate formula gives
\begin{equation*}
  (a^i_{(-r-1)}a_{l'})_{(N)}v = \sum_{n\ge0}(-1)^n\binom{-r-1}{n}
  \Bigl(a^i_{(-r-1-n)}(a_{l'})_{(N+n)}v - (-1)^{-r-1}(a_{l'})_{(-r-1+N-n)}a^i_{(n)}v\Bigr).
\end{equation*}
The sum is finite because $a^i_{(n)}v = 0$ and $(a_{l'})_{(N+n)}v=0$ for $n\gg0$. Let $v\in\mathcal W_q$. Recall
$t(x,M) = \max(x-M-1,0)$, so $t(x,-k-1) = x+k$ for $k\ge0$.

\emph{First term.} By induction $(a_{l'})_{(N+n)}v\in\mathcal W_{t(q,N+n)}$; the mode $a^i_{(-r-1-n)}$ is negative, so the
tracked estimate adds $r+n$:
\begin{equation*}
  t\bigl(t(q,N+n),\,-r-1-n\bigr) = \max(q-N-n-1,\,0)+r+n \;\ge\; \max(q-N-1,\,r+n) \;\ge\; t(q,N).
\end{equation*}
\emph{Second term.} By the tracked estimate $a^i_{(n)}v\in\mathcal W_{t(q,n)}$, $t(q,n) = \max(q-n-1,0)$; then by induction
$(a_{l'})_{(-r-1+N-n)}$ sends this into $\mathcal W_{x}$ with, using $-(-r-1+N-n)-1 = r-N+n$,
\begin{equation*}
  x = t\bigl(\max(q-n-1,0),\,-r-1+N-n\bigr) = \max\bigl(\max(q-n-1,\,0)+r-N+n,\;0\bigr)
  \;\ge\; \max(q-N-1+r,\,0) \;\ge\; t(q,N),
\end{equation*}
where the first inequality replaces $\max(q-n-1,0)$ by the smaller $q-n-1$, and the second uses $r\ge0$. Since
$\mathcal W_\bullet$ is decreasing, both terms lie in $\mathcal W_{t(q,N)}$.

\emph{Step 2: negative modes of arbitrary states.} By word spanning (\lean{hspan}) every $u\in V$ is a combination of words, so
Step~1 with $N=-n-1$ gives $u_{(-n-1)}\mathcal W_q\subseteq \mathcal W_{q+n}$ for every $u\in V$.

\emph{Step 3: Li's words.} $F^pV$ is spanned by $u^1_{(-n_1-1)}\cdots u^s_{(-n_s-1)}\lelevenvac$ with $\sum n_j\ge p$. Since
$\lelevenvac\in \mathcal W_0$, applying Step~2 letter by letter shows the word lies in $\mathcal W_{\sum n_j}\subseteq \mathcal W_p$. \qed

\begin{status}
Proved, unconditional. \texttt{\#print axioms}
\lean{VOA.}\allowbreak\lean{generatorWordsSpanLiFiltration_}\allowbreak\lean{of_}\allowbreak\lean{tracked}: \lean{[propext, Classical.}\allowbreak\lean{choice,
Quot.}\allowbreak\lean{sound]}.
\end{status}

\begin{physics}
If the generating fields strongly generate the chiral algebra, and their mutual OPEs respect the derivative count, then
the derivative count on generators is the intrinsic one. The Borcherds iterate formula used here is the mode form of
the statement that the OPE of a normal-ordered product ${:}\partial^ra\,B{:}$ with anything is computed from the OPEs of
$a$ and of $B$ (the generalised Wick theorem).
\end{physics}

\begin{theorem}[Lie-type currents satisfy the tracked estimate]\label{l11:thm:liecurrents}
Let $J^i_m := a^i_{(m)}$ and suppose there are scalars $f^{ij}{}_c$ and $\mathrm{cen}(i,j,m)$ with
\begin{equation*}
  J^i_mJ^j_p - J^j_pJ^i_m = \sum_cf^{ij}{}_c\,J^c_{m+p} + \delta_{m+p,0}\,\mathrm{cen}(i,j,m)\,\mathrm{id}\qquad(\text{all }i,j,m,p),
\end{equation*}
and $J^i_m\lelevenvac = 0$ for $m\ge0$. Then the family satisfies the tracked estimate.
\end{theorem}
\noindent File \leanpath{VOA/\allowbreak Filtered/\allowbreak GeneratorCurrentLiSpanning.lean}.
\begin{leancode}
theorem generatorModeTracked_of_current_comm (hJ : ∀ i m v, W.smul (gen i) m v = J i m v)
    (hcomm : ∀ i j m p, J i m * J j p - J j p * J i m =
      (∑ c, f i j c • J c (m + p)) +
        (if m + p = 0 then cen i j m • (1 : Module.End ℂ V) else 0))
    (hvac : ∀ i m, 0 ≤ m → J i m W.vacuum = 0) :
    GeneratorModeTracked W gen wt := by
\end{leancode}

\begin{dictionary}
\textbf{Section variables.} \lean{J : Fin g → ℤ → Module.End ℂ V}, \lean{f : Fin g → Fin g → Fin g → ℂ} and
\lean{cen : Fin g → Fin g → ℤ → ℂ} are declared by \lean{variable} in the section \lean{LieCurrents}. \lean{hJ} says
that \lean{J i m} is the mode $a^i_{(m)}$; it lets the user supply the modes in whatever form they were constructed. In
\lean{hcomm}, \lean{*} on \lean{Module.End ℂ V} is composition, \lean{•} is scalar multiplication and \lean{1} is the
identity. So \lean{hcomm} is the displayed commutator, with the central term depending on $(i,j,m)$ only when
$m+p=0$. \lean{hvac} is the vacuum condition.

\textbf{What is not assumed.} Nothing requires $f$ to be the structure constants of a Lie algebra (no antisymmetry, no
Jacobi identity), nor the central term to be of the form $m\,k\,\kappa(i,j)$, nor the $a^i$ to have weight one. Only the
\emph{shape} matters: the commutator of two generator modes is a combination of generator modes of total index $m+p$, plus
a scalar. This is the shape of the affine Kac--Moody commutator, so the theorem covers the currents of any universal or
quotient affine VOA once those relations are proved in Lean. It does not cover the Virasoro field, whose commutator
$[L_m,L_n] = (m-n)L_{m+n}+\dots$ has $m$-dependent coefficients; the Virasoro analogue is proved separately for the
Virasoro words (Chapter~\ref{ch:L20}). The conclusion is about the generators \lean{gen}; since \lean{hJ} identifies
their modes with \lean{J}, the theorem is a genuine statement about $V$, not about an auxiliary family of operators.
\end{dictionary}

\textbf{Proof on paper.} \emph{Reduction to words.} $\mathcal W_q$ is spanned by the words $a_l$ with $|l|\ge q$, and
$J^i_m$ is linear, so it suffices to show $J^i_ma_l\in\mathcal W_{t(\ell,m)}$ for every word $a_l$ of level $\ell = |l|$.
This is enough because $t(\cdot,m)$ is monotone in its first argument and $\mathcal W_\bullet$ is decreasing: if
$\ell\ge q$ then $\mathcal W_{t(\ell,m)}\subseteq\mathcal W_{t(q,m)}$.

\emph{Induction on $l$.} If $m<0$, $m = -s-1$, prepending the letter gives a word of level
$\ell+s = t(\ell,m)$. If $m\ge0$ and $l = [\,]$, $J^i_m\lelevenvac = 0$. If $m\ge0$ and $l = (j,r)::l'$ with $|l'| = \ell-r$,
so that $a_l = J^j_{-r-1}a_{l'}$, use the hypothesis on commutators with $p=-r-1$; the central term survives only when
$m+p = 0$, i.e.\ $m = r+1$:
\begin{equation*}
  [J^i_m,J^j_{-r-1}] = \sum_cf^{ij}{}_c\,J^c_{m-r-1} + \delta_{m,r+1}\,\mathrm{cen}(i,j,m)\,\mathrm{id}.
\end{equation*}
Applying both sides to $a_{l'}$ gives
\begin{equation*}
  J^i_mJ^j_{-r-1}a_{l'} = J^j_{-r-1}J^i_ma_{l'} + \sum_cf^{ij}{}_cJ^c_{m-r-1}a_{l'} + \delta_{m,r+1}\,\mathrm{cen}(i,j,m)\,a_{l'} .
\end{equation*}
By induction, $J^i_ma_{l'}\in \mathcal W_{t(\ell-r,m)}$, and prepending $J^j_{-r-1}$ adds $r$: the first term lies in
$\mathcal W_{t(\ell-r,m)+r}\subseteq \mathcal W_{t(\ell,m)}$. By induction, $J^c_{m-r-1}a_{l'}\in \mathcal W_{t(\ell-r,m-r-1)} = \mathcal W_{\max(\ell-m,0)}\subseteq
\mathcal W_{t(\ell,m)}$. When $m = r+1$, the central term is a multiple of $a_{l'}\in \mathcal W_{\ell-r}$, and $\ell-r\ge\ell-m-1$. \qed

\begin{status}
Proved, unconditional. \texttt{\#print axioms}
\lean{VOA.}\allowbreak\lean{generatorModeTracked_}\allowbreak\lean{of_}\allowbreak\lean{current_}\allowbreak\lean{comm}: \lean{[propext, Classical.}\allowbreak\lean{choice,
Quot.}\allowbreak\lean{sound]}.
\end{status}

\begin{physics}
For a current algebra, the OPE $J^i(z)J^j(w)\sim k\kappa^{ij}/(z-w)^2 + f^{ij}{}_cJ^c(w)/(z-w)$ contains only the identity and
a single current. Contracting a current with a constituent of a normal-ordered product therefore replaces one current
by one current (or by a number), and never manufactures a composite operator with fewer derivatives. This is the
Wess--Zumino--Witten version of the derivative count.
\end{physics}

\begin{theorem}[Spanning of Li's filtration passes to graded quotients]\label{l11:thm:wordquot}
Let $I\subseteq V$ be a graded ideal and $\pi:V\to V/I$ the quotient map. If $a^1,\dots,a^g$ span Li's filtration of $V$, then
$\pi(a^1),\dots,\pi(a^g)$ span Li's filtration of the quotient VOA $V/I$ (with the same weight labels).
\end{theorem}
\noindent File \leanpath{VOA/\allowbreak Filtered/\allowbreak GeneratorCurrentLiSpanning.lean}.
\begin{leancode}
theorem mkQ_genWordState (l : List (Fin g × ℕ)) :
    I.mkQ (genWordState W gen wt l) =
      genWordState (QuotVOA.quotientVOA W I hI hgr) (fun i => I.mkQ (gen i)) wt l := by

theorem generatorWordsSpanLiFiltration_quotient (h : GeneratorWordsSpanLiFiltration W gen wt) :
    GeneratorWordsSpanLiFiltration (QuotVOA.quotientVOA W I hI hgr)
      (fun i => I.mkQ (gen i)) wt := by
\end{leancode}

\begin{dictionary}
The section \lean{Quotient} declares \lean{(I : Submodule ℂ V) (hI : W.IsIdeal I) (hgr : QuotVOA.IsGradedIdeal W I)}:
a subspace that is a two-sided ideal for all products and is graded. \lean{QuotVOA.quotientVOA W I hI hgr} is the quotient
VOA and \lean{I.mkQ} the quotient map $\pi$ (Theorem~\ref{l11:thm:affineli}). The first statement says $\pi(a_l)$ is the
word $\pi(a)_l$ of the projected generators. The second is the claimed descent. There is no homogeneity hypothesis on
the $a^i$, and $I$ is any graded ideal, in particular the maximal one, so the theorem applies to simple quotients.
\end{dictionary}

\textbf{Proof on paper.} Li's filtration of $V/I$ is the image $\pi(F^pV)$
(\lean{quotientLiFiltration_}\allowbreak\lean{eq_}\allowbreak\lean{map}, Theorem~\ref{l11:thm:affineli}). Hence
$F^p(V/I) = \pi(F^pV)\subseteq\pi(\mathcal W_p)$, and $\pi(\mathcal W_p)$ is the span of the $\pi(a_l) = \pi(a)_l$, $|l|\ge p$, because $\pi$ is a
morphism of vertex algebras. \qed

\begin{status}
Proved, unconditional. \texttt{\#print axioms}
\lean{VOA.}\allowbreak\lean{generatorWordsSpanLiFiltration_}\allowbreak\lean{quotient}: \lean{[propext, Classical.}\allowbreak\lean{choice,
Quot.}\allowbreak\lean{sound]}.
\end{status}

\begin{physics}
Setting null fields to zero cannot create new operators. The simple quotient is generated by the same fields, and a
derivative count valid upstairs remains valid downstairs. (The \emph{relations} change, of course: that is what $I$ in
Theorem~\ref{l11:thm:genwordsurj} is for.)
\end{physics}

\paragraph{The resulting recipe.} For a quotient $L = V/I$ of a VOA $V$ strongly generated by currents, the three theorems
give: word spanning of $V$ and the current commutators (Theorems~\ref{l11:thm:liecurrents} and~\ref{l11:thm:tracked})
$\Rightarrow$ $V$ spans Li's filtration $\Rightarrow$ $L$ spans Li's filtration (Theorem~\ref{l11:thm:wordquot}) $\Rightarrow$
the infinite Li jet map of $L$ is surjective for every presentation $\C[x]/I'$ whose relations hold in $R_L$
(Theorem~\ref{l11:thm:genwordsurj}). This is exactly how Chapter~\ref{ch:L24} obtains \lean{CurrentJetSurjectivity} for the
seven Deligne types (Theorem~\ref{l24:thm:jetsurj}). Chapter~\ref{ch:L20} uses the last arrow for the Virasoro minimal
models, with the Virasoro words in place of current words.

\subsection{The \texorpdfstring{$\mathfrak{sl}_2$}{sl(2)} consistency instance}\label{l11:sec:genwords-sl2}

\begin{theorem}[The generic theorem recovers the $\mathfrak{sl}_2$ current description]\label{l11:thm:sl2genwords}
Let $k\neq-2$ and let $J^a = J^a_{-1}\lelevenvac\in V^k(\mathfrak{sl}_2)_1$, $a\in\{e,f,h\}$, be the three currents, all with weight
label $1$. They span Li's filtration of $V^k(\mathfrak{sl}_2)$, and $F^pV^k(\mathfrak{sl}_2) = \mathcal W_p$ for every $p$.
\end{theorem}
\noindent File \leanpath{VOA/\allowbreak Filtered/\allowbreak AffineSl2GeneratorCurrentInstance.lean}.
\begin{leancode}
theorem sl2_genWordsSpan :
    GeneratorWordsSpanLiFiltration (affineSl2VOA_closed k hk) (currentState k) (fun _ => 1) := by
  apply generatorWordsSpanLiFiltration_of_tracked
  · exact generatorModeTracked_of_current_comm (J k) fabc
      (fun a b m => (m : ℂ) * k * kappa a b) (smul_currentState_eq_J k hk)
      (fun a b m p => by rw [current_comm]; rfl)
      (fun a m hm => J_vacuum_nonneg k a hm)
  · rw [show genWordState (affineSl2VOA_closed k hk) (currentState k) (fun _ => 1) =
        wordStateA k from funext fun w => (wordStateA_eq_generatorWord k hk w).symm]
    exact wordStateA_span k

theorem sl2_liFiltration_eq_genWordPiece (p : ℕ) :
    liFiltration (affineSl2VOA_closed k hk) p =
      genWordPiece (affineSl2VOA_closed k hk) (currentState k) (fun _ => 1) p :=
  liFiltration_eq_genWordPiece _ _ _ (fun a => currentState_mem_grading k hk a)
    (sl2_genWordsSpan k hk) p
\end{leancode}

\begin{dictionary}
The declarations are \lean{VOA.}\allowbreak\lean{AffineVOA.}\allowbreak\lean{CurrentLi.}\allowbreak\lean{sl2_genWordsSpan} and
\lean{…sl2_}\allowbreak\lean{liFiltration_}\allowbreak\lean{eq_}\allowbreak\lean{genWordPiece}, with section variables
\lean{(k : ℂ) (hk : k + 2 ≠ 0)}. \lean{affineSl2VOA_closed k hk} is $V^k(\mathfrak{sl}_2)$ (Chapter~\ref{ch:L08},
Definition~\ref{l08:def:Vk}); \lean{currentState k a} is $J^a_{-1}\lelevenvac$; \lean{fun _ => 1} gives every current the weight
label $1$. The proof is shown because it \emph{is} the content: it feeds the generic
Theorems~\ref{l11:thm:tracked} and~\ref{l11:thm:liecurrents} with the $\mathfrak{sl}_2$ data of Chapter~\ref{ch:L08}. These
data are the structure constants \lean{fabc} and the central term $m\,k\,\kappa(a,b)$
(Definition~\ref{l08:def:sl2}), the commutator \lean{current_comm} (Theorem~\ref{l08:thm:currentcomm}; its
\lean{• LinearMap.id} is the \lean{• 1} of \lean{hcomm} by \lean{rfl}), vacuum annihilation \lean{J_vacuum_nonneg}, and
spanning of $V^k$ by current words \lean{wordStateA_span}. The latter is transported along
\lean{wordStateA_eq_generatorWord}, which identifies the bespoke current words of Definition~\ref{l08:def:words} with
the generic \lean{genWordState}.

\textbf{What is new, and what is not.} Mathematically nothing: the conclusion is
\lean{liFiltration_}\allowbreak\lean{eq_}\allowbreak\lean{currentWordPiece} of Theorem~\ref{l11:thm:affineli}, re-expressed with
the generic word pieces (the source docstring calls it a consistency check). Its value is as a test that the generic
hypotheses are the right ones: they are satisfied, with no extra work, by the one affine VOA for which a bespoke
proof already existed. The instance for $\mathfrak{sl}_3$ (the $A_2$ case of Chapter~\ref{ch:L24}) is
\lean{DeligneAudit.}\allowbreak\lean{A2.}\allowbreak\lean{universal_}\allowbreak\lean{genWordsSpan}, documented there.
\end{dictionary}

\textbf{Proof on paper.} Theorem~\ref{l11:thm:liecurrents} with the affine commutator
$[J^a_m,J^b_n] = \sum_cf^{ab}{}_cJ^c_{m+n}+m\,k\,\kappa(a,b)\,\delta_{m+n,0}$ and $J^a_{m}\lelevenvac = 0$ ($m\ge0$) gives the
tracked estimate. Current words span $V^k(\mathfrak{sl}_2)$, by construction of the vacuum module as a quotient of
$U(\widehat{\mathfrak{sl}}_2)$. Theorem~\ref{l11:thm:tracked} then gives spanning of Li's filtration, and
Proposition~\ref{l11:prop:wordpieceli} gives the equality, since the currents have weight $1$. \qed

\begin{status}
Proved, unconditional for $k\neq-2$. \texttt{\#print axioms} of
\lean{VOA.}\allowbreak\lean{AffineVOA.}\allowbreak\lean{CurrentLi.}\allowbreak\lean{sl2_genWordsSpan} and
\lean{VOA.}\allowbreak\lean{AffineVOA.}\allowbreak\lean{CurrentLi.}\allowbreak\lean{sl2_}\allowbreak\lean{liFiltration_}\allowbreak\lean{eq_}\allowbreak\lean{genWordPiece}:
\lean{[propext, Classical.}\allowbreak\lean{choice, Quot.}\allowbreak\lean{sound]}. No \lean{VOA.pbw}.
\end{status}

\begin{physics}
In the $SU(2)_k$ Wess--Zumino--Witten current algebra, every operator in the vacuum module is a normal-ordered polynomial
in $\partial^rJ^a$, and the Li degree is the total number of derivatives. The generic theorem reproduces this from two
facts only: the current OPE contains nothing but currents and the identity, and the currents generate.
\end{physics}

\begin{example}[Weight two of $V^k(\mathfrak{sl}_2)$]\label{l11:ex:sl2weight2}
Take $k\neq-2$, $V = V^k(\mathfrak{sl}_2)$, generators $J^e,J^f,J^h$ of weight $1$, and $I = 0$ (no relations). A word
$a_l$ with letters $(a_k,r_k)$ has weight $\sum_k(1+r_k)$ and level $\sum_kr_k$. In weight $N=2$ the words are
\begin{equation*}
  J^a_{-2}\lelevenvac\ (\text{level }1),\qquad J^a_{-1}J^b_{-1}\lelevenvac\ (\text{level }0).
\end{equation*}
By the PBW theorem for $V^k(\mathfrak{sl}_2)$, $V_2$ has dimension $3+6 = 9$: a basis is given by the three states
$J^a_{-2}\lelevenvac$ and the six ordered products $J^a_{-1}J^b_{-1}\lelevenvac$, $a\le b$. The PBW statement used is
\lean{affineSl2PBWBasis} (Chapter~\ref{ch:L08}, Theorem~\ref{l08:thm:pbw}): ordered creation monomials form a basis for
every $k\neq-2$, proved through the Wakimoto realisation without \lean{VOA.pbw}; the count $3+6$ is the weight-$2$ case
(ordered monomials of weight $2$ in the $J^a_{-n}$). By Theorem~\ref{l11:thm:sl2genwords}, $F^1V\cap V_2 = \mathcal W_1\cap V_2 =
\operatorname{span}\{J^a_{-2}\lelevenvac\}$, of dimension $3$. This also equals Zhu's $C_2(V)\cap V_2$, as it must since
$F^1V = C_2(V)$ (Theorem~\ref{l11:thm:f1}): in weight $2$ the only states $u_{(-2)}v$ are $J^a_{(-2)}\lelevenvac = J^a_{-2}\lelevenvac$
(and $\lelevenvac_{(-2)}J^a = 0$). Next,
$F^2V\cap V_2 = \mathcal W_2\cap V_2 = 0$, because a word of level $\ge2$ has weight $\ge3$. So
\begin{equation*}
  \dim\lelevengr^F_{0,2}V = 6,\qquad\dim\lelevengr^F_{1,2}V = 3 .
\end{equation*}
On the jet side, $x_a^{(r)}$ has bidegree $(r,1+r)$, so bidegree $(0,2)$ is spanned by the six monomials $x_ax_b$ and $(1,2)$
by the three $x_a^{(1)}$. Theorem~\ref{l11:thm:genwordsurj}(3) says both bidegree maps are onto. Concretely,
$x_ax_b\mapsto[J^a_{-1}J^b_{-1}\lelevenvac]$ and $x_a^{(1)}\mapsto[J^a_{-2}\lelevenvac]$, by the exact word identity of
Theorem~\ref{l11:thm:hits}. Here the dimensions match ($6=6$, $3=3$), so both maps are even bijective in weight $2$. This
bijectivity is \emph{not} a consequence of the theorems of this section, which prove surjectivity only.

Compare the old hypothesis \lean{hdec} for the same generators. In weight $2$,
$\operatorname{span}\{J^e,J^f,J^h,\lelevenvac\}$ contributes nothing (those states have weights $1$ and $0$), so the weight-$2$
part of $\operatorname{span}\{J^a,\lelevenvac\}+C_2(V)$ is $C_2(V)\cap V_2$, of dimension $3$, while $V_2$ has dimension $9$.
So \lean{hdec} fails for the currents of $V^k(\mathfrak{sl}_2)$ already in weight two, while Theorem~\ref{l11:thm:genwordsurj}
applies.
\end{example}

\phantomsection\label{l11:warn:hdec}
\begin{warning}
\textbf{Dead end: the \texttt{hdec} criteria do not apply to the quasi-lisse cases.}

\textbf{The claim.} It is tempting to obtain surjectivity of the Li jet map, for the affine VOAs of the Deligne series and
for the Virasoro minimal models, from the criteria of Theorem~\ref{l11:thm:hits}, i.e.\ the theorems with
hypothesis
\begin{plaincode}
(hdec : (⊤ : Submodule ℂ V) =
  Submodule.span ℂ (Set.range gen ∪ {W.vacuum}) ⊔ O₂ W)
\end{plaincode}
namely \lean{liPresentedJetMap_}\allowbreak\lean{hits_}\allowbreak\lean{homogeneous_}\allowbreak\lean{class_}\allowbreak\lean{of_}\allowbreak\lean{generators}
(\leanpath{VOA/\allowbreak Filtered/\allowbreak LiGeneratorWords.lean}, Theorem~\ref{l11:thm:hits}),
\lean{liPresentedJetBidegreeMap_}\allowbreak\lean{surjective_}\allowbreak\lean{of_}\allowbreak\lean{generators} and the
\lean{…_of_generators…} squeeze lemmas built on it (\leanpath{VOA/\allowbreak Filtered/\allowbreak LiJetBidegree.lean}), This does not work: they are
\emph{unusable} for every Deligne case and for the Virasoro minimal models $M(2,2n+1)$, $n\ge3$ (the $M(2,2m+3)$ of
Chapter~\ref{ch:L20} with $n=m+1$), with generator $\omega$.

\textbf{What precisely is false.} The theorems themselves are correct (they are proved); what is false in these cases is
their \emph{hypothesis} \lean{hdec}. So the theorems are vacuous there, not wrong. \lean{hdec} says that $\lelevenR$ is
linearly spanned by the images of $\lelevenvac,a^1,\dots,a^g$. Three consequences:
\begin{enumerate}
\item \textbf{\lean{hdec} forces $C_2$-cofiniteness.} This is checked in Lean against the library (axioms \lean{[propext, Classical.}\allowbreak\lean{choice,
  Quot.}\allowbreak\lean{sound]}; see the verification record, Section~\ref{l11:sec:checklist}); the check is not part
  of this release
  (the proof: the composite $\operatorname{span}(\mathrm{gen}\cup\{\lelevenvac\})\hookrightarrow V\to V/C_2(V)$ is surjective by
  \lean{hdec} and has finite-dimensional source). Here \lean{C2Cofinite W} is \lean{FiniteDimensional ℂ (A₂ W)},
  $A_2(W) = V/C_2(V)$ (Chapter~\ref{ch:L12}). Combined with \lean{VOA.}\allowbreak\lean{Sl2NonLisse.}\allowbreak\lean{not_C2Cofinite_}\allowbreak\lean{simpleSl2QL}
  (Chapter~\ref{ch:L16}, Theorem~\ref{l16:thm:nonlisse}: $L_{-4/3}(\mathfrak{sl}_2)$ is not lisse, unconditionally), this gives
  a machine-checked counterexample: \emph{no finite family whatsoever} satisfies \lean{hdec} for
  $L_{-4/3}(\mathfrak{sl}_2)$ (the two-line composition was compiled as part of the same check). On paper, the same holds for every quasi-lisse VOA that is not lisse, in particular for the
  simple affine VOAs $L_k(\mathfrak g)$ of the Deligne series at the levels of Chapter~\ref{ch:L24}. Their associated
  variety is a nilpotent orbit closure of positive dimension (Arakawa--Moreau), so $\lelevenR$ is infinite-dimensional.
\item \textbf{With the currents as generators it fails in weight two.} The weight-$2$ part of the right-hand side of
  \lean{hdec} is $C_2(V)\cap V_2$, so \lean{hdec} forces $(\lelevenR)_2 = 0$. But for $L_k(\mathfrak g)$ the weight-$2$ part of
  $\lelevenR$ is a quotient of $\operatorname{Sym}^2\mathfrak g$ by the weight-two null relations, and it surjects onto
  the degree-$2$ functions on the associated variety. At the Deligne levels that variety is the minimal nilpotent orbit
  closure, whose degree-$2$ functions form the non-zero representation $V(2\theta)$; so $(\lelevenR)_2\neq0$. For $V^k(\mathfrak{sl}_2)$ the count is
  Example~\ref{l11:ex:sl2weight2}. (Paper argument; the Lean check in (1) is the machine-checked form.)
\item \textbf{It fails even in lisse cases with too few generators.} For $M(2,2n+1)$ one has
  $\lelevenR\cong\C[T]/(T^{n})$ with $T = [\omega]$ of weight $2$ (Chapter~\ref{ch:L20}, where the same ring is written
  $\C[T]/(T^{m+1})$). For $n\ge3$, $T^2\neq0$ has weight $4$
  and is not in $\operatorname{span}\{1,T\}$, so \lean{hdec} with the single generator $\omega$ fails. Enlarging the family
  to lifts of a linear basis of $\lelevenR$ would satisfy \lean{hdec}, but the jet presentation would then carry many
  redundant variables, and it still cannot reach non-lisse cases.
\end{enumerate}
\textbf{The replacement} is Theorem~\ref{l11:thm:genwordsurj}, whose hypothesis is spanning of Li's filtration by
generator words, proved for currents by Theorems~\ref{l11:thm:tracked}--\ref{l11:thm:wordquot} and for the Virasoro
words in Chapter~\ref{ch:L20}. By Remark~\ref{l11:rem:hdechspan} its hypothesis is, on paper, exactly algebra generation of
$\lelevenR$, which the single generator $\omega$ of $M(2,2n+1)$ does satisfy. The module docstring of
\leanpath{VOA/\allowbreak Filtered/\allowbreak GeneratorWordLiSurjectivity.lean} records the same limitation. The docstrings
of the \lean{hdec}-based theorems in \leanpath{LiGeneratorWords.lean} and \leanpath{LiJetBidegree.lean} do not; this
is a documentation defect of the library.
\end{warning}

%======================================================================
\section{The symbol toolkit: leading terms, filtered words, certified arithmetic}\label{l11:sec:symbols}
%======================================================================

The proofs that certain states are \emph{linearly independent} (a PBW-type statement for affine
$\mathfrak{sl}_3$ Verma modules, in \leanpath{VOA/\allowbreak Examples/\allowbreak AffineSl3VermaBasis.lean}, whose chapter is in
preparation; non-degeneracy of free-field realisations,
Chapter~\ref{ch:L19}) all follow one pattern: map the states to polynomials, identify a leading (or lowest)
homogeneous term --- a \emph{symbol} --- and show that the symbols are independent. This section documents
the generic, VOA-free part of that pattern, which lives in the files of this chapter. The files are short and
the theorems are standard commutative algebra; the point of the fidelity checks is that the statements are
the right ones and carry no hidden hypothesis.

\begin{definition}[Leading and minimum symbols]\label{l11:def:symbol}
Let $w:\sigma\to\N$ (resp.\ $w:\sigma\to\Z$) be a weight on the variables of $\C[x_s : s\in\sigma]$, and write
$|e|_w = \sum_se_sw(s)$ for the weight of the monomial $x^e$.
\begin{enumerate}
\item $Q$ is a \emph{leading symbol of weight $d$} of $P$ if $Q$ is homogeneous of weight $d$ and every monomial
  of $P-Q$ has weight $<d$.
\item $Q$ is a \emph{minimum symbol of weight $d$} of $P$ if $Q$ is homogeneous of weight $d$ and every monomial
  of $P-Q$ has weight $\ge d+1$; here $\mathrm{AtLeast}_w(d)$ is the span of monomials of weight $\ge d$.
\end{enumerate}
\end{definition}
\noindent Files \leanpath{VOA/\allowbreak Core/\allowbreak PolynomialLeadingSymbol.lean},
\leanpath{VOA/\allowbreak Core/\allowbreak PolynomialMinimumWeight.lean},
\leanpath{VOA/\allowbreak Core/\allowbreak PolynomialMinimumSymbol.lean}.
\begin{leancode}
structure Symbol (w : σ → ℕ) (d : ℕ) (P Q : MvPolynomial σ ℂ) : Prop where
  homogeneous : Q.IsWeightedHomogeneous w d
  remainder : P - Q ∈ Below w d

def AtLeast (w : σ → ℤ) (d : ℤ) : Submodule ℂ (MvPolynomial σ ℂ) :=
  Submodule.span ℂ {P | ∃ e : σ →₀ ℕ, d ≤ Finsupp.weight w e ∧ P = monomial e 1}

structure Symbol (w : σ → ℤ) (d : ℤ) (P Q : MvPolynomial σ ℂ) : Prop where
  homogeneous : Q.IsWeightedHomogeneous w d
  remainder : P-Q ∈ AtLeast w (d+1)
\end{leancode}

\begin{dictionary}
\lean{Below w d} is the set of polynomials all of whose monomials with nonzero coefficient have weight $<d$
(\lean{Finsupp.weight w e} is $|e|_w$); \lean{IsWeightedHomogeneous} is Mathlib's homogeneity. The two \lean{Symbol}
structures live in different namespaces (\lean{VOA.PolynomialLeading} with $\N$-weights and
\lean{VOA.PolynomialMinimum} with $\Z$-weights). \lean{AtLeast} is defined as a span of monomials;
\lean{mem_iff} proves that this is the same as a bound on every nonzero coefficient, with no finiteness of
$\sigma$ needed. Signed weights are used so that a variable of weight $-1$ (``creation raises length by one'')
and differentiation by a variable of larger weight can be handled without truncated subtraction. Note that a
symbol may be $0$ (e.g.\ when the declared weight is larger than the true one), which is why the independence
theorems below assume independence of the symbols.
\end{dictionary}

\begin{physics}
The weight of a monomial is a count of constituents (number of creation operators, number of fields) and the symbol is the term with
the extreme count: the leading term in a semiclassical or free-field expansion of a composite operator.
\end{physics}


\begin{theorem}[Symbols multiply; independent symbols give independent polynomials]\label{l11:thm:symbols}
\begin{enumerate}
\item If $Q,V$ are leading symbols of weights $a,b$ of $P,U$, then $QV$ is a leading symbol of weight $a+b$ of $PU$;
  likewise for minimum symbols (\lean{AtLeast} is multiplicative), and $\partial/\partial x_v$ lowers a minimum bound
  by $w(v)$.
\item If $(P_i)_{i\in\iota}$ is any family with leading (or minimum) symbols $(Q_i)$ of weights $d_i$, and the $Q_i$
  are linearly independent, then the $P_i$ are linearly independent. No bound on the $d_i$ and no finiteness of
  $\iota$ or $\sigma$ is needed.
\item Every nonzero polynomial has a nonzero minimum symbol; multiplication by a fixed nonzero polynomial preserves
  linear independence of all monomials.
\end{enumerate}
\end{theorem}
\noindent Files as above, and \leanpath{VOA/\allowbreak Core/\allowbreak PolynomialNonzeroMinimum.lean}.
\begin{leancode}
theorem Symbol.mul {a b : ℕ} {P Q U V : MvPolynomial σ ℂ}
    (h : Symbol w a P Q) (g : Symbol w b U V) : Symbol w (a+b) (P*U) (Q*V) := by

theorem mul_mem {w : σ → ℤ} {d e : ℤ} {P Q : MvPolynomial σ ℂ}
    (hp : P ∈ AtLeast w d) (hq : Q ∈ AtLeast w e) : P*Q ∈ AtLeast w (d+e) := by

theorem pderiv_mem {w : σ → ℤ} (v : σ) {d : ℤ} {P : MvPolynomial σ ℂ}
    (hp : P ∈ AtLeast w d) : pderiv v P ∈ AtLeast w (d-w v) := by

theorem linearIndependent_of_symbols {ι : Type*} (d : ι → ℕ)
    (P Q : ι → MvPolynomial σ ℂ) (h : ∀ i, Symbol w (d i) (P i) (Q i))
    (hq : LinearIndependent ℂ Q) : LinearIndependent ℂ P := by

theorem linearIndependent_of_symbols {ι : Type*} (d : ι → ℤ)
    (P Q : ι → MvPolynomial σ ℂ) (h : ∀ i, Symbol w (d i) (P i) (Q i))
    (hq : LinearIndependent ℂ Q) : LinearIndependent ℂ P := by

theorem exists_nonzero_symbol (w : σ → ℤ) (P : MvPolynomial σ ℂ) (hP : P≠0) :
    ∃ d : ℤ,∃ Q : MvPolynomial σ ℂ,Q≠0 ∧ Symbol w d P Q := by

theorem monomials_mul_independent (Q : MvPolynomial σ ℂ) (hQ : Q≠0) :
    LinearIndependent ℂ (fun e : σ →₀ ℕ => monomial e (1 : ℂ)*Q) := by
\end{leancode}

\begin{dictionary}
All statements are as described. \lean{LinearIndependent ℂ P} is Mathlib's linear independence of a family
(no nontrivial finite linear relation). In (2), the proof only ever looks at a \emph{finite} relation, so the
family may be infinite. \lean{pderiv v} is $\partial/\partial x_v$. The leading-symbol file also proves
\lean{SeriesSymbol}, the coefficientwise extension to power series with polynomial coefficients (\lean{SeriesSymbol.mul}
uses only finite antidiagonal sums).
\end{dictionary}

\textbf{Proof on paper.} (1) Expand $PU = QV + (P-Q)V + P(U-V)$; every monomial of the last two terms has weight
$<a+b$. (2) Given a finite relation $\sum c_iP_i = 0$, let $n$ be the largest (resp.\ smallest) $d_i$ with
$c_i\ne0$, and take the weight-$n$ homogeneous component: it is $\sum_{d_i=n}c_iQ_i = 0$, contradicting
independence of the $Q_i$. (3) Take the lowest weight occurring in the finite support; $\C[x]$ is a domain. \qed

\begin{status}
Proved, unconditional; \texttt{\#print axioms} of both \lean{linearIndependent_}\allowbreak\lean{of_}\allowbreak\lean{symbols} and of
\lean{exists_nonzero_symbol}: standard three.
\end{status}

\begin{physics}
This is the algebraic form of the semiclassical argument ``if the leading classical terms of a set of operators are
independent, so are the operators''. The weight is the filtration degree (length, number of creation operators,
or $R$-charge), and the symbol is the leading term in a free-field realisation.
\end{physics}

\begin{theorem}[Filtered words have product symbols]\label{l11:thm:filteredwords}
Let $F_0\subseteq F_1\subseteq\cdots$ be subspaces of $V$, $(T_a)_{a\in\alpha}$ operators with $T_aF_d\subseteq F_{d+1}$,
$v\in F_0$, and $\varphi:V\to\C[x_s]$ linear with $\varphi(v) = 1$. Suppose each $T_a$ has a linear symbol $f_a$
(homogeneous of weight $-1$ for the weight $w\equiv-1$) in the sense that $\varphi(T_au)-f_a\varphi(u)\in\mathrm{AtLeast}(-d)$
for $u\in F_d$. Then for every word $a_1\cdots a_\ell$, $\varphi(T_{a_1}\cdots T_{a_\ell}v)$ has minimum symbol
$f_{a_1}\cdots f_{a_\ell}$ of weight $-\ell$. Consequently, for a linear order on $\alpha$, if the commutative products
$\prod_af_a^{e_a}$ (over multiplicity functions $e$) are linearly independent, so are the ordered states
$T_{a_1}\cdots T_{a_\ell}v$ ($a_1\le\cdots\le a_\ell$). The same holds on a tensor product $V\otimes W$ with the
total length filtration and the sum of the factor symbols.
\end{theorem}
{\sloppy \noindent Files \leanpath{VOA/\allowbreak Core/\allowbreak OrderedExponentWord.lean},
\leanpath{VOA/\allowbreak Core/\allowbreak FilteredWordSymbols.lean},
\leanpath{VOA/\allowbreak Core/\allowbreak FilteredTensorSymbols.lean}.\par}
\begin{leancode}
def orderedExponentWord (e : α →₀ ℕ) : List α :=
  e.toMultiset.sort (· ≤ ·)

theorem orderedExponentWord_prod {R : Type*} [CommMonoid R] (e : α →₀ ℕ) (f : α → R) :
    ((orderedExponentWord e).map f).prod=e.prod (fun a n => f a ^ n) := by

def filteredWord (T : α → End ℂ V) (v : V) : List α → V
  | [] => v
  | a::l => T a (filteredWord T v l)

theorem filteredOrderedWord_independent [LinearOrder α]
    (F : ℕ → Submodule ℂ V) (T : α → End ℂ V) (v : V)
    (φ : V →ₗ[ℂ] MvPolynomial σ ℂ) (f : α → MvPolynomial σ ℂ)
    (hv : v∈F 0) (hφ : φ v=1)
    (hT : ∀ a d u,u∈F d → T a u∈F (d+1))
    (hf : ∀ a,(f a).IsWeightedHomogeneous (fun _ => (-1 : ℤ)) (-1))
    (he : ∀ a d u,u∈F d → φ (T a u)-f a*φ u∈
      PolynomialMinimum.AtLeast (fun _ => (-1 : ℤ)) (-(d : ℤ)))
    (hind : LinearIndependent ℂ (fun e : α →₀ ℕ => ((orderedExponentWord e).map f).prod)) :
    LinearIndependent ℂ (fun e : α →₀ ℕ => filteredWord T v (orderedExponentWord e)) := by

def tensorLength (F : ℕ → Submodule ℂ V) (G : ℕ → Submodule ℂ W) (d : ℕ) :
    Submodule ℂ (V ⊗[ℂ] W) := Submodule.span ℂ
  {z | ∃ i j v w,i+j≤d ∧ v∈F i ∧ w∈G j ∧ z=v⊗ₜ[ℂ]w}

theorem tensorSum_polynomial_error (F : ℕ → Submodule ℂ V) (G : ℕ → Submodule ℂ W)
    (φ : V →ₗ[ℂ] MvPolynomial σ ℂ) (ψ : W →ₗ[ℂ] MvPolynomial σ ℂ)
    (T : End ℂ V) (U : End ℂ W) (f g : MvPolynomial σ ℂ)
    (hφ : ∀ d v,v∈F d → φ v∈PolynomialMinimum.AtLeast (fun _ => (-1 : ℤ)) (-(d : ℤ)))
    (hψ : ∀ d w,w∈G d → ψ w∈PolynomialMinimum.AtLeast (fun _ => (-1 : ℤ)) (-(d : ℤ)))
    (hT : ∀ d v,v∈F d → φ (T v)-f*φ v∈PolynomialMinimum.AtLeast (fun _ => (-1 : ℤ)) (-(d : ℤ)))
    (hU : ∀ d w,w∈G d → ψ (U w)-g*ψ w∈PolynomialMinimum.AtLeast (fun _ => (-1 : ℤ)) (-(d : ℤ)))
    (d : ℕ) {z : V⊗[ℂ]W} (hz : z∈tensorLength F G d) :
    tensorPolynomialProduct φ ψ ((T.rTensor W+U.lTensor V) z)-
        (f+g)*tensorPolynomialProduct φ ψ z ∈
      PolynomialMinimum.AtLeast (fun _ => (-1 : ℤ)) (-(d : ℤ)) := by
\end{leancode}

\begin{dictionary}
\lean{e : α →₀ ℕ} is a finitely supported multiplicity function; \lean{e.}\allowbreak\lean{toMultiset.}\allowbreak\lean{sort (· ≤ ·)} is the increasing
word with each letter repeated $e_a$ times, and \lean{orderedExponentWord_}\allowbreak\lean{prod} says its commutative evaluation is
$\prod_af(a)^{e_a}$. \lean{filteredWord T v l} applies the operators of \lean{l} with the head outermost. In
\lean{filteredOrderedWord_}\allowbreak\lean{independent} the hypotheses are exactly those of the theorem: \lean{hv}, \lean{hφ}
(normalisation), \lean{hT} (length raises by one), \lean{hf} (linear symbols of weight $-1$), \lean{he} (the
symbol error is length-preserving) and \lean{hind} (independence of the commutative symbols); the conclusion is
independence of the ordered states. \lean{φ} need not be injective. In the tensor file, \lean{tensorLength F G d}
is the span of $v\otimes w$ with $v\in F_i$, $w\in G_j$, $i+j\le d$; \lean{rTensor}, \lean{lTensor} are $T\otimes1$,
$1\otimes U$.
\end{dictionary}

\textbf{Proof on paper.} Induction on the word using Theorem~\ref{l11:thm:symbols}(1): the new error and the old
error multiplied by $f_a$ both have weight $\ge-\ell+1$. Independence then follows from
Theorem~\ref{l11:thm:symbols}(2) with minimum symbols. For tensor products, the symbol of $v\otimes w$ is the
product of the factor symbols and errors add. \qed

\begin{status}
Proved, unconditional; \texttt{\#print axioms} of \lean{VOA.}\allowbreak\lean{filteredOrderedWord_}\allowbreak\lean{independent} and
\lean{VOA.}\allowbreak\lean{tensorSum_}\allowbreak\lean{polynomial_}\allowbreak\lean{error}: standard three. Consumers: the affine $\mathfrak{sl}_3$ files \leanpath{VOA/\allowbreak Examples/\allowbreak AffineSl3VermaBasis.lean} and
\leanpath{VOA/\allowbreak Examples/\allowbreak AffineSl3HighestTensorFiltration.lean} (not documented in this edition).
\end{status}

\begin{physics}
A PBW basis of a Verma module or a vacuum module is a set of ordered monomials in creation operators. In a free-field
realisation each creation operator becomes, to leading order, multiplication by a classical field; the theorem says
that independence of the classical monomials implies independence of the quantum ordered states --- the
mathematical content of ``the free-field realisation is faithful at leading order''.
\end{physics}

\begin{theorem}[Kernel vectors at a simple determinant root]\label{l11:thm:simpleroot}
Let $M(z)$ be a square matrix of polynomials over a commutative domain $R$, and $z_0\in R$ with $\det M(z_0) = 0$ and
$(\det M)'(z_0)\neq0$. Then there are a left kernel vector $r$ and a right kernel vector $c$ of $M(z_0)$ with
$r^{\mathsf T}M'(z_0)c\neq0$.
\end{theorem}
\noindent File \leanpath{VOA/\allowbreak Core/\allowbreak PolynomialMatrixSimpleRoot.lean}.
\begin{leancode}
theorem kernel_pair_of_simple_root [NoZeroDivisors R]
    (M : Matrix ι ι R[X]) (z : R)
    (hz : M.det.eval z = 0) (hd : M.det.derivative.eval z ≠ 0) :
    ∃ r c : ι → R,
      (∀ j, ∑ i, r i * (M i j).eval z = 0) ∧
      (∀ i, ∑ j, (M i j).eval z * c j = 0) ∧
      (∑ i, ∑ j, (r i * c j) * (M i j).derivative.eval z) ≠ 0 := by
\end{leancode}

\begin{dictionary}
\lean{M.det.eval z} is $\det M(z)$, \lean{derivative} is $d/dz$; \lean{(M i j).eval z} is the $(i,j)$ entry at $z$.
The three conjuncts are $r^{\mathsf T}M(z_0) = 0$, $M(z_0)c = 0$, and $\sum_{i,j}r_ic_jM_{ij}'(z_0)\neq0$.
\lean{[NoZeroDivisors R]} is the domain hypothesis. Exactly the statement above.
\end{dictionary}

\textbf{Proof on paper.} Jacobi's formula $(\det M)' = \operatorname{tr}(\operatorname{adj}(M)M')$ shows
$\operatorname{adj}M(z_0)\neq0$; at a root $M\operatorname{adj}M = \operatorname{adj}M\,M = 0$, so rows and columns of
the adjugate are kernel vectors; differentiating $M\operatorname{adj}M = \det M\cdot1$ and multiplying by $\operatorname{adj}M$
gives $\operatorname{adj}M\,M'\operatorname{adj}M = (\det M)'(z_0)\operatorname{adj}M\neq0$. \qed

\begin{status}
Proved, unconditional; standard three. Consumer: \leanpath{VOA/\allowbreak Examples/\allowbreak AffineSl2JointSimpleRoot.lean}
(Chapter~\ref{ch:L09}).
\end{status}

\begin{physics}
A simple zero of a Kac--Kazhdan-type Gram determinant at a special value of a parameter (level, central charge)
signals a single new null vector; the nonzero pairing through the derivative is the first-order statement that the
null vector is ``created'' at exactly that value, which is what deformation and rigidity arguments use.
\end{physics}

\begin{theorem}[Coefficient bookkeeping for jets]\label{l11:thm:jetconv}
\begin{enumerate}
\item Extending sequences $f,g:\N\to R$ by zero to $\Z$, the integer convolution $\sum_{p\in\Z}f(p)g(n-p)$ (a sum with
  finite support) is the $n$-th coefficient of $\bigl(\sum f_kt^k\bigr)\bigl(\sum g_kt^k\bigr)$, and $0$ for $n<0$; in
  creation-mode indexing, $\sum_pf(-p-a)g(p-n)$ is the coefficient of $t^{-n-a}$.
\item A derivation $D$ of a commutative $\C$-algebra $R$ extends coefficientwise to a derivation of $R[[t]]$.
\end{enumerate}
\end{theorem}
\noindent Files \leanpath{VOA/\allowbreak Core/\allowbreak IntegerJetConvolution.lean},
\leanpath{VOA/\allowbreak Core/\allowbreak PowerSeriesCoefficientDerivation.lean}.
\begin{leancode}
theorem creation_convolution (f g : ℕ → R) (a n : ℤ) :
    (∑ᶠ p : ℤ, extend f (-p-a) * extend g (p-n)) =
      extend (fun r => PowerSeries.coeff r (PowerSeries.mk f * PowerSeries.mk g)) (-n-a) := by

def powerSeriesCoefficientDerivation (D : Derivation ℂ R R) :
    Derivation ℂ (PowerSeries R) (PowerSeries R) where

theorem coeff_powerSeriesCoefficientDerivation (D : Derivation ℂ R R)
    (P : PowerSeries R) (n : ℕ) :
    PowerSeries.coeff n (powerSeriesCoefficientDerivation D P) = D (PowerSeries.coeff n P) := by
\end{leancode}

\begin{dictionary}
\lean{extend f p} is $f(p)$ for $p\ge0$ and $0$ otherwise; \lean{∑ᶠ} is a sum over the (finite) support, and the file
proves the support is finite (\lean{convolution_finite}), so the junk value $0$ of an infinite \lean{finsum} never
occurs. \lean{PowerSeries.mk f} is $\sum_kf(k)t^k$. \lean{Derivation ℂ R R} is Mathlib's $\C$-linear derivation.
\end{dictionary}

\textbf{Proof on paper.} (1) Reindex the finite sum. (2) The Leibniz rule holds coefficient by coefficient because
the coefficients of a product are finite antidiagonal sums. \qed

\begin{status}
Proved, unconditional; standard three. Consumers: the affine $\mathfrak{sl}_3$ classical-limit files \leanpath{VOA/\allowbreak Examples/\allowbreak AffineSl3ClassicalCoefficient.lean}
and \leanpath{VOA/\allowbreak Examples/\allowbreak AffineSl3FirstClassicalDerivation.lean} (not documented in this edition).
\end{status}

\begin{physics}
The coefficient of $z^n$ in a normal-ordered product of two creation fields is a finite convolution of their modes;
this lemma is the bookkeeping that identifies the classical limit of a product of fields with the product of jets.
\end{physics}

\begin{definition}[Certified sparse rational polynomials]\label{l11:def:sparse}
A \emph{sparse polynomial} over $\Q$ in variables $\iota$ (linearly ordered) is a list of pairs (sorted monomial,
nonzero rational coefficient), sorted and without repetitions; \lean{normalizeQ} sorts and merges a raw list of
terms; \lean{evalSparse x P} evaluates at $x:\iota\to A$ in any commutative $\Q$-algebra.
\end{definition}
\noindent File \leanpath{VOA/\allowbreak Core/\allowbreak SparseRational.lean}.
\begin{leancode}
abbrev RawPoly (ι : Type*) := List (ℚ × List ι)

abbrev Poly (ι : Type*) := List (Monomial ι × ℚ)

def normalizeQ [LinearOrder ι] : RawPoly ι → Poly ι
  | [] => []
  | (q, w) :: P => insertQ q w (normalizeQ P)

noncomputable def evalSparse {A : Type*} [CommRing A] [Algebra ℚ A]
    (x : ι → A) (P : Poly ι) : A :=
  (P.map fun t => algebraMap ℚ A t.2 * evalWord x t.1).sum

theorem evalSparse_normalizeQ [LinearOrder ι]
    {A : Type*} [CommRing A] [Algebra ℚ A] (x : ι → A) :
    ∀ P : RawPoly ι, evalSparse x (normalizeQ P) = evalRaw x P
  | [] => by simp [normalizeQ, evalSparse, evalRaw]
  | (q, w) :: P => by
      rw [normalizeQ, evalSparse_insertQ, evalSparse_normalizeQ x P]
      simp [evalRaw]

theorem evalRaw_toRaw [LinearOrder ι]
    {A : Type*} [CommRing A] [Algebra ℚ A] (x : ι → A) (P : Poly ι) :
    evalRaw x (toRaw P) = evalSparse x P := by

theorem evalSparse_addQ [LinearOrder ι]
    {A : Type*} [CommRing A] [Algebra ℚ A] (x : ι → A) :
    ∀ P Q : Poly ι,
      evalSparse x (addQ P Q) = evalSparse x P + evalSparse x Q
  | [], Q => by simp [addQ, evalSparse]
  | (w, c) :: P, [] => by simp [addQ, evalSparse]
  | (w, c) :: P, (v, d) :: Q => by
      rw [addQ]
      generalize hcmp : compareWord w v = cmp
      cases cmp with
      | lt =>
          change algebraMap ℚ A c * evalWord x w +
              evalSparse x (addQ P ((v, d) :: Q)) =
            (algebraMap ℚ A c * evalWord x w + evalSparse x P) +
              (algebraMap ℚ A d * evalWord x v + evalSparse x Q)
          rw [evalSparse_addQ x P ((v, d) :: Q)]
          change _ + (evalSparse x P +
            (algebraMap ℚ A d * evalWord x v + evalSparse x Q)) = _
          ring
      | gt =>
          change algebraMap ℚ A d * evalWord x v +
              evalSparse x (addQ ((w, c) :: P) Q) =
            (algebraMap ℚ A c * evalWord x w + evalSparse x P) +
              (algebraMap ℚ A d * evalWord x v + evalSparse x Q)
          rw [evalSparse_addQ x ((w, c) :: P) Q]
          change _ + (algebraMap ℚ A c * evalWord x w + evalSparse x P +
            evalSparse x Q) = _
          ring
      | eq =>
          have hwv : w = v := (compareWord_eq_iff w v).mp hcmp
\end{leancode}

\begin{dictionary}
Lists of pairs are an executable stand-in for Mathlib's \lean{MvPolynomial}, whose constructors use classical
equality and do not reduce in the kernel. The correctness theorems say that normalisation, addition, scaling
(\lean{evalSparse_scaleQ}), sums and substitution (\lean{evalSparse_bindQ}) commute with evaluation in \emph{every}
commutative $\Q$-algebra, so a polynomial identity checked on the sparse side by \lean{decide}/\lean{norm_num} holds
after evaluation. No \lean{native_decide} is used (the insertion sort is written structurally so that the kernel can
reduce it). The file ends with a small \lean{example} checking that $x_1x_0 - x_0x_1$ normalises to $0$.
\end{dictionary}

\textbf{Proof on paper.} Each list operation implements the corresponding polynomial operation; induction on the
lists. \qed

\begin{status}
Proved, unconditional; \texttt{\#print axioms} \lean{VOA.}\allowbreak\lean{SparseRational.}\allowbreak\lean{evalSparse_}\allowbreak\lean{bindQ}: standard three. Consumer: the
certified Adamovi\'c--Milas identity for the triplet algebra $W(2)$ (Chapter~\ref{ch:L18}).
\end{status}

\begin{physics}
Null-vector identities in a chiral algebra (here the triplet algebra's relation that makes it $C_2$-cofinite) are
large but explicit polynomial identities among modes; certified sparse arithmetic is how such a CAS-style
computation is made into a kernel-checked proof.
\end{physics}


%======================================================================
\section{Presenting a vertex algebra by generators and OPEs (De Sole--Kac)}\label{l11:sec:dsk}
%======================================================================

\subsection{The mathematical background}

Physicists specify a chiral algebra by listing strong generators $a^i(z)$ of conformal weights $\Delta_i$ and their
singular OPEs
\begin{equation*}
a^i(z)a^j(w)\sim\sum_{n\ge0}\frac{(a^i_{(n)}a^j)(w)}{(z-w)^{n+1}},
\end{equation*}
where each coefficient $a^i_{(n)}a^j$ is a normally ordered polynomial in the generators and their derivatives
(``non-linear'' when products occur, as for the $W_3$ algebra). De Sole and Kac (2005) made this precise. Write
$[a_\lambda b] = \sum_{n\ge0}\frac{\lambda^n}{n!}a_{(n)}b$ (the $\lambda$-bracket). Let $R$ be the free
$\C[T]$-module on the generators and $\mathcal T(R)$ its tensor algebra.
\begin{itemize}
\item \emph{Definition 3.1 (non-linear conformal algebra).} A $\lambda$-bracket $R\otimes R\to\C[\lambda]\otimes
  \mathcal T(R)$ that is sesquilinear, $[Ta_\lambda b] = -\lambda[a_\lambda b]$, $[a_\lambda Tb] = (\lambda+T)[a_\lambda b]$, and
  satisfies a grading condition $\Delta([a_\lambda b])<\Delta(a)+\Delta(b)$. (Definition~3.3 adds skew-symmetry and the Jacobi
  identity, making a \emph{non-linear Lie conformal algebra}.)
\item \emph{Lemma 3.2.} The normally ordered product $N$ and the bracket extend uniquely to all of $\mathcal T(R)$.
\item \emph{Definition 3.4.} $M(R)$ is spanned by $A\otimes\bigl((b\otimes c - p(b,c)c\otimes b)\otimes D -
  N(\int_{-T}^0[b_\lambda c]\,d\lambda,\ D)\bigr)$ for $b,c\in R$, $A,D\in\mathcal T(R)$; it is a \emph{left} ideal.
\item \emph{Theorem 3.9.} If the bracket is Lie, $U(R) = \mathcal T(R)/M(R)$ carries a unique vertex algebra structure
  with $Y(a,z)$ prescribed for $a\in R$, and ordered monomials form a basis (PBW).
\end{itemize}
The library builds this machinery in stages. We follow it in order, and we are careful to say at each point what is a
definition, what is proved, what is refuted, and what is a contract that nothing fulfils.

\subsection{The tensor algebra, its weights and translation}

\begin{definition}[$\mathcal T(R)$ with divided-power generators; weights]\label{l11:def:TR}
Let $\iota$ index the strong generators. $R$ has the $\C$-basis $T^ka^i/k!$, indexed by pairs $(i,k)\in\iota\times\N$;
$\mathcal T(R)$ is the free associative algebra on these pairs, with its basis of words. With weights $\Delta_i\in\N$,
the letter $(i,k)$ has weight $\Delta_i+k$ and a word the sum of its letters' weights; $\mathrm{filt}(d)$ is the span of words
of weight $\le d$ ($d\in\Z$; it is $0$ for $d<0$). The translation $T$ is the unique derivation of $\mathcal T(R)$ with
$T(T^ka^i/k!) = (k+1)\,T^{k+1}a^i/(k+1)!$.
\end{definition}
\noindent File \leanpath{VOA/\allowbreak Core/\allowbreak NonlinearConformal.lean}.
\begin{leancode}
abbrev GenIdx (ι : Type u) : Type u := ι × ℕ

abbrev TR (ι : Type u) : Type u := FreeAlgebra ℂ (GenIdx ι)

noncomputable def gv (p : GenIdx ι) : TR ι := wbasis ι (FreeMonoid.of p)

def wtIdx : GenIdx ι → ℕ := fun p => wt p.1 + p.2

def wtWord : FreeMonoid (GenIdx ι) → ℕ := fun w => (w.toList.map (wtIdx wt)).sum

noncomputable def filt (d : ℤ) : Submodule ℂ (TR ι) :=
  Submodule.span ℂ (wbasis ι '' {w | ((wtWord wt w : ℤ)) ≤ d})

noncomputable def TGen (p : GenIdx ι) : TR ι := ((p.2 + 1 : ℕ) : ℂ) • gv (p.1, p.2 + 1)

noncomputable def Ttr : TR ι →ₗ[ℂ] TR ι where
  toFun := fun x => (phiT x).snd
  map_add' := fun x y => by rw [map_add, TrivSqZeroExt.snd_add]
  map_smul' := fun c x => by rw [RingHom.id_apply, map_smul, TrivSqZeroExt.snd_smul]

theorem Ttr_mul (x y : TR ι) : Ttr (x * y) = x * Ttr y + Ttr x * y := by

theorem Ttr_gv (p : GenIdx ι) : Ttr (gv p) = ((p.2 + 1 : ℕ) : ℂ) • gv (p.1, p.2 + 1) := by

theorem Ttr_mem_filt (wt : ι → ℕ) {d : ℤ} {x : TR ι} (hx : x ∈ filt wt d) :
    Ttr x ∈ filt wt (d + 1) := by
\end{leancode}

\begin{dictionary}
\lean{FreeAlgebra ℂ X} is the free associative unital $\C$-algebra on the set $X$, i.e.\ the tensor algebra of the vector
space with basis $X$; so \lean{TR ι} is $\mathcal T(R)$ with $R$ free over $\C[T]$ on $\iota$. \lean{gv (i, k)} is the
letter $T^ka^i/k!$ (the library's ``divided power'' normalisation; De Sole--Kac use $T^ka^i$, a harmless rescaling of the
same basis). \lean{wbasis ι} is the word basis (indexed by \lean{FreeMonoid}, i.e.\ lists). The derivation is built by the
dual-numbers trick: \lean{phiT} is the algebra map $x\mapsto x+\varepsilon T(x)$ into $\mathcal T(R)\oplus\varepsilon\mathcal T(R)$
(\lean{TrivSqZeroExt}), and \lean{Ttr} is its $\varepsilon$-component. \lean{Ttr_mul} is the Leibniz rule (written
$x\,Ty + Tx\,y$), \lean{Ttr_gv} the value on letters, and \lean{Ttr_mem_filt} says $T$ raises weight by one.

\textbf{Restriction relative to De Sole--Kac.} $R$ is assumed \emph{free} over $\C[T]$ (no torsion); central elements
such as $c$ are handled as scalar multiples of the unit word $1\in\mathcal T(R)$. This covers every example in the
library and in most of physics.
\end{dictionary}

\begin{physics}
$\mathcal T(R)$ is the space of all formal normal-ordered words ${:}\partial^{k_1}a^{i_1}\cdots\partial^{k_r}a^{i_r}{:}$
before any relation is imposed --- not even commutativity of normal ordering. The word ${:}L\,L{:}$ and its reorderings are
different elements here. $T$ is $\partial_z$ acting on composite operators by the Leibniz rule.
\end{physics}

\begin{definition}[OPE data: De Sole--Kac Definition 3.1, free-generator form]\label{l11:def:opedata}
An \emph{OPE datum} on generators $\iota$ consists of weights $\Delta_i\ge1$, parities $p_i\in\Z/2$, and elements
$a^i_{(n)}a^j\in\mathcal T(R)$ for $n\in\N$ (the singular OPE table) such that
$a^i_{(n)}a^j\in\mathrm{filt}(\Delta_i+\Delta_j-n-1)$. In particular $a^i_{(n)}a^j = 0$ for $n\ge\Delta_i+\Delta_j$.
\end{definition}
\noindent File \leanpath{VOA/\allowbreak Core/\allowbreak NonlinearConformal.lean}.
\begin{leancode}
structure OPEData (ι : Type u) where
  wt : ι → ℕ
  wt_pos : ∀ i, 0 < wt i
  par : ι → ZMod 2
  ope : ι → ι → ℕ → TR ι
  ope_wt : ∀ i j (n : ℕ), ope i j n ∈ filt wt ((wt i : ℤ) + (wt j : ℤ) - (n : ℤ) - 1)

theorem ope_eq_zero (i j : ι) {n : ℕ} (h : D.wt i + D.wt j ≤ n) : D.ope i j n = 0 := by

noncomputable def virasoroOPE (c : ℂ) : OPEData Unit where
  wt := fun _ => 2
  wt_pos := fun _ => by norm_num
  par := fun _ => 0
  ope := fun _ _ n =>
    if n = 0 then gv ((), 1)
    else if n = 1 then (2 : ℂ) • gv ((), 0)
    else if n = 3 then (c / 2) • (1 : TR Unit)
    else 0
  ope_wt := by
    intro i j n
    split_ifs with h0 h1 h3
    · subst h0
      exact gv_mem_filt _ (by norm_num [wtIdx])
    · subst h1
      exact Submodule.smul_mem _ _ (gv_mem_filt _ (by norm_num [wtIdx]))
    · subst h3
      exact Submodule.smul_mem _ _ (one_mem_filt _ (by norm_num))
    · exact Submodule.zero_mem _
\end{leancode}

\begin{dictionary}
\lean{ope i j n} is $a^i_{(n)}a^j$ (an element of \lean{TR ι}); \lean{ope_wt} is the grading condition, with integer
arithmetic so that a negative bound forces $0$ (\lean{ope_eq_zero}). \lean{virasoroOPE c} has one generator
(\lean{Unit}, the one-element type) $L$ of weight $2$ and table $L_{(0)}L = TL$ (the letter \lean{gv ((), 1)}),
$L_{(1)}L = 2L$, $L_{(3)}L = \frac c2\lelevenvac$, all others $0$. Since $[L_\lambda L] = (T+2\lambda)L+\frac c{12}\lambda^3$ and
$\lambda^3 = 3!\,\lambda^{(3)}$, the coefficient $L_{(3)}L = 6\cdot\frac c{12} = \frac c2$ is correct.

\textbf{Fidelity.} (a) Sesquilinearity is not a field: it is \emph{built into} the extension of the table to all of $R$
(Definition~\ref{l11:def:brk}), so it holds by construction. (b) The grading condition is stated with the conformal
weight, $\Delta(a_{(n)}b)\le\Delta a+\Delta b-n-1$; De Sole--Kac's condition is the weaker
$\Delta([a_\lambda b])<\Delta a+\Delta b$ for a filtration in which $T$ has degree $0$. The library proves that its condition
implies theirs (\lean{ope_mem_zfilt}, Section~\ref{l11:sec:dskgrading}). For a VOA, the library's condition is the natural
one. (c) \textbf{Skew-symmetry and the Jacobi identity are not required.} \lean{OPEData} is a non-linear
\emph{conformal} algebra (Definition~3.1), not a non-linear \emph{Lie} conformal algebra (Definition~3.3). Any table
satisfying the weight bound is accepted, including tables that come from no vertex algebra; for those the quotient
constructed below will in general be degenerate. This is faithful to Definition~3.1 and is not a defect, but a reader must
not read ``\lean{OPEData}'' as ``consistent OPE''. (d) The docstring proposes handling half-integer weights by doubling all
weights; then the bound $\Delta_i+\Delta_j-n-1$ is no longer the conformal one (it becomes weaker), so termination still
holds but the bound is not sharp. No library datum uses this.
\end{dictionary}

\begin{physics}
\lean{virasoroOPE c} is the textbook stress-tensor OPE
$T(z)T(w)\sim\frac{c/2}{(z-w)^4}+\frac{2T(w)}{(z-w)^2}+\frac{\partial T(w)}{z-w}$. The grading condition is dimensional
analysis: the coefficient of $(z-w)^{-n-1}$ has dimension $\Delta_i+\Delta_j-n-1$, so the OPE of two fields of finite dimension
has finitely many singular terms. Jacobi (associativity of the OPE) is the crossing-symmetry constraint that a consistent
chiral algebra must satisfy; \lean{OPEData} does not impose it.
\end{physics}

\begin{definition}[The $\lambda$-bracket on all of $R$]\label{l11:def:brk}
For a function $f:\N\to\mathcal T(R)$ (the slices $n\mapsto a_{(n)}b$) put $(L f)(n) = -n f(n-1)$ and
$(R f)(n) = T f(n)+n f(n-1)$, implementing $(Ta)_{(n)}b = -n\,a_{(n-1)}b$ and $a_{(n)}(Tb) = T(a_{(n)}b)+n\,a_{(n-1)}b$. Then
\begin{equation*}
\bigl(T^ka^i/k!\bigr)_{(n)}\bigl(T^la^j/l!\bigr) := \frac1{k!\,l!}\bigl(L^kR^l(a^i_{(\cdot)}a^j)\bigr)(n).
\end{equation*}
\end{definition}
\noindent File \leanpath{VOA/\allowbreak Core/\allowbreak NonlinearConformal.lean}.
\begin{leancode}
noncomputable def Lstep (f : ℕ → TR ι) : ℕ → TR ι := fun n => -(n : ℂ) • f (n - 1)

noncomputable def Rstep (f : ℕ → TR ι) : ℕ → TR ι := fun n => Ttr (f n) + (n : ℂ) • f (n - 1)

noncomputable def brk (D : OPEData ι) (p q : GenIdx ι) (n : ℕ) : TR ι :=
  (((p.2).factorial : ℂ) * ((q.2).factorial : ℂ))⁻¹ •
    Lstep^[p.2] (Rstep^[q.2] (D.ope p.1 q.1)) n
\end{leancode}

\begin{dictionary}
\lean{Lstep^[k]} is the $k$-fold iterate. In natural-number arithmetic \lean{n - 1} is $0$ at $n=0$, but it is multiplied
by $n = 0$, so the convention is harmless. The two steps commute (\lean{Lstep_Rstep}), so the order of iteration does not
matter.
\end{dictionary}

\begin{physics}
Sesquilinearity is the statement that $\partial_z$ of a field acts on the OPE by differentiating: the OPE of $\partial a$ with $b$
is the derivative of the OPE of $a$ with $b$ in the first variable. Once one knows the OPE of the generators, the OPE of all
their derivatives follows.
\end{physics}

\begin{theorem}[Properties of the extended bracket]\label{l11:thm:brk}
For all letters $p = (i,k)$, $q = (j,l)$ and $n\in\N$:
\begin{enumerate}
\item on bottom letters the bracket is the table: $(a^i)_{(n)}(a^j) = $ \lean{D.ope i j n};
\item sesquilinearity: $(k+1)\,(T^{k+1}a^i/(k+1)!)_{(n)}q = -n\,p_{(n-1)}q$ and
  $(l+1)\,p_{(n)}(T^{l+1}a^j/(l+1)!) = T(p_{(n)}q)+n\,p_{(n-1)}q$;
\item grading: $p_{(n)}q\in\mathrm{filt}(\Delta(p)+\Delta(q)-n-1)$, hence $p_{(n)}q = 0$ for $n\ge\Delta(p)+\Delta(q)$.
\end{enumerate}
\end{theorem}
\noindent File \leanpath{VOA/\allowbreak Core/\allowbreak NonlinearConformal.lean}.
\begin{leancode}
theorem brk_bot (D : OPEData ι) (i j : ι) (n : ℕ) : brk D (i, 0) (j, 0) n = D.ope i j n := by

theorem brk_left (D : OPEData ι) (i : ι) (k : ℕ) (q : GenIdx ι) (n : ℕ) :
    (((k : ℂ) + 1) • brk D (i, k + 1) q n) = -(n : ℂ) • brk D (i, k) q (n - 1) := by

theorem brk_right (D : OPEData ι) (p : GenIdx ι) (j : ι) (l : ℕ) (n : ℕ) :
    (((l : ℂ) + 1) • brk D p (j, l + 1) n)
      = Ttr (brk D p (j, l) n) + (n : ℂ) • brk D p (j, l) (n - 1) := by

theorem brk_mem_filt (D : OPEData ι) (p q : GenIdx ι) (n : ℕ) :
    brk D p q n ∈ filt D.wt ((wtIdx D.wt p : ℤ) + (wtIdx D.wt q : ℤ) - n - 1) := by

theorem brk_vanish (D : OPEData ι) (p q : GenIdx ι) {n : ℕ}
    (h : wtIdx D.wt p + wtIdx D.wt q ≤ n) : brk D p q n = 0 := by
\end{leancode}

\begin{dictionary}
Exactly as stated; the factors $(k+1)$, $(l+1)$ come from the divided-power normalisation. The library also checks the
Virasoro example: \lean{virasoro_brk_three} ($L_{(3)}L = \frac c2\lelevenvac$) and \lean{virasoro_brk_T}
($(TL)_{(1)}L = -L_{(0)}L = -TL$), the latter \emph{derived} from sesquilinearity.
\end{dictionary}

\textbf{Proof on paper.} (1) $k=l=0$. (2) One more application of $L$ or $R$, and the factorials. (3) Each application of $L$
or $R$ raises the weight bound by one (the $T$ in $R$ raises weight by one; the shift $n\mapsto n-1$ in both steps raises
$-n-1$ by one), while the letters' weights rise by one. \qed

\begin{status}
Proved, unconditional; \texttt{\#print axioms} \lean{VOA.}\allowbreak\lean{NonlinearConformal.}\allowbreak\lean{brk_}\allowbreak\lean{mem_}\allowbreak\lean{filt}: standard three.
\end{status}

\begin{physics}
The full singular OPE between all derivatives of the generators is determined by the generators' OPE, and still has only
finitely many singular terms: dimensional analysis survives differentiation.
\end{physics}

\subsection{Relations, admissible ideals and the quotient}

\begin{definition}[The quantum-commutator relation; admissible relation ideals]\label{l11:def:rel}
For letters $p,q$ put
\begin{equation*}
\textstyle\int_{-T}^0[p_\lambda q]\,d\lambda = \sum_{n\ge0}\frac{(-1)^n}{(n+1)!}T^{n+1}(p_{(n)}q),\qquad
\mathrm{rel}(p,q) = p\,q - (-1)^{p(p)p(q)}q\,p - \int_{-T}^0[p_\lambda q]\,d\lambda .
\end{equation*}
An \emph{admissible relation ideal} is a subspace $M\subseteq\mathcal T(R)$ closed under left multiplication by $\mathcal T(R)$,
stable under $T$, and containing every $\mathrm{rel}(p,q)$. $M_{\mathrm{bot}}$ is the smallest one: the span of all
$A\cdot\mathrm{rel}(p,q)$.
\end{definition}
\noindent File \leanpath{VOA/\allowbreak Core/\allowbreak NonlinearConformal.lean}.
\begin{leancode}
noncomputable def opeInt (D : OPEData ι) (p q : GenIdx ι) : TR ι :=
  ∑ n ∈ Finset.range (wtIdx D.wt p + wtIdx D.wt q),
    ((-1 : ℂ) ^ n / ((n + 1).factorial : ℂ)) •
      (((Ttr : End ℂ (TR ι)) ^ (n + 1)) (brk D p q n))

noncomputable def rel (D : OPEData ι) (p q : GenIdx ι) : TR ι :=
  gv p * gv q - koszul (D.par p.1) (D.par q.1) • (gv q * gv p) - opeInt D p q

structure MIdeal (D : OPEData ι) where
  car : Submodule ℂ (TR ι)
  left : ∀ (A x : TR ι), x ∈ car → A * x ∈ car
  Tstable : ∀ x ∈ car, Ttr x ∈ car
  rel_mem : ∀ p q : GenIdx ι, rel D p q ∈ car

noncomputable def MRbotCar (D : OPEData ι) : Submodule ℂ (TR ι) :=
  Submodule.span ℂ (Set.range (fun z : TR ι × GenIdx ι × GenIdx ι => z.1 * rel D z.2.1 z.2.2))

theorem MRbot_le {D : OPEData ι} (M : MIdeal D) : (MRbot D).car ≤ M.car := by
\end{leancode}

\begin{dictionary}
\lean{Finset.}\allowbreak\lean{range (wtIdx D.}\allowbreak\lean{wt p + wtIdx D.}\allowbreak\lean{wt q)} is the range $0\le n<\Delta(p)+\Delta(q)$, beyond which every term vanishes
(Theorem~\ref{l11:thm:brk}(3)); so \lean{opeInt} is exactly $\int_{-T}^0[p_\lambda q]d\lambda$, using
$\int_{-T}^0\lambda^n/n!\,d\lambda = (-1)^nT^{n+1}/(n+1)!$. \lean{koszul a b} is the sign $(-1)^{ab}$
(\leanpath{VOA/\allowbreak Foundations/\allowbreak Parity.lean}). An \lean{MIdeal} is carried as a $\C$-subspace with a
\emph{left}-multiplication field: De Sole--Kac's $M(R)$ is a left ideal, not two-sided, so the quotient is a vector space and
not an algebra. \lean{MRbotCar} spans $A\cdot\mathrm{rel}(p,q)$ over all $A\in\mathcal T(R)$ and all letters $p,q$ (all divided
powers, not only bottom generators). \lean{MRbot} packages it with its three closure proofs; \lean{MRbot_le} is minimality.

\textbf{Relation to Definition 3.4.} Since $N(X,1) = X$, $M_{\mathrm{bot}}$ is the ``$D=1$'' part of De Sole--Kac's $M(R)$; the
library shows below that it is \emph{strictly} smaller and the wrong ideal.
\end{dictionary}

\begin{physics}
$\mathrm{rel}(a,b)$ is the familiar statement that normal ordering is commutative only up to derivatives of the singular OPE:
${:}ab{:}-{:}ba{:} = \sum_{n\ge0}\frac{(-1)^n}{(n+1)!}\partial^{n+1}(a_{(n)}b)$. For the stress tensor:
${:}TT'{:}$-type reorderings produce derivatives of $T$ with no central term, because the central term is a constant and its
derivatives vanish.
\end{physics}

\begin{theorem}[$T$-stability of the relations]\label{l11:thm:Trel}
$T(\mathrm{rel}(p,q)) = (k+1)\,\mathrm{rel}(T^{k+1}a^i/(k+1)!,\,q)+(l+1)\,\mathrm{rel}(p,\,T^{l+1}a^j/(l+1)!)$ for $p = (i,k)$,
$q = (j,l)$. Consequently $M_{\mathrm{bot}}$ is $T$-stable.
\end{theorem}
\noindent File \leanpath{VOA/\allowbreak Core/\allowbreak NonlinearConformal.lean}.
\begin{leancode}
theorem Ttr_rel (D : OPEData ι) (p q : GenIdx ι) :
    Ttr (rel D p q)
      = ((p.2 : ℂ) + 1) • rel D (p.1, p.2 + 1) q
        + ((q.2 : ℂ) + 1) • rel D p (q.1, q.2 + 1) := by
\end{leancode}

\begin{dictionary}
Exactly as stated. This is why relations must be imposed for all divided powers: $T$ of a bottom relation is a combination of
relations involving $TL$.
\end{dictionary}

\textbf{Proof on paper.} $T$ is a derivation, so $T(pq) = (Tp)q+p(Tq)$, and $T\int_{-T}^0[p_\lambda q] =
\int_{-T}^0[Tp_\lambda q]+\int_{-T}^0[p_\lambda Tq]$ by sesquilinearity: at each $n$ the two rules contribute $-n\,p_{(n-1)}q$ and
$+n\,p_{(n-1)}q$, which cancel, and the one extra top term vanishes by the weight bound. \qed

\begin{status}
Proved, unconditional; standard three.
\end{status}

\begin{physics}
Translation invariance of the operator algebra: the derivative of a relation among composite operators is again a
consequence of the relations.
\end{physics}

\begin{definition}[The quotient $U = \mathcal T(R)/M$, vacuum, translation, creation half]\label{l11:def:UR}
For an admissible $M$: $U = \mathcal T(R)/M$ (a vector space), $|0\rangle = [1]$, $T_U$ the descent of $T$, and for each letter $p$
the \emph{creation half} of its field: the $\End(U)$-valued series whose $z^j$-coefficient is left multiplication by
$T^j p/j!$ for $j\ge0$ and $0$ for $j<0$.
\end{definition}
\noindent File \leanpath{VOA/\allowbreak Core/\allowbreak NonlinearConformal.lean}.
\begin{leancode}
abbrev UR (M : MIdeal D) : Type u := TR ι ⧸ M.car

noncomputable def vacU (M : MIdeal D) : UR M := Submodule.Quotient.mk 1

noncomputable def TrU (M : MIdeal D) : End ℂ (UR M) := descendU M Ttr M.Tstable

noncomputable def leftMulU (M : MIdeal D) (A : TR ι) : End ℂ (UR M) :=

noncomputable def dpow (p : GenIdx ι) (k : ℕ) : TR ι :=
  (((p.2 + k).choose k : ℕ) : ℂ) • gv (p.1, p.2 + k)

noncomputable def creHalf (M : MIdeal D) (p : GenIdx ι) : FormalDist (End ℂ (UR M)) :=
  fun j => if 0 ≤ j then leftMulU M (dpow p j.toNat) else 0
\end{leancode}

\begin{dictionary}
\lean{descendU M f hf} is the map induced on the quotient by an $f$ preserving $M$ (\lean{Submodule.mapQ}). Left
multiplication descends because $M$ is a left ideal. \lean{dpow p k} is $T^k(T^ma^i/m!)/k! = \binom{m+k}{k}\,T^{m+k}a^i/(m+k)!$.
\lean{FormalDist (End ℂ (UR M))} is a two-sided formal series $\sum_jf_jz^j$ with operator coefficients
(Chapter~\ref{ch:L01}, Definition~\ref{l01:def:formaldist}); \lean{creHalf M p j} is the $z^j$-coefficient, i.e.\ (in the mode
convention) $p_{(-j-1)}$. The definition sets the annihilation range ($j<0$, i.e.\ modes $p_{(n)}$ with $n\ge0$) to $0$
because it is not constructed here.
\end{dictionary}

\begin{physics}
$\mathcal T(R)/M$ is the would-be space of states: normal-ordered words in the generators modulo the reordering relations. The
vacuum is the empty word, $T_U$ is $L_{-1}$, and the creation modes $a_{(-j-1)}$ of a generator act by putting $\partial^ja/j!$ in front of a
word --- the state--operator correspondence for creation operators.
\end{physics}


\begin{theorem}[Spanning, non-vacuity, covariance of the creation half]\label{l11:thm:URbasic}
For every admissible $M$: the word states $[w]$ span $U$; the creation half is $T_U$-covariant in the creation range,
$[T_U,p_{(-m-1)}] = (m+1)p_{(-m-2)}$ for $m\ge0$. For $M = M_{\mathrm{bot}}$ and \emph{every} OPE datum, $|0\rangle\neq0$, so $U\neq0$.
\end{theorem}
\noindent File \leanpath{VOA/\allowbreak Core/\allowbreak NonlinearConformal.lean}.
\begin{leancode}
theorem wstU_span (M : MIdeal D) : Submodule.span ℂ (Set.range (wstU M)) = ⊤ := by

theorem creHalf_cov_nonneg (M : MIdeal D) (p : GenIdx ι) {m : ℤ} (hm : 0 ≤ m) :
    TrU M * creHalf M p m - creHalf M p m * TrU M = ((m + 1 : ℤ)) • creHalf M p (m + 1) := by

noncomputable def aug : TR ι →ₐ[ℂ] ℂ := FreeAlgebra.lift ℂ (fun _ => (0 : ℂ))

theorem vacU_ne_zero_bot (D : OPEData ι) : vacU (MRbot D) ≠ 0 :=
\end{leancode}

\begin{dictionary}
\lean{wstU M w} is the image of the word $w$. \lean{aug} is the augmentation $\mathcal T(R)\to\C$ killing every letter; the
non-vacuity proof shows it kills $M_{\mathrm{bot}}$. The covariance identity is the $z^m$-coefficient form of $[T,a(z)] =
\partial_za(z)$ restricted to $m\ge0$; the case $m=-1$ involves $p_{(0)}$, which is not constructed.
\end{dictionary}

\textbf{Proof on paper.} Words span $\mathcal T(R)$. $[T,\,\cdot\,]$ of left multiplication by $A$ is left multiplication by $TA$
(Leibniz), and $T\,\mathrm{dpow}(p,m) = (m+1)\mathrm{dpow}(p,m+1)$. For non-vacuity: every relation is killed by the augmentation
(the quadratic terms contain letters; the integral terms are in the image of $T$, which lands in the augmentation ideal),
hence so is $M_{\mathrm{bot}}$, while $\mathrm{aug}(1) = 1$. \qed

\begin{status}
Proved, unconditional; \texttt{\#print axioms} \lean{VOA.}\allowbreak\lean{NonlinearConformal.}\allowbreak\lean{vacU_}\allowbreak\lean{ne_}\allowbreak\lean{zero_}\allowbreak\lean{bot}: standard three. No PBW.
\end{status}

\begin{physics}
The vacuum survives: imposing the commutation relations of normal ordering never sets the identity operator to zero,
because every relation involves at least one field. (Imposing a null relation that \emph{is} the identity --- like
$\omega = 0$ at $c\neq0$, Section~\ref{l11:sec:doctrine} --- can kill the vacuum.) The creation modes of a field acting on the
vacuum module are just ``multiply on the left by the field'', which is why they come for free.
\end{physics}

\subsection{Reconstruction from generating fields}

Before turning to the contract on $U$, we document the general reconstruction theorem that the staging is designed to
feed. Chapter~\ref{ch:L03} refers to it as the word-field corollary of operator-first reconstruction. It upgrades the operator-first
reconstruction of Chapter~\ref{ch:L03} (Theorem~\ref{l03:thm:reconstruct}) from ``a spanning family of states, each with its
field'' to ``a family of generating fields whose monomials span''.

\begin{definition}[Word fields and word states]\label{l11:def:wordfld}
Let $V$ be a vector space with $\lelevenvac\in V$ and fields $(g_k)_{k\in K}$. For a word $l = k_1k_2\cdots k_r$ put
\begin{equation*}
s_l = (g_{k_1})_{(-1)}(g_{k_2})_{(-1)}\cdots(g_{k_r})_{(-1)}\lelevenvac,\qquad
f_l(z) = {:}g_{k_1}(z){:}g_{k_2}(z)\cdots{:}g_{k_r}(z)\,\mathrm{id}{:}\cdots{:}{:},
\end{equation*}
with $s_\emptyset = \lelevenvac$, $f_\emptyset = \mathrm{id}$: $s_l$ is the iterated $(-1)$-mode and $f_l$ the iterated normal-ordered
product, nested to the right.
\end{definition}
\noindent File \leanpath{VOA/\allowbreak Core/\allowbreak NonlinearConformal.lean}.
\begin{leancode}
noncomputable def wordFldF {κ : Type*} (gen : κ → FormalDist (End ℂ V))
    (hgen : ∀ i, IsField (gen i)) : List κ → Fld V
  | [] => Fld.one
  | (i :: l) => Fld.nprod ⟨gen i, hgen i⟩ (wordFldF gen hgen l)

noncomputable def wordFld {κ : Type*} (gen : κ → FormalDist (End ℂ V))
    (hgen : ∀ i, IsField (gen i)) (l : List κ) : FormalDist (End ℂ V) :=
  (wordFldF gen hgen l).1

noncomputable def wordSt {κ : Type*} (gen : κ → FormalDist (End ℂ V)) (vac : V) : List κ → V
  | [] => vac
  | (i :: l) => gen i 0 (wordSt gen vac l)
\end{leancode}

\begin{dictionary}
\lean{Fld V} is the type of fields on $V$ (a series together with a proof of \lean{IsField}); \lean{Fld.one} is the identity
field and \lean{Fld.nprod a b} the normal-ordered product ${:}a(z)b(z){:}$ (\leanpath{VOA/\allowbreak Core/\allowbreak LocalSystem.lean},
Chapter~\ref{ch:L03}, Definition~\ref{l03:def:fld}). \lean{gen i 0} is the $z^0$-coefficient of $g_i$, which in mode convention is
$(g_i)_{(-1)}$, as in the definition.
\end{dictionary}

\begin{physics}
$f_l$ is the composite operator ${:}g_{k_1}g_{k_2}\cdots g_{k_r}{:}$ and $s_l$ the corresponding state: the state--operator map
on monomials in the generators. Derivatives are included automatically when the family $(g_k)$ contains the derivatives of the
generators, as it does in $\mathcal T(R)$, whose letters are $T^ka^i/k!$.
\end{physics}

\begin{theorem}[Reconstruction from generators]\label{l11:thm:recongen}
Let $V$ be a complex vector space, $\lelevenvac\in V$, $T\in\End V$ with $T\lelevenvac = 0$, and let $(g_k)_{k\in K}$ be fields on $V$
such that
\begin{enumerate}
\item $g_k$ and $g_{k'}$ are mutually local for all $k,k'$;
\item $g_k(z)\lelevenvac\in V[[z]]$ (no negative powers of $z$) for all $k$;
\item each $g_k$ is $T$-covariant: $[T,g_k(z)] = \partial_zg_k(z)$;
\item the word states $s_l$ span $V$.
\end{enumerate}
Then there is a linear map $Y:V\to\End(V)[[z^{\pm1}]]$ with $Y(s_l,z) = f_l(z)$ for every word $l$, such that every $Y(u,z)$ is a
field, all $Y(u,z)$ are pairwise local, $Y(u,z)\lelevenvac\in u+zV[[z]]$, and every $Y(u,z)$ is $T$-covariant.
\end{theorem}
\noindent File \leanpath{VOA/\allowbreak Core/\allowbreak NonlinearConformal.lean}
(\lean{VOA.}\allowbreak\lean{NonlinearConformal.}\allowbreak\lean{reconstruct_}\allowbreak\lean{of_}\allowbreak\lean{generators}).
\begin{leancode}
theorem reconstruct_of_generators {κ : Type*} (vac : V) (T : End ℂ V)
    (gen : κ → FormalDist (End ℂ V))
    (hgen : ∀ i, IsField (gen i))
    (hloc : ∀ i j, Local (gen i) (gen j))
    (hvn : ∀ i, ∀ q : ℤ, q < 0 → gen i q vac = 0)
    (hTvac : T vac = 0)
    (hcov : ∀ i, DCovariant T (gen i))
    (hspan : Submodule.span ℂ (Set.range (wordSt gen vac)) = ⊤) :
    ∃ YA : V →ₗ[ℂ] FormalDist (End ℂ V),
      (∀ l : List κ, YA (wordSt gen vac l) = wordFld gen hgen l) ∧
      (∀ u : V, IsField (YA u)) ∧
      (∀ u v : V, Local (YA u) (YA v)) ∧
      (∀ u : V, (YA u) 0 vac = u) ∧
      (∀ u : V, ∀ q : ℤ, q < 0 → (YA u) q vac = 0) ∧
      (∀ u : V, DCovariant T (YA u)) :=
  reconstruct vac (wordSt gen vac) (wordFld gen hgen) T hspan
    (wordFld_isField gen hgen) (wordFld_local gen hgen hloc)
    (wordFld_vac_base gen hgen vac hvn)
    (fun l q hq => wordFld_vac_neg gen hgen vac hvn l q hq)
    hTvac (wordFld_cov gen hgen hcov)
\end{leancode}

\begin{dictionary}
\lean{κ} is the index set $K$; \lean{List κ} the words; \lean{hgen} says each $g_k$ is a field; \lean{hloc} is (1); \lean{hvn} is
(2), since the $z^q$-coefficient with $q<0$ applied to the vacuum must vanish; \lean{hTvac} is $T\lelevenvac = 0$; \lean{hcov} is (3)
(\lean{DCovariant}, Definition~\ref{l03:def:dcov}); \lean{hspan} is (4). The conclusion is an existential over a linear map
\lean{YA : V →ₗ[ℂ] FormalDist (End ℂ V)} with six conjuncts: $Y(s_l) = f_l$; fields; locality; $(YA\,u)_0\lelevenvac = u$ (the
$z^0$-coefficient on the vacuum is $u$); vanishing of negative powers on the vacuum; $T$-covariance. The term-mode proof is
literally \lean{reconstruct} applied to the family of all word states and word fields.

\textbf{Comparison with the literature.} This is the existence part of Kac's reconstruction theorem (\emph{Vertex algebras for
beginners}, Thm.~4.5) and of Frenkel--Ben-Zvi Thm.~4.4.1: from translation-covariant, mutually local generating fields whose
monomials span, one gets a vertex algebra. Two differences. (i) Kac and FBZ use monomials in \emph{all} negative modes
$a^{k}_{(-j-1)}$; the library uses only $(-1)$-modes, so the family $(g_k)$ must itself contain the derivatives (in $U$ the letters
are $T^ja/j!$, so this is automatic). (ii) The conclusion does not package a \lean{PreVertexAlgebra}: no grading and no vacuum axiom
$Y(\lelevenvac) = \mathrm{id}$ are produced (the latter is the empty word: $Y(s_\emptyset) = f_\emptyset = \mathrm{id}$, so it does
hold when $\lelevenvac$ is a word state, which it always is). Uniqueness of $Y$ is not stated; it follows from Goddard uniqueness.
\textbf{Could it be vacuous?} No: the hypotheses are satisfiable by every library vertex algebra with its generating fields, and
the library uses this theorem for the Moonshine fields (\leanpath{VOA/\allowbreak Moonshine/\allowbreak VnatFields.lean} imports
the file).
\end{dictionary}

\textbf{Proof on paper.} Apply the operator-first reconstruction (Theorem~\ref{l03:thm:reconstruct}) to the family $(s_l,f_l)_l$.
Its hypotheses: the $s_l$ span (4); each $f_l$ is a field, because normal-ordered products of fields are fields; the $f_l$ are
pairwise local, because the set of fields generated from a pairwise-local set by normal-ordered products is pairwise local (Dong's
lemma, packaged as \lean{Core.}\allowbreak\lean{LocalSystem.}\allowbreak\lean{generated_}\allowbreak\lean{local}, Chapter~\ref{ch:L03}, Theorem~\ref{l03:thm:genlocal}); $f_l(z)\lelevenvac$
has no negative powers and constant term $s_l$, by induction on $l$ using (2) and the formula
${:}a(z)b(z){:}\lelevenvac = a(z)_+b(z)\lelevenvac$; and each $f_l$ is $T$-covariant, because $T$-covariance is preserved by normal
ordering (Leibniz rule for $\partial_z$). \qed

\begin{status}
Proved, unconditional; \texttt{\#print axioms} \lean{VOA.}\allowbreak\lean{NonlinearConformal.}\allowbreak\lean{reconstruct_}\allowbreak\lean{of_}\allowbreak\lean{generators}: standard three (run for this
chapter).
\end{status}

\begin{physics}
This is the form of reconstruction used whenever a chiral algebra is generated by finitely many fields: the states are the
normal-ordered monomials ${:}g_{k_1}\cdots g_{k_r}{:}$ applied to the vacuum, which span the space of states. Every operator of the
theory is a normal-ordered composite of the generating currents (for a $W$-algebra: of $T,W_3,\dots$, and their derivatives), and
one has to check only the locality of the generators (finite-order poles in their OPEs), their translation covariance, and that
the monomials span.
\end{physics}

Two bridges, valid on any vector space, turn weight bookkeeping and an OPE table into the hypotheses of
Theorem~\ref{l11:thm:recongen}.

\begin{theorem}[Fields from a weight filtration; locality from the OPE]\label{l11:thm:bridges}
\begin{enumerate}
\item If $V$ is exhausted by subspaces $F_d$ ($d\in\Z$) with $F_d = 0$ for $d<0$, and the $z^j$-coefficient of $a$ maps $F_d$ into
  $F_{d+\Delta+j}$, then $a$ is a field.
\item If $[a_{(m)},b_{(n)}]$-type commutators are given by finitely many OPE coefficient fields $c_0,\dots,c_{N-1}$ through the
  falling-factorial (commutator) formula, then $a$ and $b$ are mutually local.
\end{enumerate}
\end{theorem}
\noindent File \leanpath{VOA/\allowbreak Core/\allowbreak NonlinearConformal.lean}.
\begin{leancode}
theorem isField_of_filt (a : FormalDist (End ℂ V)) (F : ℤ → Submodule ℂ V) (Δ : ℤ)
    (hbot : ∀ d : ℤ, d < 0 → F d = ⊥)
    (hcov : ∀ x : V, ∃ d : ℤ, x ∈ F d)
    (hmap : ∀ (j d : ℤ) (x : V), x ∈ F d → a j x ∈ F (d + Δ + j)) :
    IsField a := by

theorem local_of_ope (a b : FormalDist (End ℂ V)) (N : ℕ) (c : Fin N → FormalDist (End ℂ V))
    (h : ∀ m n : ℤ, a m * b n - b n * a m
        = ∑ j : Fin N, (∏ i ∈ Finset.range (j : ℕ), (-m - 1 - (i : ℤ))) • c j (m + n + 1 + j)) :
    Local a b := by
\end{leancode}

\begin{dictionary}
In (2), \lean{a m}, \lean{b n} are $z^m$-, $w^n$-coefficients; \lean{∏ i ∈ Finset.}\allowbreak\lean{range j, (-m - 1 - i)} is the falling factorial
$(-m-1)(-m-2)\cdots(-m-j)$, i.e.\ $j!\binom{-m-1}{j}$; the hypothesis is the commutator formula
$[a(z),b(w)] = \sum_{j<N}c_j(w)\partial_w^{(j)}\delta(z,w)$ written coefficientwise in the library's $z$-power indexing.
\lean{local_iff_deltaSum} (Chapter~\ref{ch:L01}) converts it to locality of order $N$.
\end{dictionary}

\textbf{Proof on paper.} (1) A vector of filtration degree $d$ is killed by coefficients with $d+\Delta+j<0$. (2) A finite sum of
derivatives of $\delta(z-w)$ is annihilated by $(z-w)^N$. \qed

\begin{status}
Proved, unconditional; standard three. Neither bridge is used by the Virasoro instance below (which transports locality).
\end{status}

\begin{physics}
(1) is the statement that a field of definite conformal weight acting on a state of finite energy can lower the energy only
finitely often. (2) is the classic equivalence ``finite singular OPE $\Leftrightarrow$ locality''.
\end{physics}

\subsection{The contract \texorpdfstring{\lean{GenFields}}{GenFields}, and a refuted route}

\begin{definition}[Generating-field contract]\label{l11:def:genfields}
For an admissible ideal $M$, a \emph{\lean{GenFields M}} is a family of fields $\{Y(p,z)\}$ on $U = \mathcal T(R)/M$, one per letter $p$,
which are pairwise local, $T_U$-covariant, kill $|0\rangle$ in negative powers of $z$, and whose creation half is the one of
Definition~\ref{l11:def:UR}.
\end{definition}
\noindent File \leanpath{VOA/\allowbreak Core/\allowbreak NonlinearConformal.lean}.
\begin{leancode}
structure GenFields {D : OPEData ι} (M : MIdeal D) where
  fld : GenIdx ι → FormalDist (End ℂ (UR M))
  isFld : ∀ p, IsField (fld p)
  loc : ∀ p q, Local (fld p) (fld q)
  cre : ∀ (p : GenIdx ι) (j : ℤ), 0 ≤ j → fld p j = creHalf M p j
  cre_neg : ∀ p, ∀ j : ℤ, j < 0 → (fld p) j (vacU M) = 0
  cov : ∀ p, DCovariant (TrU M) (fld p)

theorem UR_YA :
    ∃ YA : UR M →ₗ[ℂ] FormalDist (End ℂ (UR M)),
      (∀ l : List (GenIdx ι), YA (wordSt G.fld (vacU M) l) = wordFld G.fld G.isFld l) ∧
      (∀ u : UR M, IsField (YA u)) ∧
      (∀ u v : UR M, Local (YA u) (YA v)) ∧
      (∀ u : UR M, (YA u) 0 (vacU M) = u) ∧
      (∀ u : UR M, ∀ j : ℤ, j < 0 → (YA u) j (vacU M) = 0) ∧
      (∀ u : UR M, DCovariant (TrU M) (YA u)) :=
\end{leancode}

\begin{dictionary}
The only freedom left in \lean{fld} is the annihilation range (\lean{cre} fixes all $z^j$, $j\ge0$); \lean{UR_YA} is
Theorem~\ref{l11:thm:recongen} applied to \lean{G.fld}, with spanning supplied by \lean{wstU_span}. It is a theorem
\emph{conditional on a term} \lean{G : GenFields M} (the \lean{variable} binder). The only general witness, \lean{emptyGenFields},
has no generators.
\end{dictionary}

\begin{physics}
\lean{GenFields} is exactly what a physicist means by ``the OPE defines the fields'': it asks for the annihilation modes of each
generator, compatible with the creation modes already present. Providing it for a nonlinear OPE in general is De Sole--Kac's
Lemma~3.2 and \S4 --- the hard theorem.
\end{physics}

\begin{warning}
\textbf{Dead end: the left ideal $M_{\mathrm{bot}}$ is the wrong ideal (refuted in Lean).} The original staging offered
$M_{\mathrm{bot}}$ as the ideal over which to populate \lean{GenFields}. For the Virasoro datum this is \emph{false}:
\begin{leancode}
theorem not_nonempty_genFields_MRbot_virasoro (c : ℂ) :
    ¬ Nonempty (GenFields (MRbot (virasoroOPE c))) := by
\end{leancode}
(\leanpath{VOA/\allowbreak Core/\allowbreak NonlinearVirasoro.lean}). Reason: \lean{cre} forces $L_{(-m-1)}$ to be left
multiplication by $T^mL/m!$ for all $m\ge0$, so $[L_{(-m-1)},L_{(-n-1)}]$ is left multiplication by $g_mg_n - g_ng_m$
(with $g_k = T^kL/k!$). In a vertex algebra this commutator must be a finite combination of single modes, which requires the
relation $\mathrm{rel}(g_m,g_n)$ to be usable \emph{inside} a word, i.e.\ multiplied on the right. The coordinate of the
length-three word $[g_{N+1},g_0,g_0]$ vanishes on $M_{\mathrm{bot}}$ but locality of order $N$ would force it to be $\pm1$.
Accordingly $M_{\mathrm{bot}}$ is strictly smaller than the two-sided ideal (witness $\mathrm{rel}(L,TL)\cdot L$):
\begin{leancode}
theorem MRbot_lt_MRtwo_virasoro (c : ℂ) :
    (MRbot (virasoroOPE c)).car < (MRtwo (virasoroOPE c)).car := by
\end{leancode}
Any theorem with a hypothesis \lean{G : GenFields (MRbot _}\allowbreak\lean{)} for the Virasoro datum is vacuous. Both statements:
\texttt{\#print axioms}: standard three. (In De Sole--Kac's own $M(R)$ the $D\neq1$ generators supply exactly the missing
right multiplications.)
\end{warning}

\begin{definition}[The two-sided relation ideal]\label{l11:def:MRtwo}
$M_{\mathrm{two}}$ is the span of all $A\cdot\mathrm{rel}(p,q)\cdot B$, $A,B\in\mathcal T(R)$: the two-sided ideal generated by the
relations. It is admissible.
\end{definition}
\noindent File \leanpath{VOA/\allowbreak Core/\allowbreak NonlinearVirasoro.lean}.
\begin{leancode}
noncomputable def MRtwoCar (D : OPEData ι) : Submodule ℂ (TR ι) :=
  Submodule.span ℂ (Set.range (fun z : TR ι × GenIdx ι × GenIdx ι × TR ι =>
    z.1 * rel D z.2.1 z.2.2.1 * z.2.2.2))

noncomputable def MRtwo (D : OPEData ι) : MIdeal D where
\end{leancode}

\begin{dictionary}
For a \emph{linear} datum (every OPE value in $\C\cdot1\oplus R$, as for Virasoro), De Sole--Kac's (3.1)--(3.2) give $N(X,D) = XD$ for
$X\in R\oplus\C1$, so Definition~3.4 is exactly $M_{\mathrm{two}}$; the library proves this reconciliation for any interface term
(\lean{MRfullCar_eq_MRtwoCar}, Section~\ref{l11:sec:doctrine}). For nonlinear data $M_{\mathrm{two}}$ is \emph{not} De Sole--Kac's
$M(R)$. The definition mentions only the OPE table (through \lean{rel}): no model, no kernel of a map.
\end{dictionary}

\begin{physics}
For a Lie-algebra-like (linear) OPE such as Virasoro or affine currents, the presented chiral algebra is the universal enveloping
algebra of the negative-mode Lie algebra: one may reorder modes anywhere inside a word using the commutation relations. That is
exactly the two-sided ideal.
\end{physics}

\subsection{Generators plus OPE equals Borcherds, for Virasoro}

\begin{theorem}[The Virasoro relations in closed form; independence of $c$]\label{l11:thm:virrel}
For the Virasoro datum and letters $g_k = T^kL/k!$:
\begin{equation*}
\mathrm{rel}(g_k,g_l) = g_kg_l - g_lg_k - (l-k)\,g_{k+l+2}.
\end{equation*}
In particular the relations, and hence $M_{\mathrm{two}}$, do not depend on the central charge $c$.
\end{theorem}
\noindent File \leanpath{VOA/\allowbreak Core/\allowbreak NonlinearVirasoro.lean}; the independence of $M_{\mathrm{two}}$
from $c$ was checked in Lean against the library (the check is not part of this release).
\begin{leancode}
theorem virasoro_rel (c : ℂ) (k l : ℕ) :
    rel (virasoroOPE c) ((), k) ((), l)
      = gv ((), k) * gv ((), l) - gv ((), l) * gv ((), k)
        - ((l : ℂ) - k) • gv ((), k + l + 2) := by
\end{leancode}

\begin{dictionary}
Under $g_k\mapsto L_{-k-2}$ this is the Witt bracket $[L_{-k-2},L_{-l-2}] = (l-k)L_{-k-l-4}$ (no central term since
$-k-l-4\neq0$). The central term of the OPE disappears because $\int_{-T}^0$ applies at least one $T$ to $L_{(3)}L = \frac c2\cdot1$,
and $T1 = 0$. The check outside the library states the consequence: \lean{MRtwoCar (virasoroOPE c) = MRtwoCar (virasoroOPE c')} for all $c,c'$.
\end{dictionary}

\textbf{Proof on paper.} First take $l=0$. By left sesquilinearity, $(g_k)_{(n)}L = \frac{(-1)^k}{k!}n(n-1)\cdots(n-k+1)\,L_{(n-k)}L$,
so only $n = k$ ($L_{(0)}L = TL$), $n = k+1$ ($L_{(1)}L = 2L$) and $n = k+3$ (the central term) contribute. Applying
$\frac{(-1)^n}{(n+1)!}T^{n+1}$ gives $(k+2)g_{k+2}$, $-2(k+1)g_{k+2}$ and $0$ (since $T^{k+4}1 = 0$), so
$\int_{-T}^0[g_{k\,\lambda}L] = -k\,g_{k+2}$. The general case follows by induction on $l$ using the $T$-equivariance of
$\int_{-T}^0$ (the identity behind Theorem~\ref{l11:thm:Trel}). \qed

\begin{status}
Proved, unconditional; standard three (library; the independence of $M_{\mathrm{two}}$ in the check
outside the library).
\end{status}

\begin{physics}
As a vector space, the universal Virasoro vacuum module is the same for all $c$ (a PBW space on $L_{-2},L_{-3},\dots$); the central
charge lives entirely in the annihilation modes ($L_{n}$, $n\ge-1$) acting on it. So any ``presentation'' that only fixes the vector
space and creation operators cannot see $c$.
\end{physics}

\begin{theorem}[$M_{\mathrm{two}} = \ker\Phi$; the presentation is the Virasoro VOA]\label{l11:thm:virpres}
Let $\Phi:\mathcal T(R)\to V(c,0)$ be the linear map $g_{k_1}\cdots g_{k_r}\mapsto L_{-k_1-2}\cdots L_{-k_r-2}\lelevenvac$ into the
universal Virasoro vacuum module. Then:
\begin{enumerate}
\item $M_{\mathrm{two}} = \ker\Phi$; so $\bar\Phi:\mathcal T(R)/M_{\mathrm{two}}\to V(c,0)$ is a linear isomorphism;
\item the sorted word states $[g_{k_1}\cdots g_{k_r}]$, $k_1\ge\cdots\ge k_r$, form a basis of $\mathcal T(R)/M_{\mathrm{two}}$ (spanning and
  independence);
\item there are generating fields \lean{genFieldsVirTwo c : GenFields (MRtwo (virasoroOPE c))}, and an isomorphism
  $e:\mathcal T(R)/M_{\mathrm{two}}\to V(c,0)$ sending $|0\rangle\mapsto\lelevenvac$, $T_U\mapsto L_{-1}$, and each generating field of
  $g_k$ to $Y(L_{-k-2}\lelevenvac,z)$ --- in particular the field of $g_0$ to $Y(\omega,z)$;
\item the modes $L_m := (\text{field of }g_0)_{(m+1)}$ on $\mathcal T(R)/M_{\mathrm{two}}$ satisfy the Virasoro relations with central charge
  $c$, and $\mathcal T(R)/M_{\mathrm{two}}$ carries a \lean{PreVertexAlgebra} structure whose vertex operator at $g_k$ is the generating field
  (the library's pre-vertex algebra record \lean{PreVertexAlgebra}, Chapter~\ref{ch:L03}, which does not include translation
  covariance; it is not the library's translation-covariant \lean{VertexAlgebra} of \leanpath{Core/\allowbreak GenuineVertexAlgebra.lean}.
  Here it is transported from the VOA $V(c,0)$, so covariance holds by transport, but the type does not record it).
\end{enumerate}
\end{theorem}
\noindent Files \leanpath{VOA/\allowbreak Examples/\allowbreak NonlinearVirasoroTransport.lean},
\leanpath{VOA/\allowbreak Examples/\allowbreak NonlinearVirasoroGenuine.lean}.
\begin{leancode}
noncomputable def psiVir : TR Unit →ₐ[ℂ] UVir :=
  FreeAlgebra.lift ℂ (fun p : GenIdx Unit => ι ℂ (L (-(p.2 : ℤ) - 2)))

noncomputable def PhiVir (c : ℂ) : TR Unit →ₗ[ℂ] VirUniv c where
  toFun X := Submodule.Quotient.mk (psiVir X)

theorem MRtwo_eq_ker_PhiVir (c : ℂ) :
    (MRtwo (virasoroOPE c)).car = LinearMap.ker (PhiVir c) := by

theorem sortedWstUw_span (c : ℂ) :
    Submodule.span ℂ (Set.range (fun w : {w : List ℕ // SortedWord w} => wstUw c w.1)) = ⊤ := by

theorem sortedWordSt_linearIndependent_two (c : ℂ) :
    LinearIndependent ℂ (fun w : {w : List ℕ // SortedWord w} => wstUw c w.1) := by

noncomputable def fldVirTwo (c : ℂ) (p : GenIdx Unit) :
    FormalDist (End ℂ (UR (MRtwo (virasoroOPE c)))) :=
  conjDist (eVirTwo c) (virYA c (PhiVir c (gv p)))

theorem virasoro_presentation_equivalence_genuine (c : ℂ) :
    ∃ e : UR (MRtwo (virasoroOPE c)) ≃ₗ[ℂ] VirUniv c,
      e (vacU (MRtwo (virasoroOPE c))) = (virasoroVOA c).vacuum
      ∧ (∀ u : UR (MRtwo (virasoroOPE c)),
            e (TrU (MRtwo (virasoroOPE c)) u) = Ltrans c (e u))
      ∧ (∀ (k : ℕ) (j : ℤ) (u : UR (MRtwo (virasoroOPE c))),
            e ((genFieldsVirTwo c).fld ((), k) j u)
              = (virasoroVOA c).Y (negWordVir c [k]) j (e u))
      ∧ (∀ (j : ℤ) (u : UR (MRtwo (virasoroOPE c))),
            e ((genFieldsVirTwo c).fld ((), 0) j u)
              = (virasoroVOA c).Y (virasoroVOA c).conformalVector j (e u)) :=

theorem fldVirTwo_virasoro (c : ℂ) (m n : ℤ) :
    mode (fldVirTwo c ((), 0)) (m + 1) * mode (fldVirTwo c ((), 0)) (n + 1)
        - mode (fldVirTwo c ((), 0)) (n + 1) * mode (fldVirTwo c ((), 0)) (m + 1)
      = ((m : ℂ) - n) • mode (fldVirTwo c ((), 0)) (m + n + 1)
        + (if m + n = 0 then ((m : ℂ) ^ 3 - m) / 12 * c else 0) •
            (1 : End ℂ (UR (MRtwo (virasoroOPE c)))) := by

noncomputable def virasoroURVATwo (c : ℂ) : PreVertexAlgebra (UR (MRtwo (virasoroOPE c))) :=
  transportVA (virasoroVOA c).toPreVertexAlgebra (eVirTwo c)
\end{leancode}

\begin{dictionary}
\lean{UVir} is the universal enveloping algebra of the Virasoro Lie algebra, \lean{ι ℂ (L n)} the image of $L_n$ in it, and
\lean{VirUniv c} the universal vacuum module $V(c,0)$ (Chapter~\ref{ch:L06}, Definition~\ref{l06:def:Vc}); \lean{psiVir} is the algebra
map $g_k\mapsto L_{-k-2}$ and \lean{PhiVir c X} the class of \lean{psiVir X} applied to the vacuum. \lean{wstUw c w} is the word state
of the list $w$ of indices; \lean{SortedWord w} means $w$ is weakly decreasing (Definition~\ref{l06:def:sorted}). \lean{virYA c} is the
vertex operator map of the universal Virasoro VOA (Theorem~\ref{l06:thm:virYA}), \lean{negWordVir c [k]} $= L_{-k-2}\lelevenvac$,
\lean{Ltrans c} $= L_{-1}$, \lean{virasoroVOA c} the universal Virasoro VOA (Definition~\ref{l06:def:virVOA}). \lean{conjDist e a} is
the conjugated series $e^{-1}\circ a(z)\circ e$; \lean{transportVA A e} transports a vertex algebra structure along a linear
isomorphism.

\textbf{What is genuine, and what is by construction.} The ideal \lean{MRtwo (virasoroOPE c)} is defined from the OPE table alone;
(1) and (2) are honest theorems (spanning by a straightening argument in the quotient, which is where two-sidedness is used;
independence imported from the $\mathtt{pbw}$-free partition basis of Chapter~\ref{ch:L06}). In (3)--(4), however, the fields are
\emph{defined} by conjugating the already constructed $Y$ of $V(c,0)$ across $e = \bar\Phi$ (\lean{fldVirTwo}), and
\lean{virasoroURVATwo} is the transported \lean{PreVertexAlgebra} (the library's record without translation covariance; it inherits
covariance from $V(c,0)$ only because it is transported, and its type does not say so). So the field-matching clauses of
\lean{virasoro_}\allowbreak\lean{presentation_}\allowbreak\lean{equivalence_}\allowbreak\lean{genuine} hold essentially by construction; their content is that the transported fields
satisfy the contract --- above all \lean{cre}, which pins their creation half to left multiplication in the presentation (the
numerical identity $\mathrm{dpow}((),k)\,n = \binom{k+n}n g_{k+n}$ against the coefficient of $\frac1{k!}\partial^kY(\omega,z)$). The
annihilation half is \emph{not} derived from the OPE (De Sole--Kac Lemma~3.2 and \S4 are not formalised; the docstrings say so), and
\lean{local_of_ope} is not used. Combined with Theorem~\ref{l11:thm:virrel}, the honest summary is: \emph{the vector space and the
creation operators presented by one weight-two generator with the Virasoro OPE are those of $V(c,0)$, for every $c$; the vertex
algebra structure on it is the transported one.} The theorem name ``generators $+$ OPE $=$ Borcherds'' (file headers) is therefore
stronger than what is proved, though the proved statement is correct and non-circular.

\textbf{Non-vacuity.} $U\neq0$ and the sorted monomials are independent, so the presentation is not the zero space; the Virasoro
relations in (4) hold with the given $c$ (this is where $c$ re-enters, through the transported annihilation modes).
\end{dictionary}

\textbf{Proof on paper.} $\Phi$ kills the relations because they are the Witt relations, which hold in $U(\mathrm{Vir})$; and $\Phi$ is
$T$-equivariant and left-$\mathcal T(R)$-linear, so $M_{\mathrm{two}}\subseteq\ker\Phi$. For the converse, the swap identity
$[g_ng_jW] = [g_jg_nW]+(j-n)[g_{n+j+2}W]$ in the quotient (a right multiple of a relation) and induction on the energy show that
sorted words span $\mathcal T(R)/M_{\mathrm{two}}$; their images $L_{-k_1-2}\cdots L_{-k_r-2}\lelevenvac$ are linearly independent in
$V(c,0)$ for every $c$ (partition basis). So $\bar\Phi$ maps a spanning set bijectively onto an independent set: it is injective,
and it is surjective because the creation monomials span $V(c,0)$. Everything else is transport of structure along $\bar\Phi$. \qed

\begin{status}
Proved, unconditional, \texttt{pbw}-free; \texttt{\#print axioms} of \lean{VOA.}\allowbreak\lean{VirModule.}\allowbreak\lean{MRtwo_}\allowbreak\lean{eq_}\allowbreak\lean{ker_}\allowbreak\lean{PhiVir},
\lean{…sortedWstUw_span}, \lean{…sortedWordSt_}\allowbreak\lean{linearIndependent_}\allowbreak\lean{two}, \lean{…virasoro_}\allowbreak\lean{presentation_}\allowbreak\lean{equivalence_}\allowbreak\lean{genuine},
\lean{…fldVirTwo_virasoro}, \lean{…virasoroURVATwo}: standard three (run for this chapter). The earlier versions over
\lean{MVir c := ker Φ} (\lean{virasoro_}\allowbreak\lean{presentation_}\allowbreak\lean{equivalence}, \lean{genFieldsVir}; file
\leanpath{NonlinearVirasoroTransport.lean}) are circular by their own docstring and superseded by the ``\lean{Two}'' versions.
\end{status}

\begin{physics}
``The Virasoro algebra is generated by $T(z)$ with the $TT$ OPE.'' What the library verifies is the counting half of that sentence:
the states built from $T$ and its derivatives, modulo the reordering relations dictated by the singular OPE, are exactly the
Virasoro descendants of the vacuum, one for each partition into parts $\ge2$. The dynamical half --- that the OPE alone determines
the annihilation modes and their locality --- is the content of the De Sole--Kac construction, which is staged but not completed.
\end{physics}

\subsection{Gradings on \texorpdfstring{$\mathcal T(R)$}{T(R)}: exact weight and the De Sole--Kac filtration}\label{l11:sec:dskgrading}

\begin{definition}[Exact weight; conformal data; positive weight; the $\zeta$-filtration]\label{l11:def:gr}
$\mathrm{gr}(d)$ is the span of words of weight exactly $d$; an OPE datum is \emph{conformal} if $a^i_{(n)}a^j\in\mathrm{gr}(\Delta_i+\Delta_j-n-1)$;
$\mathrm{grPos}$ is the span of words of positive weight. The \emph{$\zeta$-weight} of a letter $T^ka^i/k!$ is $\Delta_i$ (the power of $T$ is
ignored), and $\mathrm{zfilt}(d)$ is the span of words of $\zeta$-weight $\le d$.
\end{definition}
\noindent Files \leanpath{VOA/\allowbreak Core/\allowbreak NonlinearVirasoro.lean}, \leanpath{VOA/\allowbreak Core/\allowbreak NonlinearZeta.lean}.
\begin{leancode}
noncomputable def gr (wt : ι → ℕ) (d : ℕ) : Submodule ℂ (TR ι) :=
  Submodule.span ℂ (wbasis ι '' {w | wtWord wt w = d})

def OPEData.IsConformal (D : OPEData ι) : Prop :=
  ∀ i j (n : ℕ), D.ope i j n ∈ gr D.wt (D.wt i + D.wt j - n - 1)

noncomputable def grPos (wt : ι → ℕ) : Submodule ℂ (TR ι) :=
  Submodule.span ℂ (wbasis ι '' {w | 0 < wtWord wt w})

def zwtIdx : GenIdx ι → ℕ := fun p => wt p.1

noncomputable def zfilt (d : ℤ) : Submodule ℂ (TR ι) :=
  Submodule.span ℂ (wbasis ι '' {w | ((zwtWord wt w : ℤ)) ≤ d})
\end{leancode}

\begin{dictionary}
\lean{IsConformal} is an \emph{extra} predicate, so that genuinely non-conformal data (allowed by De Sole--Kac) remain expressible.
In \lean{IsConformal} the subtraction is in $\N$; for $n+1>\Delta_i+\Delta_j$ the coefficient is already $0$, so truncation imposes nothing.
\end{dictionary}

\begin{physics}
The exact weight is the conformal dimension $L_0$; positive weight means ``not proportional to the identity''. The $\zeta$-weight counts
how composite an operator is, ignoring derivatives.
\end{physics}


\begin{theorem}[Gradings: what holds]\label{l11:thm:gradings}
\begin{enumerate}
\item The Virasoro datum is conformal.
\item For conformal data, $\mathrm{rel}(p,q)$ is homogeneous of weight $\Delta(p)+\Delta(q)$.
\item $T$ preserves $\mathrm{zfilt}(d)$, and the library's grading condition implies De Sole--Kac's: $a^i_{(n)}a^j\in\mathrm{zfilt}(\Delta_i+\Delta_j-1)$.
\item $x\in\mathrm{grPos}$ iff its weight-$0$ component vanishes.
\end{enumerate}
\end{theorem}
\noindent Files \leanpath{VOA/\allowbreak Core/\allowbreak NonlinearVirasoro.lean}, \leanpath{VOA/\allowbreak Core/\allowbreak NonlinearZeta.lean}.
\begin{leancode}
theorem virasoroOPE_isConformal (c : ℂ) : (virasoroOPE c).IsConformal := by

theorem rel_mem_gr {D : OPEData ι} (hC : D.IsConformal) (p q : GenIdx ι) :
    rel D p q ∈ gr D.wt (wtIdx D.wt p + wtIdx D.wt q) := by

theorem Ttr_mem_zfilt (wt : ι → ℕ) {d : ℤ} {x : TR ι} (hx : x ∈ zfilt wt d) :
    Ttr x ∈ zfilt wt d := by

theorem ope_mem_zfilt (D : OPEData ι) (i j : ι) (n : ℕ) :
    D.ope i j n ∈ zfilt D.wt ((D.wt i : ℤ) + (D.wt j : ℤ) - 1) :=

theorem mem_grPos_iff (wt : ι → ℕ) (x : TR ι) :
    x ∈ grPos wt ↔ grProj wt 0 x = 0 := by
\end{leancode}

\begin{dictionary}
\lean{grProj wt 0 x} is the weight-$0$ component (a finite sum, computed from the word-basis coordinates). (3) is the precise sense in
which \lean{OPEData} feeds De Sole--Kac's Lemma~3.2 recursion, which runs on a filtration where $T$ has degree $0$: the library's
conformal-weight filtration gives $T$ degree $+1$, so a recursion on it could not terminate. The $\zeta$ results are unconditional but,
by the file's own account, not yet used by any proof outside the interface fields \lean{N_zfilt}, \lean{P_zfilt} below.
\end{dictionary}

\textbf{Proof on paper.} (1) Inspection of the table. (2) Each term of $\mathrm{rel}$ has weight $\Delta(p)+\Delta(q)$: the products by
additivity, $T^{n+1}(p_{(n)}q)$ by the conformal condition plus $n+1$. (3) $T$ changes only the $T$-power of a letter. Since
$\zeta\le\Delta$ wordwise and $\Delta_i+\Delta_j-n-1\le\Delta_i+\Delta_j-1$. (4) Decompose into weight components. \qed

\begin{status}
Proved, unconditional; standard three.
\end{status}

\begin{physics}
Two different ``sizes'' of a composite operator: its conformal dimension (derivatives count) and its ``number of constituents''
weighted by the generators' dimensions (derivatives do not count). The second is the right induction parameter for constructing the
normal-ordered product of composites, because differentiating does not make an operator more composite.
\end{physics}


%======================================================================
\section{The non-linear layer: an interface with no instance}\label{l11:sec:doctrine}
%======================================================================

For OPE data whose singular terms contain genuine normally ordered products (the $W_3$ algebra, $\mathcal N=2$ chiral
algebras, nonlinear $W$-algebras), $M(R)$ cannot be written down without De Sole--Kac's Lemma~3.2. The library stages
this ``Stage A'' as an \emph{interface}: a structure listing all the properties of the extended normally ordered product
$N$ and bracket $P$, over which everything downstream is parametrised.

\begin{warning}
\textbf{No instance of \lean{NPExtension} exists anywhere in the library (open gap).} A search of the whole of \leanpath{VOA/}
finds no term of type \lean{NPExtension D} for any \lean{D}; the existence--uniqueness theorem (``Stage B'') is not
formalised. Every theorem of this section that takes \lean{NP : NPExtension D} as a hypothesis is therefore, as a Lean
proposition, \emph{vacuously true today}: it would remain provable if \lean{NPExtension} were uninhabitable. The library's
docstrings say exactly this (``every downstream statement quantified over \lean{NP : NPExtension D} is a conditional whose
antecedent nobody has supplied''). We have \emph{not} checked that \lean{NPExtension} is satisfiable; its 17 fields were
audited against De Sole--Kac by the authors (per the docstring), and one near-inconsistency (the shared cutoff in
\lean{wickRight}) is documented there. A reader must treat this whole section as a precise \emph{specification}, not as
mathematics established in Lean.
\end{warning}

\begin{definition}[The extension interface, De Sole--Kac Lemma 3.2]\label{l11:def:npext}
An \lean{NPExtension D} consists of bilinear maps $N:\mathcal T(R)\times\mathcal T(R)\to\mathcal T(R)$ and
$P:\mathcal T(R)\times\mathcal T(R)\to\mathcal T(R)^{\N}$ (the slices $P(A,B)(n) = A_{(n)}B$) satisfying: $N(1,A) = N(A,1) = A$;
$N(a,B) = aB$ for a letter $a$; $P(1,A) = P(A,1) = 0$; $P$ restricts to the OPE table on bottom letters; extended
sesquilinearity $P(TA,B) = L(P(A,B))$, $P(A,TB) = R(P(A,B))$; $T$ is a derivation of $N$; quasi-associativity, the left and
right Wick formulas (De Sole--Kac (3.3), (3.6), (3.7), in slice form with explicit cutoffs); the grading bounds (3.8) in the
$\zeta$-filtration; and the conformal-weight termination bound $P(A,B)(n)\in\mathrm{filt}(\Delta A+\Delta B-n-1)$.
\end{definition}
\noindent File \leanpath{VOA/\allowbreak Core/\allowbreak NPExtension.lean}.
\begin{leancode}
structure NPExtension (D : OPEData ι) where
  N : TR ι →ₗ[ℂ] TR ι →ₗ[ℂ] TR ι
  P : TR ι →ₗ[ℂ] TR ι →ₗ[ℂ] (ℕ → TR ι)
  N_one_left : ∀ A : TR ι, N 1 A = A
  N_one_right : ∀ A : TR ι, N A 1 = A
  N_gv : ∀ (p : GenIdx ι) (B : TR ι), N (gv p) B = gv p * B
  P_one_left : ∀ (A : TR ι) (n : ℕ), P 1 A n = 0
  P_one_right : ∀ (A : TR ι) (n : ℕ), P A 1 n = 0
  P_bot : ∀ i j : ι, P (gv (i, 0)) (gv (j, 0)) = D.ope i j
  P_Tleft : ∀ A B : TR ι, P (Ttr A) B = Lstep (P A B)
  P_Tright : ∀ A B : TR ι, P A (Ttr B) = Rstep (P A B)
  Ttr_N : ∀ A B : TR ι, Ttr (N A B) = N (Ttr A) B + N A (Ttr B)
  quasiAssoc : ∀ (p : GenIdx ι) (w : FreeMonoid (GenIdx ι)) (C : TR ι) (m : ℕ),
    (∀ n, m ≤ n → P (wbasis ι w) C n = 0) →
    (∀ n, m ≤ n → P (gv p) C n = 0) →
    N (gv p * wbasis ι w) C
      = N (gv p) (N (wbasis ι w) C)
        + ∑ n ∈ Finset.range m, N (Tdiv (n + 1) (gv p)) (P (wbasis ι w) C n)
        + koszul (D.par p.1) (parWord D w) •
            ∑ n ∈ Finset.range m, N (Tdiv (n + 1) (wbasis ι w)) (P (gv p) C n)
  wickLeft : ∀ (p q : GenIdx ι) (C : TR ι) (n : ℕ),
    P (gv p) (gv q * C) n
      = N (P (gv p) (gv q) n) C
        + koszul (D.par p.1) (D.par q.1) • N (gv q) (P (gv p) C n)
        + ∑ k ∈ Finset.range n,
            ((n.choose k : ℕ) : ℂ) • P (P (gv p) (gv q) k) C (n - 1 - k)
  wickRight : ∀ (p : GenIdx ι) (w : FreeMonoid (GenIdx ι)) (C : TR ι) (n m : ℕ),
    (∀ k, m ≤ k → P (wbasis ι w) C k = 0) →
    (∀ k, m ≤ k → P (gv p) C k = 0) →
    P (gv p * wbasis ι w) C n
      = ∑ j ∈ Finset.range m, N (Tdiv j (gv p)) (P (wbasis ι w) C (n + j))
        + koszul (D.par p.1) (parWord D w) •
            (∑ j ∈ Finset.range m, N (Tdiv j (wbasis ι w)) (P (gv p) C (n + j))
              + ∑ j ∈ Finset.range n, P (wbasis ι w) (P (gv p) C j) (n - 1 - j))
  N_zfilt : ∀ {dA dB : ℤ} {A B : TR ι}, A ∈ zfilt D.wt dA → B ∈ zfilt D.wt dB →
    N A B ∈ zfilt D.wt (dA + dB)
  P_zfilt : ∀ {dA dB : ℤ} {A B : TR ι}, A ∈ zfilt D.wt dA → B ∈ zfilt D.wt dB →
    ∀ n : ℕ, P A B n ∈ zfilt D.wt (dA + dB - 1)
  P_filt : ∀ {dA dB : ℤ} {A B : TR ι}, A ∈ filt D.wt dA → B ∈ filt D.wt dB →
    ∀ n : ℕ, P A B n ∈ filt D.wt (dA + dB - n - 1)
\end{leancode}

\begin{dictionary}
\lean{TR ι →ₗ[ℂ] TR ι →ₗ[ℂ] TR ι} is a bilinear map; \lean{P A B n} is the coefficient of $\lambda^n/n!$ in
$P_\lambda(A,B)$. \lean{quasiAssoc}, \lean{wickLeft} and \lean{wickRight} transcribe (3.3), (3.6), (3.7) with finite cutoffs \lean{m}
justified by the vanishing hypotheses. The field docstring of \lean{quasiAssoc} records a deliberate sign choice: the Koszul
sign is $p(a,B)$ where De Sole--Kac's printed (3.3) has $p(a,b)$; the authors argue the printed sign is a misprint. \lean{Tdiv k} is $T^k/k!$,
\lean{parWord D w} the parity of a word. The library derives \lean{P_gv_eq_brk} (on letters, $P$ is the bracket of
Definition~\ref{l11:def:brk}) and \lean{P_eventually_zero} from these fields. \textbf{Fidelity.} 16 fields transcribe De
Sole--Kac Lemma~3.2 / Lemma~4.2; \lean{P_filt} is an addition (needed to make slices terminate). De Sole--Kac's
Definition~3.1 also assumes the seed bracket is parity-preserving; \lean{OPEData} does not, so the library names this as
\lean{OPEData.}\allowbreak\lean{SeedParityCompatible} and proves it propagates to all letters (\lean{brk_mem_parPiece}). Exact-weight
gradedness of $N,P$ is a separate mixin \lean{NPGraded}.
\end{dictionary}

\begin{physics}
$N(A,B)$ is the normal-ordered product of two composite operators and $P(A,B)$ their full singular OPE; the Wick formulas are the
rules for computing OPEs of composites (the ``generalised Wick theorem'' of Bais et al.\ used to compute $W$-algebra OPEs), and
quasi-associativity is the rule for rebracketing nested normal-ordered products. Their existence and uniqueness for a non-linear
Lie conformal algebra is the content of the De Sole--Kac construction.
\end{physics}

\begin{definition}[The full $M(R)$ over the interface]\label{l11:def:MRfull}
Given \lean{NP}, the full Definition~3.4 generator is $\mathrm{relD}(p,q;B) = pqB - (-1)^{p(p)p(q)}qpB - N(\int_{-T}^0[p_\lambda q],B)$,
and $M_{\mathrm{full}}$ is the span of all $A\cdot\mathrm{relD}(p,q;w)$ over words $w$.
\end{definition}
\noindent File \leanpath{VOA/\allowbreak Core/\allowbreak NonlinearMR.lean}.
\begin{leancode}
noncomputable def relDF (p q : GenIdx ι) : TR ι →ₗ[ℂ] TR ι where
  toFun B := gv p * gv q * B - koszul (D.par p.1) (D.par q.1) • (gv q * gv p * B)
    - NP.N (opeInt D p q) B

noncomputable def MRfullCar : Submodule ℂ (TR ι) :=
  Submodule.span ℂ (Set.range
    (fun z : TR ι × GenIdx ι × GenIdx ι × FreeMonoid (GenIdx ι) =>
      z.1 * relDF NP z.2.1 z.2.2.1 (wbasis ι z.2.2.2)))

def OPEData.IsLinear (D : OPEData ι) : Prop := ∀ i j (n : ℕ), D.ope i j n ∈ linSub ι

theorem MRfullCar_eq_MRtwoCar (hD : D.IsLinear) (NP : NPExtension D) :
    MRfullCar NP = MRtwoCar D := by

theorem MRfull_virasoro (c : ℂ) (NP : NPExtension (virasoroOPE c)) :
    MRfullCar NP = MRtwoCar (virasoroOPE c) :=

theorem vacU_ne_zero_MRfull (hC : D.IsConformal) (hg : NPGraded NP) :
    vacU (MRfull NP).toMIdeal ≠ 0 :=
\end{leancode}

\begin{dictionary}
This is Definition~3.4 verbatim (with $D$ ranging over the word basis, which suffices by linearity). For a linear datum
(\lean{IsLinear}: every OPE value in \lean{linSub ι} $=\C1\oplus\operatorname{span}\{\text{letters}\}$) the full ideal is
$M_{\mathrm{two}}$, for any interface term. \lean{vacU_ne_zero_MRfull} is the graded non-vacuity guard, with hypotheses
\lean{IsConformal} and \lean{NPGraded NP}. The closure of $M(R)$ under $N$ and $P$ on both sides (De Sole--Kac Cor.~4.5, which
uses the Jacobi identity) is only stated, as the \lean{Prop}-structure \lean{NPClosure}. The docstring notes that the composite
$M_{\mathrm{full}} = \ker\Phi$ for Virasoro is not formalised (it would be one \lean{Eq.trans} in an \leanpath{Examples} file).
\end{dictionary}

\textbf{Proof on paper.} For $X\in\C1\oplus R$, $N(X,B) = XB$ by the two unit/letter axioms, and for linear data
$\int_{-T}^0[p_\lambda q]\in\C1\oplus R$, so $\mathrm{relD}(p,q;B) = \mathrm{rel}(p,q)B$. The augmentation kills positive-weight words. \qed

\begin{status}
Proved, \emph{conditional on an \lean{NPExtension} term, which does not exist}; \texttt{\#print axioms} of
\lean{MRfullCar_eq_MRtwoCar}, \lean{MRfull_virasoro}, \lean{vacU_ne_zero_MRfull}: standard three. \lean{virasoroOPE_isLinear} is
unconditional.
\end{status}

\begin{physics}
For a Lie-type OPE nothing nonlinear happens, and the general construction must reduce to the universal enveloping algebra of the
mode Lie algebra; that consistency check is what is proved, modulo the existence of the general construction.
\end{physics}

\begin{definition}[Presentations with strong relations]\label{l11:def:vapres}
A \lean{VAPresentation} is an OPE datum together with a finite list of homogeneous elements of $\mathcal T(R)$ of positive weight, the
\emph{strong relations} (null fields). Given \lean{NP}, the relation ideal is the smallest subspace containing them and closed under
$T$ and under $N$ and every slice of $P$ on both sides; the presentation ideal is $M_{\mathrm{full}}+\mathrm{RelIdeal}$ and the
presented object is $\mathcal T(R)$ modulo it.
\end{definition}
\noindent File \leanpath{VOA/\allowbreak Core/\allowbreak VAPresentation.lean}.
\begin{leancode}
structure VAPresentation (ι : Type u) extends OPEData ι where
  relIdx : Type
  relFin : Fintype relIdx
  relVal : relIdx → TR ι
  relWt : relIdx → ℕ
  relWt_pos : ∀ r, 0 < relWt r
  relVal_homog : ∀ r, relVal r ∈ gr wt (relWt r)

inductive InRelIdeal (P : VAPresentation ι) (NP : NPExtension P.toOPEData) : TR ι → Prop
  | of (r : P.relIdx) : InRelIdeal P NP (P.relVal r)
  | zero : InRelIdeal P NP 0
  | add {x y : TR ι} : InRelIdeal P NP x → InRelIdeal P NP y → InRelIdeal P NP (x + y)
  | smul (c : ℂ) {x : TR ι} : InRelIdeal P NP x → InRelIdeal P NP (c • x)
  | tstab {x : TR ι} : InRelIdeal P NP x → InRelIdeal P NP (Ttr x)
  | nLeft (A : TR ι) {x : TR ι} : InRelIdeal P NP x → InRelIdeal P NP (NP.N A x)
  | nRight (A : TR ι) {x : TR ι} : InRelIdeal P NP x → InRelIdeal P NP (NP.N x A)
  | pLeft (A : TR ι) (n : ℕ) {x : TR ι} : InRelIdeal P NP x → InRelIdeal P NP (NP.P A x n)
  | pRight (A : TR ι) (n : ℕ) {x : TR ι} : InRelIdeal P NP x → InRelIdeal P NP (NP.P x A n)

noncomputable def MPresCar : Submodule ℂ (TR ι) := MRfullCar NP ⊔ RelIdeal P NP

noncomputable abbrev UPres : Type u := UR (MPres P NP).toMIdeal
\end{leancode}

\begin{dictionary}
\lean{relVal r} is the $r$-th relation as \emph{literal data} (an explicit element of $\mathcal T(R)$, never a kernel); \lean{relWt},
\lean{relVal_homog} its weight and homogeneity. \lean{relFin} (finiteness) is recorded but used by no theorem. \lean{InRelIdeal} is an
inductive predicate: the smallest set closed under the nine listed constructors. This is Linshaw's $M'(\mathcal A)$-closure for
degenerate non-linear conformal algebras.
\end{dictionary}

\begin{physics}
This is how a $W$-algebra or a quotient such as a minimal model is specified by physicists: generators, their OPEs, and the null
fields that must be set to zero (``the null field $N = {:}TT{:}-\frac3{10}\partial^2T$ vanishes'' for Lee--Yang). The relation ideal is
everything those null fields generate under OPE with anything.
\end{physics}

\begin{theorem}[L2 recovery; the vacuum guard and its counterexample]\label{l11:thm:vapres}
For every \lean{NP}:
\begin{enumerate}
\item with no relations, the presentation ideal is $M_{\mathrm{full}}$, and for linear data $M_{\mathrm{two}}$;
\item if the datum is conformal, \lean{NPGraded NP}, and \lean{RelBracketVacuumFree} (every bracket slice of an element of the relation
  ideal has no weight-$0$ component), then $|0\rangle\neq0$ in the presented object;
\item for the Virasoro datum at $c\neq0$ with the single relation $\omega$ (i.e.\ $L$), $1\in\mathrm{RelIdeal}$.
\end{enumerate}
\end{theorem}
\noindent File \leanpath{VOA/\allowbreak Core/\allowbreak VAPresentation.lean}.
\begin{leancode}
theorem MPres_empty_linear (h : IsEmpty P.relIdx) (hD : P.toOPEData.IsLinear) :
    MPresCar P NP = MRtwoCar P.toOPEData := by

def RelBracketVacuumFree : Prop :=
  ∀ (A x : TR ι) (n : ℕ), InRelIdeal P NP x →
    NP.P A x n ∈ grPos P.wt ∧ NP.P x A n ∈ grPos P.wt

theorem vacU_ne_zero_MPres (hC : P.toOPEData.IsConformal) (hg : NPGraded NP)
    (hvf : RelBracketVacuumFree P NP) : vacU (MPres P NP).toMIdeal ≠ 0 :=

theorem one_mem_relIdeal_omegaRel (c : ℂ) (hc : c ≠ 0)
    (NP : NPExtension (virasoroOPE c)) :
    (1 : TR Unit) ∈ RelIdeal (omegaRelPresentation c) NP := by
\end{leancode}

\begin{dictionary}
In (3) the presentation \lean{omegaRelPresentation c} lists $L$ itself as a relation; then $P(L,L)(3) = \frac c2\cdot1$ is in the
relation ideal, so the vacuum dies. This refutes the design's ``model-free'' guard (conformal $+$ homogeneous positive-weight
relations $\Rightarrow$ $1\notin M$), and shows that \lean{RelBracketVacuumFree} is not implied by \lean{IsConformal} and \lean{NPGraded}
(\lean{not_}\allowbreak\lean{relBracketVacuumFree_}\allowbreak\lean{omegaRel}). \textbf{But} all four statements take \lean{NP} as input, so --- as the docstring stresses
--- as Lean propositions they ``refute nothing by themselves'' until an interface term exists. The mathematics behind (3) is
unconditional: setting $T(z) = 0$ at $c\neq0$ contradicts $T(z)T(w)\sim\frac{c/2}{(z-w)^4}$.
\end{dictionary}

\textbf{Proof on paper.} (1) The inductive closure of the empty set is $\{0\}$. (2) Under the hypotheses, every generator of the
presentation ideal lies in $\mathrm{grPos}$, which the augmentation kills. (3) \lean{pLeft} applied to the relation $L$ with $A = L$,
$n=3$, and \lean{P_bot}, give $\frac c2\cdot1$; scale. \qed

\begin{status}
Proved, conditional on \lean{NP : NPExtension …} (uninhabited); \texttt{\#print axioms}: standard three. A recorded dead end (the
model-free vacuum guard), refuted on paper unconditionally and in Lean conditionally.
\end{status}

\begin{physics}
Imposing a null relation is safe only if the relation really is null in some consistent theory; an arbitrary positive-weight
relation can generate the identity through its OPE. The stress tensor is the cleanest example: it is never null when $c\neq0$.
\end{physics}

\begin{definition}[Evaluation, nullity certificates, \lean{TheVOA}, realizations]\label{l11:def:doctrine}
For a VOA $W$ and $g:\iota\to W$, $\mathrm{ev}:\mathcal T(R)\to W$ sends a word $(i_1,k_1)\cdots$ to
$(T^{k_1}g_{i_1}/k_1!)_{(-1)}\cdots\lelevenvac$. A relation is \emph{null} for an invariant form $B$ if its evaluation pairs to zero with
everything. \lean{TheVOA P NP W} is the simple quotient of a given VOA structure $W$ on the presented space. A \emph{realization} of a
presentation in $W$ is $g:\iota\to W$ with homogeneous images of the declared weights, satisfying the OPE list, killing the
relations, and with spanning word states. \lean{simple_}\allowbreak\lean{realization_}\allowbreak\lean{unique} is the proposition that any two simple realizations of a
presentation are isomorphic as vertex algebras.
\end{definition}
\noindent File \leanpath{VOA/\allowbreak Core/\allowbreak PresentationDoctrine.lean}.
\begin{leancode}
noncomputable def evTR (W : VertexOperatorAlgebra 𝒱) (g : ι → 𝒱) : TR ι →ₗ[ℂ] 𝒱 :=
  (wbasis ι).constr ℂ (fun w => evWord W g (FreeMonoid.toList w))

def NullCertificate (P : VAPresentation ι) (W : VertexOperatorAlgebra 𝒱)
    (B : InvariantForm W) (g : ι → 𝒱) (r : P.relIdx) : Prop :=
  ∀ v : 𝒱, B.form (evTR W g (P.relVal r)) v = 0

theorem mem_maxGradedIdeal_of_null {P : VAPresentation ι} {W : VertexOperatorAlgebra 𝒱}
    {B : InvariantForm W} {g : ι → 𝒱} {r : P.relIdx}
    (hB : B.IsFullyInvariant) (hn : B.Normalized)
    (hnull : NullCertificate P W B g r) :
    evTR W g (P.relVal r) ∈ QuotVOA.maxGradedIdeal W := by

noncomputable def TheVOA (P : VAPresentation ι) (NP : NPExtension P.toOPEData)
    (W : VertexOperatorAlgebra (UPres P NP)) :
    VertexOperatorAlgebra (UPres P NP ⧸ QuotVOA.maxGradedIdeal W) :=
  QuotVOA.maxQuotientVOA W

structure RealizesPresentation (P : VAPresentation ι) (W : VertexOperatorAlgebra 𝒱) where
  gen : ι → 𝒱
  gen_homog : ∀ i, gen i ∈ W.grading (P.wt i : ℤ)
  hOPE : ∀ i j (n : ℕ), W.smul (gen i) (n : ℤ) (gen j) = evTR W gen (P.ope i j n)
  hRel : ∀ r, evTR W gen (P.relVal r) = 0
  spanning : Submodule.span ℂ
    (Set.range fun w : FreeMonoid (GenIdx ι) => evTR W gen (wbasis ι w)) = ⊤

def simple_realization_unique : Prop :=
  ∀ {κ : Type u} (P : VAPresentation κ) {𝒱₁ 𝒱₂ : Type v}
    [AddCommGroup 𝒱₁] [Module ℂ 𝒱₁] [AddCommGroup 𝒱₂] [Module ℂ 𝒱₂]
    (W₁ : VertexOperatorAlgebra 𝒱₁) (W₂ : VertexOperatorAlgebra 𝒱₂),
    RealizesPresentation P W₁ → RealizesPresentation P W₂ →
    W₁.IsSimple → W₂.IsSimple →
    ∃ e : 𝒱₁ ≃ₗ[ℂ] 𝒱₂,
      e W₁.vacuum = W₂.vacuum ∧
      ∀ (u v : 𝒱₁) (n : ℤ), e (W₁.smul u n v) = W₂.smul (e u) n (e v)
\end{leancode}

\begin{dictionary}
\lean{InvariantForm W}, \lean{IsFullyInvariant}, \lean{Normalized} and Li's radical theorem are Chapter~\ref{ch:L04}
(Definitions~\ref{l04:def:invform}, \ref{l04:def:fullinv}, Theorem~\ref{l04:thm:lirad}); \lean{QuotVOA.}\allowbreak\lean{maxGradedIdeal W} is the
maximal proper graded ideal and \lean{maxQuotientVOA} the simple quotient (Definition~\ref{l04:def:maxideal}). \lean{evTR} is defined
by \lean{Module.Basis.constr} on the word basis. \lean{mem_}\allowbreak\lean{maxGradedIdeal_}\allowbreak\lean{of_}\allowbreak\lean{null} is a genuine theorem: a null relation evaluates into
the maximal ideal (Li's radical theorem).
\end{dictionary}

\begin{warning}
\textbf{\lean{TheVOA} is not tied to its presentation (moderate), and the uniqueness statement is unproved (open gap).} \lean{TheVOA P NP W}
takes as input an \emph{arbitrary} \lean{W : VertexOperatorAlgebra (UPres P NP)}: any VOA structure whatsoever on the underlying vector
space of the presented object, with no requirement that its vertex operators extend the creation half, satisfy the OPE table, or kill
the relations. Its value is \lean{QuotVOA.}\allowbreak\lean{maxQuotientVOA W}, which depends on \lean{P} and \lean{NP} only through the carrier type. So
the name ``THE vertex algebra of a presentation'' is a convention attached to a datum the library cannot yet produce (the docstring says
the structure ``arrives via the \lean{GenFields} contract and conformal data in later stages''), not a construction. And since
\lean{UPres P NP} needs \lean{NP}, \lean{TheVOA} has no instance at all. Separately, \lean{simple_}\allowbreak\lean{realization_}\allowbreak\lean{unique} is a
\lean{def : Prop} --- a statement, \emph{not} a theorem --- and nothing proves it; its intended proof route (universal property plus Li's
theorem, after Roitman) is described only in prose. Note also that \lean{RealizesPresentation} mentions only the non-negative products
(\lean{hOPE} for \lean{n : ℕ}) and that the conclusion of \lean{simple_}\allowbreak\lean{realization_}\allowbreak\lean{unique} preserves vacuum and all products but not the
conformal vector.
\end{warning}

\textbf{Proof on paper} (of \lean{mem_}\allowbreak\lean{maxGradedIdeal_}\allowbreak\lean{of_}\allowbreak\lean{null}). For a fully invariant normalised form, the radical is the maximal proper
graded ideal (Li); a null vector lies in the radical. \qed

\begin{status}
\lean{mem_}\allowbreak\lean{maxGradedIdeal_}\allowbreak\lean{of_}\allowbreak\lean{null}, \lean{relIdealNull_}\allowbreak\lean{iff_}\allowbreak\lean{forms_}\allowbreak\lean{vanish}: proved; standard three. \lean{TheVOA}: definition, no
instance. \lean{simple_}\allowbreak\lean{realization_}\allowbreak\lean{unique}: unproved proposition, no consumer.
\end{status}

\begin{physics}
Nullity is the statement ``this composite operator has zero two-point function with everything'', which in a unitary or at least
nondegenerate theory means it vanishes. ``The'' chiral algebra of a presentation is the simple quotient: the one in which every null
field has been set to zero. Uniqueness of simple realizations is the mathematical form of ``a chiral algebra is determined by its
generators, their OPEs and its null relations'' --- believed by every physicist, and proved for lattice VOAs by Roitman, but not
formalised here.
\end{physics}

%======================================================================
\section{Li's \texorpdfstring{$\Delta$}{Delta}-operator on the Fock space}\label{l11:sec:lidelta}
%======================================================================

Two files filed with this chapter concern a different Li construction: the operator $\Delta(h,z)$ of Li (1997) that produces new
(twisted, or charged) modules from old by $Y_\Delta(u,z) = Y(\Delta(h,z)u,z)$. On the rank-one Fock space $\C[X_0,X_1,\dots]$ of
Chapter~\ref{ch:L05} (Definition~\ref{l05:def:fock}), where $a_n = nk\,\partial/\partial X_{n-1}$ for $n\ge1$ and $a_0 = 0$,
$\Delta(\lambda,z) = z^{\lambda a_0}\exp\bigl(-\sum_{n\ge1}\frac\lambda na_n(-z)^{-n}\bigr)$ is the substitution
$X_i\mapsto X_i+\lambda k(-1)^iw^{i+1}$, $w = z^{-1}$.

\begin{definition}[Li's $\Delta$ as a substitution, one and two variables]\label{l11:def:lidelta}
$\Delta_\lambda:\C[X_i]\to\C[w][X_i]$ is the algebra map $X_i\mapsto X_i+\lambda k(-1)^iw^{i+1}$. The two-variable version
$\Delta_\lambda^{(2)}:\C[X_i]\to(\C[w][[z_0]])[X_i]$ substitutes the $\iota_{z,z_0}$-expansion of $\lambda k(-1)^i(z+z_0)^{-i-1}$.
\end{definition}
\noindent Files \leanpath{VOA/\allowbreak Examples/\allowbreak LiDelta.lean}, \leanpath{VOA/\allowbreak Examples/\allowbreak LiDelta2.lean}.
\begin{leancode}
noncomputable def deltaCoeff (k lam : ℂ) (i : ℕ) : Polynomial ℂ :=
  Polynomial.C (lam * k * (-1) ^ i) * Polynomial.X ^ (i + 1)

noncomputable def deltaSub (k lam : ℂ) : Fock →ₐ[ℂ] FockW :=
  aeval (fun i : ℕ => (X i : FockW) + MvPolynomial.C (deltaCoeff k lam i))

noncomputable def expandNegPow (i : ℕ) : CoeffW₂ :=
  PowerSeries.mk fun j =>
    Polynomial.C (gBinom (-(i : ℤ) - 1) j) * Polynomial.X ^ (i + 1 + j)

noncomputable def deltaSub₂ (k lam : ℂ) : Fock →ₐ[ℂ] FockW₂ :=
  aeval fun i : ℕ => (X i : FockW₂) + MvPolynomial.C (deltaCoeff₂ k lam i)
\end{leancode}

\begin{dictionary}
\lean{Polynomial.X} is $w$; \lean{gBinom (-(i : ℤ) - 1) j} is $\binom{-i-1}j$, so \lean{expandNegPow i} is
$\sum_j\binom{-i-1}jw^{i+1+j}z_0^j = (z+z_0)^{-i-1}$ expanded in non-negative powers of $z_0$. The exponent check:
$-\sum_n\frac\lambda nnk\,\partial_{X_{n-1}}(-1)^nz^{-n}$ translates $X_{n-1}$ by $\lambda k(-1)^{n-1}w^n$, i.e.\ $X_i$ by
$\lambda k(-1)^iw^{i+1}$, matching \lean{deltaCoeff}. Faithful.
\end{dictionary}

\begin{physics}
The substitution $X_i\mapsto X_i+\text{const}\cdot w^{i+1}$ is the effect of conjugating by a vertex operator of momentum $\lambda$:
each oscillator acquires a classical background $\langle\partial^i\phi\rangle\propto\lambda z^{-i-1}$, the field of a charge at the origin.
\end{physics}


\begin{theorem}[Group law, inverse, specialisation]\label{l11:thm:lidelta}
$\Delta_\mu\circ\Delta_\lambda = \Delta_{\lambda+\mu}$ (as substitutions), $\Delta_0$ is the inclusion, $\Delta_{-\lambda}\circ\Delta_\lambda$ is the
inclusion, $\Delta_\lambda$ commutes with $\partial/\partial X_j$, and setting $z_0 = 0$ in $\Delta^{(2)}_\lambda$ gives $\Delta_\lambda$.
\end{theorem}
\noindent Files as above.
\begin{leancode}
theorem deltaSubW_comp_deltaSub (k lam mu : ℂ) :
    ((deltaSubW k mu).restrictScalars ℂ).comp (deltaSub k lam) = deltaSub k (lam + mu) := by

theorem deltaSubW_neg_comp_deltaSub (k lam : ℂ) :
    ((deltaSubW k (-lam)).restrictScalars ℂ).comp (deltaSub k lam) = includeW := by

theorem deltaSub_pderiv (k lam : ℂ) (j : ℕ) (p : Fock) :
    deltaSub k lam (MvPolynomial.pderiv j p)
      = MvPolynomial.pderiv j (deltaSub k lam p) := by

theorem deltaSub₂_specialize (k lam : ℂ) (p : Fock) :
    specializeZ₀ (deltaSub₂ k lam p) = deltaSub k lam p := by
\end{leancode}

\begin{dictionary}
\lean{deltaSubW k μ} is the $\C[w]$-linear extension of $\Delta_\mu$ to $\C[w][X_i]$; \lean{restrictScalars ℂ} views it as $\C$-linear.
\lean{includeW} is the inclusion $\C[X]\to\C[w][X]$. \lean{specializeZ₀} sets $z_0 = 0$ in the coefficients.
\end{dictionary}

\textbf{Proof on paper.} Translations compose additively; partial derivatives commute with constant translations; $z_0 = 0$ in
$(z+z_0)^{-i-1}$ gives $z^{-i-1}$. \qed

\begin{status}
Proved, unconditional; \texttt{\#print axioms} of \lean{VOA.}\allowbreak\lean{Heisenberg.}\allowbreak\lean{deltaSub₂_}\allowbreak\lean{specialize},
\lean{VOA.}\allowbreak\lean{Heisenberg.}\allowbreak\lean{deltaSubW_}\allowbreak\lean{comp_}\allowbreak\lean{deltaSub}: standard three. \textbf{Definitional layer only}: Li's conjugation identity
$\Delta(\lambda,z)Y(u,z_0)\Delta(\lambda,z)^{-1} = Y(\Delta(\lambda,z+z_0)u,z_0)$ and the resulting charged module $F_\lambda$ are not proved
here (the file says so); consumers are \leanpath{FockBC2.lean}, \leanpath{FockBaseChange.lean} (Chapter~\ref{ch:L05}).
\end{status}

\begin{physics}
$\Delta(\lambda,z)$ is spectral flow by the $U(1)$ current: it shifts the zero mode $a_0$ by $\lambda k$ and turns the vacuum Fock module
into the charged Fock module $F_\lambda$, the sector of momentum $\lambda$ of the free boson. The group law is the additivity of
momentum. It is the standard route to the first logarithmic modules (by differentiating in $\lambda$).
\end{physics}


%======================================================================
\section{Referee's checklist}\label{l11:sec:checklist}
%======================================================================

\begingroup\small
\begin{longtable}{p{0.27\textwidth}p{0.27\textwidth}p{0.22\textwidth}p{0.14\textwidth}}
\toprule
Lean declaration & Intended meaning & Hypotheses / status & Fidelity check\\
\midrule\endhead
\lean{GradedFlag}, \lean{GoodFiltration} & filtration of weight spaces; mode-compatible & definitions; \textbf{every VOA has one} (checked outside the library); weaker than Li--Arakawa ``good'' & Def.~\ref{l11:def:flag}\\
\lean{GoodFiltration.}\allowbreak\lean{quotient} & good filtrations descend to quotients & graded ideal; unconditional & Thm.~\ref{l11:thm:flagquot}\\
\lean{LinearFiltration}, \lean{IsStrictFilteredAlong} & filtered / strict maps & definitions & Def.~\ref{l11:def:linfilt}\\
\lean{IsFilteredAlong.}\allowbreak\lean{isStrictFilteredAlong_}\allowbreak\lean{of_}\allowbreak\lean{finrank_}\allowbreak\lean{eq_}\allowbreak\lean{range_}\allowbreak\lean{inf} & dimension criterion for strictness & injective, finite pieces; unconditional & Thm.~\ref{l11:thm:strict}\\
\lean{liFiltration}, \lean{Gfil} & Li's decreasing filtration & definition; same as Li 2005 & Def.~\ref{l11:def:gfil}\\
\lean{sum_}\allowbreak\lean{liWeightLayer_}\allowbreak\lean{finrank_}\allowbreak\lean{eq_}\allowbreak\lean{grading} & $\sum_p\dim\lelevengr_{p,N} = \dim V_N$ & none; unconditional & Thm.~\ref{l11:thm:litele}\\
\lean{adjointGfil_smul_mem}, \lean{liFiltration_}\allowbreak\lean{smul_}\allowbreak\lean{nonneg_}\allowbreak\lean{mem} & Li's estimates $F^p_{(n)}F^q$ & unconditional & Thm.~\ref{l11:thm:limult}\\
\lean{liWeightLayer}\allowbreak\lean{GCommRing}, \lean{liAssociatedGraded}\allowbreak\lean{Translation_leibniz} & $\lelevengr^FV$ commutative differential algebra & unconditional; Poisson bracket not packaged here (see \lean{liPoissonVertexAlgebra}, Ch.~\ref{ch:L20}) & Thm.~\ref{l11:thm:ligr}\\
\lean{O₂_le_liFiltration_one}, \lean{liFiltration_one_eq_O₂} & $F^1 = C_2$, $\lelevengr_0 = \lelevenR$ & none; equality proved in library (\leanpath{LiC2Bridge.lean}) & Thm.~\ref{l11:thm:f1}\\
\lean{Presentation.}\allowbreak\lean{IsExactPresentation}, \lean{quotientEquiv} & exact presentations; first iso.\ theorem & surjective; unconditional & Def.~\ref{l11:def:compres}\\
\lean{FiniteJet.jetIdeal}, \lean{PresentedJet} & order-$m$ jet scheme of $\C[x]/I$ & definitions; standard & Def.~\ref{l11:def:jets}\\
\lean{jetIdeal_}\allowbreak\lean{span_}\allowbreak\lean{eq_}\allowbreak\lean{jetIdealOfSet}, \lean{jetWeightPieceOrderEquiv} & jets: generators, stabilisation & positive weights; unconditional & Thm.~\ref{l11:thm:jetweights}\\
\lean{liPresentedJetMap} & canonical $J_m\lelevenR\to\lelevengr^FV$ & relations vanish on $x$ & Def.~\ref{l11:def:jetmap}\\
\lean{liPresentedJetMap_}\allowbreak\lean{injective_}\allowbreak\lean{iff_}\allowbreak\lean{polynomial_}\allowbreak\lean{kernel} & kernel criterion & unconditional & Thm.~\ref{l11:thm:jetmap}\\
\lean{liPresentedJetMap_}\allowbreak\lean{hits_}\allowbreak\lean{homogeneous_}\allowbreak\lean{class_}\allowbreak\lean{of_}\allowbreak\lean{generators} & jet map hits every Li layer, $p\le m$ & \lean{hdec}: \textbf{forces $C_2$-cofinite} & Thm.~\ref{l11:thm:hits}\\
\lean{liFiniteJetPolynomialMap_}\allowbreak\lean{kernel_}\allowbreak\lean{eq_}\allowbreak\lean{of_}\allowbreak\lean{bidegree_}\allowbreak\lean{hilbert_}\allowbreak\lean{squeeze} & classical freeness from a squeeze & three hypotheses (the content) & Thm.~\ref{l11:thm:freeness}\\
\lean{FiniteHilbertSqueeze.}\allowbreak\lean{pointwise_}\allowbreak\lean{eq_}\allowbreak\lean{of_}\allowbreak\lean{sandwich_}\allowbreak\lean{of_}\allowbreak\lean{sum_}\allowbreak\lean{eq} & dimension squeeze & unconditional & Thm.~\ref{l11:thm:squeeze}\\
\lean{AffineVOA.}\allowbreak\lean{CurrentLi.}\allowbreak\lean{liFiltration_}\allowbreak\lean{eq_}\allowbreak\lean{currentWordPiece} & currents span Li filtration of $V^k(\mathfrak{sl}_2)$ & $k\neq-2$; unconditional & Thm.~\ref{l11:thm:affineli}\\
\lean{GeneratorWords}\allowbreak\lean{SpanLiFiltration}, \lean{genWordPiece}, \lean{liFiltration_}\allowbreak\lean{eq_}\allowbreak\lean{genWordPiece} & $F^pV\subseteq\mathcal W_p$ (Li spanning; $\mathcal W_p$ the word piece); then $F^p=\mathcal W_p$ & definition; \lean{wt} unused until \lean{hgen} & Def.~\ref{l11:def:genwordpiece}, Prop.~\ref{l11:prop:wordpieceli}\\
\lean{liInfinite}\allowbreak\lean{PresentedJetMap_}\allowbreak\lean{surjective_}\allowbreak\lean{of_}\allowbreak\lean{genWordSpan}, \lean{liPresentedJet}\allowbreak\lean{BidegreeMap_}\allowbreak\lean{surjective_}\allowbreak\lean{of_}\allowbreak\lean{genWordSpan} & word spanning $\Rightarrow$ $J_\infty$ and bidegree jet maps onto & \lean{hgen}, \lean{hspan} (= Li spanning), \lean{hI}; unconditional; any VOA, no $C_2$-cofiniteness; comparison with \lean{hdec} on paper only (Rem.~\ref{l11:rem:hdechspan}) & Thm.~\ref{l11:thm:genwordsurj}\\
\lean{generatorWords}\allowbreak\lean{SpanLiFiltration_}\allowbreak\lean{of_}\allowbreak\lean{tracked} & word spanning $+$ tracked estimate $\Rightarrow$ $F^p\subseteq\mathcal W_p$ & \lean{htrack}, \lean{hspan} (here = word spanning, $\mathcal W_0=V$); unconditional & Thm.~\ref{l11:thm:tracked}\\
\lean{generatorMode}\allowbreak\lean{Tracked_}\allowbreak\lean{of_}\allowbreak\lean{current_}\allowbreak\lean{comm} & Lie-type currents obey the tracked estimate & commutator shape, vacuum; no Jacobi needed; unconditional & Thm.~\ref{l11:thm:liecurrents}\\
\lean{generatorWords}\allowbreak\lean{SpanLiFiltration_}\allowbreak\lean{quotient} & spanning descends to graded quotients & graded ideal; unconditional & Thm.~\ref{l11:thm:wordquot}\\
\lean{AffineVOA.}\allowbreak\lean{CurrentLi.}\allowbreak\lean{sl2_genWordsSpan}, \lean{sl2_}\allowbreak\lean{liFiltration_}\allowbreak\lean{eq_}\allowbreak\lean{genWordPiece} & generic theorem recovers Thm.~\ref{l11:thm:affineli} & $k\neq-2$; unconditional; consistency check & Thm.~\ref{l11:thm:sl2genwords}\\
\lean{hdec}-based criteria (\leanpath{LiGeneratorWords}, \leanpath{LiJetBidegree}) & surjectivity from span modulo $C_2(V)$ & \textbf{\lean{hdec} implies \lean{C2Cofinite}: unusable for quasi-lisse and Deligne cases} & warning in §\ref{l11:sec:genwords}\\
\lean{PolynomialMinimum.}\allowbreak\lean{linearIndependent_}\allowbreak\lean{of_}\allowbreak\lean{symbols}, \lean{filteredOrderedWord_}\allowbreak\lean{independent} & independent symbols $\Rightarrow$ independent states & unconditional & Thms.~\ref{l11:thm:symbols}, \ref{l11:thm:filteredwords}\\
\lean{PolynomialMatrix.}\allowbreak\lean{kernel_}\allowbreak\lean{pair_}\allowbreak\lean{of_}\allowbreak\lean{simple_}\allowbreak\lean{root} & kernel pair at a simple determinant root & domain; unconditional & Thm.~\ref{l11:thm:simpleroot}\\
\lean{SparseRational.}\allowbreak\lean{evalSparse_}\allowbreak\lean{normalizeQ} etc. & certified sparse arithmetic & unconditional & Def.~\ref{l11:def:sparse}\\
\lean{NonlinearConformal.}\allowbreak\lean{OPEData}, \lean{virasoroOPE} & De Sole--Kac Def.~3.1 (no Jacobi) & definition & Def.~\ref{l11:def:opedata}\\
\lean{brk_mem_filt}, \lean{Ttr_rel} & bracket grading; $T$-stability of relations & unconditional & Thms.~\ref{l11:thm:brk}, \ref{l11:thm:Trel}\\
\lean{vacU_ne_zero_bot} & $\mathcal T(R)/M_{\mathrm{bot}}\neq0$ & unconditional & Thm.~\ref{l11:thm:URbasic}\\
\lean{reconstruct_}\allowbreak\lean{of_}\allowbreak\lean{generators} & reconstruction from generating fields & local, covariant, spanning; unconditional & Thm.~\ref{l11:thm:recongen}\\
\lean{UR_YA} & vertex map on $\mathcal T(R)/M$ & \lean{G : GenFields M} & Def.~\ref{l11:def:genfields}\\
\lean{not_}\allowbreak\lean{nonempty_}\allowbreak\lean{genFields_}\allowbreak\lean{MRbot_}\allowbreak\lean{virasoro} & refutes the $M_{\mathrm{bot}}$ route & \textbf{refutation}; unconditional & warning after Def.~\ref{l11:def:genfields}\\
\lean{virasoro_rel} (+ $c$-independence checked outside the library) & Witt relations; $c$-independent & unconditional & Thm.~\ref{l11:thm:virrel}\\
\lean{VirModule.}\allowbreak\lean{MRtwo_}\allowbreak\lean{eq_}\allowbreak\lean{ker_}\allowbreak\lean{PhiVir} & presentation ideal $=\ker\Phi$ & unconditional, \texttt{pbw}-free & Thm.~\ref{l11:thm:virpres}\\
\lean{VirModule.}\allowbreak\lean{virasoro_}\allowbreak\lean{presentation_}\allowbreak\lean{equivalence_}\allowbreak\lean{genuine} & $\mathcal T(R)/M_{\mathrm{two}}\cong V(c,0)$ with fields & unconditional; fields by transport & Thm.~\ref{l11:thm:virpres}\\
\lean{ope_mem_zfilt}, \lean{rel_mem_gr} & gradings on $\mathcal T(R)$ & unconditional (\lean{IsConformal} for the latter) & Thm.~\ref{l11:thm:gradings}\\
\lean{NPExtension} & De Sole--Kac Lemma 3.2 interface & \textbf{no instance exists} & Def.~\ref{l11:def:npext}\\
\lean{MRfullCar_eq_MRtwoCar}, \lean{vacU_ne_zero_MPres}, \lean{one_}\allowbreak\lean{mem_}\allowbreak\lean{relIdeal_}\allowbreak\lean{omegaRel} & full $M(R)$, vacuum guard, $\omega$-counterexample & \textbf{vacuous today} (need \lean{NP}) & Def.~\ref{l11:def:MRfull}, Thm.~\ref{l11:thm:vapres}\\
\lean{mem_}\allowbreak\lean{maxGradedIdeal_}\allowbreak\lean{of_}\allowbreak\lean{null} & null relations lie in the maximal ideal & fully invariant normalised form & Def.~\ref{l11:def:doctrine}\\
\lean{TheVOA}, \lean{simple_}\allowbreak\lean{realization_}\allowbreak\lean{unique} & ``the'' VOA; uniqueness of simple realizations & \textbf{arbitrary input $W$; unproved \lean{Prop}} & warning in §\ref{l11:sec:doctrine}\\
\lean{Heisenberg.deltaSub}, \lean{deltaSub₂_specialize} & Li's $\Delta$ (spectral flow) as substitution & unconditional; definitional layer & Thm.~\ref{l11:thm:lidelta}\\
\bottomrule
\end{longtable}
\endgroup

Entries in \textbf{bold} in the third column are the ones a referee must not miss: a structure with no instance (so theorems
over it are vacuous today), an unproved \lean{Prop}, a refuted route, or a hypothesis narrower than the literature.

All \texttt{\#print axioms} results quoted in this chapter were obtained at an earlier commit (the snapshot
against which the chapter was first written; later re-probes are stated where they were made) by running scratch files that
import the relevant modules, with \texttt{lake env lean -DmaxSynthPendingDepth=3} (never a bare \texttt{lake env lean}); they were
not re-run for this edition, and no full \texttt{lake build} was run for this chapter. Every library declaration checked reports \lean{[propext, Classical.}\allowbreak\lean{choice,
Quot.}\allowbreak\lean{sound]} (one reports only \lean{[propext, Quot.sound]}); none uses \lean{VOA.pbw}. The Lean checks of this chapter made outside the library (the coarse good
filtration, the first theorem without \lean{hLW}, the $c$-independence of $M_{\mathrm{two}}$, and $F^1V = C_2(V)$) compiled with the
same command and the same three axioms; they are not part of this release.

The declarations of Section~\ref{l11:sec:genwords} were checked at an earlier commit on the branch \texttt{deligne-lean}
(the commit at which that section's quotations were first taken), not re-run for this edition. A scratch file
importing \texttt{VOA.\allowbreak Filtered.\allowbreak AffineSl2GeneratorCurrentInstance} was run with the same command; it
printed the axioms of the eleven declarations named in that section's status boxes and compiled the Lean
check of the warning box (\lean{hdec} forces $C_2$-cofiniteness). All report \lean{[propext, Classical.}\allowbreak\lean{choice,
Quot.}\allowbreak\lean{sound]}. That check is not part of this release.

%======================================================================
\section{Notes and references}\label{l11:sec:notes}
%======================================================================

\textbf{$C_2$, Zhu's algebra and the associated variety.} Y.~Zhu, ``Modular invariance of characters of vertex operator algebras'',
\emph{J. Amer.\ Math.\ Soc.} 9 (1996), 237--302 (the $C_2$-algebra, §4.4). H.~Li, ``Abelianizing vertex algebras'', \emph{Comm.\
Math.\ Phys.} 259 (2005), 391--411, arXiv:math/0409140 (the filtration $F_pV$, its multiplicativity, $F_1V = C_2(V)$, the vertex
Poisson structure on $\lelevengr^FV$, and the surjection from the arc algebra). T.~Arakawa, ``A remark on the $C_2$-cofiniteness
condition on vertex algebras'', \emph{Math.\ Z.} 270 (2012), 559--575, arXiv:1004.1492 (associated variety $X_V$, Li filtration,
lisse $\Leftrightarrow$ $\dim X_V = 0$). T.~Arakawa, ``Associated varieties of modules over Kac--Moody algebras and
$C_2$-cofiniteness of W-algebras'', \emph{IMRN} 2015, arXiv:1004.1554 (nilpotent cones at admissible level). T.~Arakawa and
A.~Moreau, \emph{Arc spaces and vertex algebras} (lecture notes) for jet schemes. J.~van~Ekeren and R.~Heluani, ``Chiral homology
of elliptic curves and the Zhu algebra'', \emph{Comm.\ Math.\ Phys.} 386 (2021), arXiv:1804.00017, and J.~Andrews, J.~van~Ekeren,
R.~Heluani, ``The singular support of the Ising model'' (2020) (classical freeness, and its failure).
H.~Karel and H.~Li, ``Certain generating subspaces for vertex operator algebras'', \emph{J. Algebra} 217 (1999). T.~Arakawa and A.~Moreau, ``Joseph ideals and lisse minimal W-algebras'', \emph{J. Inst.\ Math.\ Jussieu} 17 (2018),
397--417, arXiv:1506.00710 (the associated variety of $L_k(\mathfrak g)$ for the Deligne series at $k=-h^\vee/6-1$ is the minimal
nilpotent orbit closure; used in the warning of Section~\ref{l11:sec:genwords}).

\textbf{Spanning sets.} M.~R.~Gaberdiel and A.~Neitzke, ``Rationality, quasirationality and finite W-algebras'', \emph{Comm.\
Math.\ Phys.} 238 (2003), 305--331, arXiv:hep-th/0009235 (ordered PBW-type spanning sets from $C_2$). G.~Buhl, ``A spanning set
for VOA modules'', \emph{J. Algebra} 254 (2002), 125--151.

\textbf{Presentations and non-linear Lie conformal algebras.} A.~De~Sole and V.~G.~Kac, ``Freely generated vertex algebras and
non-linear Lie conformal algebras'', \emph{Comm.\ Math.\ Phys.} 254 (2005), 659--694, arXiv:math-ph/0312042 (Definition~3.1,
Lemma~3.2, Definition~3.4, Theorem~3.9; the formulas transcribed in \lean{NPExtension}). A.~De~Sole and V.~G.~Kac, ``Finite vs affine
W-algebras'', \emph{Japan.\ J. Math.} 1 (2006), 137--261, arXiv:math-ph/0511055 (the review cited as [DSK-rev] in the source).
V.~G.~Kac, \emph{Vertex Algebras for Beginners}, 2nd ed., AMS University Lecture Series 10 (1998), Thm.~4.5 (reconstruction).
E.~Frenkel and D.~Ben-Zvi, \emph{Vertex Algebras and Algebraic Curves}, 2nd ed., AMS (2004), Thm.~4.4.1. A.~R.~Linshaw,
``Universal two-parameter $\mathcal W_\infty$-algebra and vertex algebras of type $\mathcal W(2,3,\dots,N)$'', \emph{Compos.\ Math.} 157
(2021), arXiv:1710.02275 (§3.1: null fields of degenerate non-linear conformal algebras). M.~Roitman, ``Combinatorics of free vertex
algebras'', \emph{J. Algebra} 255 (2002), 297--323 (presentations of lattice vertex algebras). H.~Li, ``Local systems of twisted
vertex operators, vertex operator superalgebras and twisted modules'', \emph{Contemp.\ Math.} 193 (1996), 203--236 (the
$\Delta$-operator).

\textbf{Physics.} C.~Beem, M.~Lemos, P.~Liendo, W.~Peelaers, L.~Rastelli and B.~C.~van~Rees, ``Infinite chiral symmetry in four
dimensions'', \emph{Comm.\ Math.\ Phys.} 336 (2015), 1359--1433, arXiv:1312.5344. C.~Beem and L.~Rastelli, ``Vertex operator
algebras, Higgs branches, and modular differential equations'', \emph{JHEP} 08 (2018) 114, arXiv:1707.07679 (Higgs branch $=$
associated variety; the $R$-filtration). A.~Arabi~Ardehali, C.~Beem, M.~Lemos and L.~Rastelli, ``Graded unitarity in the SCFT/VOA
correspondence'', arXiv:2507.23781 (graded flags and good filtrations; the source of \lean{GoodFiltration}, not available in the
repository). F.~A.~Bais, P.~Bouwknegt, M.~Surridge and K.~Schoutens, ``Extensions of the Virasoro algebra constructed from
Kac--Moody algebras using higher order Casimir invariants'', \emph{Nucl.\ Phys.\ B} 304 (1988), 348--370 (the generalised Wick
theorem). P.~Bouwknegt and K.~Schoutens, ``W-symmetry in conformal field theory'', \emph{Phys.\ Rep.} 223 (1993), 183--276.

\sloppy
\textbf{Project records.} The staging of the presentation layer (Stages A1--A5 and the deferred Stage B) is documented in the module
docstrings of \leanpath{NonlinearConformal.lean}, \leanpath{NonlinearVirasoro.lean}, \leanpath{NonlinearZeta.lean},
\leanpath{NPExtension.lean}, \leanpath{NonlinearMR.lean}, \leanpath{VAPresentation.lean}, \leanpath{PresentationDoctrine.lean} and
of the two Virasoro example files. The Lean checks of this chapter made outside the library are not part of this release.

\endgroup
